% Generated mechanically from the frozen examples.tex target.
% Do not edit this generated file; edit the bound locale source instead.

% OK, start here.
%

\setcounter{chapter}{109}
\chapter{例}



\phantomsection
\label{examples-section-phantom}


\section{引言}
\label{examples-section-introduction}

\noindent
本章将收录若干阐明相关理论的例子。





\section{空的极限}
\label{examples-section-empty-limit}

\noindent
本例由 Waterhouse 给出，见 \cite{Waterhouse}。
设 $S$ 是不可数集。对每个有限子集
$T \subset S$，考虑单射 $T \to \mathbf{N}$ 所成的集合 $M_T$。
若 $T \subset T' \subset S$ 均有限，则限制映射 $M_{T'} \to M_T$
是满射。于是，以 $S$ 的有限子集所成的有向偏序集为指标，
我们得到一个转移映射均为满射的逆系统。
但 $\lim M_T = \emptyset$，因为极限中的一个元素将给出单射
$S \to \mathbf{N}$。



\section{零极限}
\label{examples-section-zero-limit}

\noindent
设 $(S_i)_{i \in I}$ 是由非空集合构成的有向逆系统，
其转移映射均为满射，并且
$\lim S_i = \emptyset$；见第 \ref{examples-section-empty-limit} 节。
设 $K$ 是域，并令
$$
V_i = \bigoplus\nolimits_{s \in S_i} K
$$
则转移映射 $V_i \to V_j$ 在 $i \geq j$ 时是满射。
然而，$\lim V_i = 0$。事实上，若 $v = (v_i)$ 是极限中的元素，
则 $v_i$ 的支撑将是有限子集 $T_i \subset S_i$，并且
$\lim T_i \not = \emptyset$；见
《范畴》中的引理 \ref{categories-lemma-nonempty-limit}。

\medskip\noindent
对每个 $i$，考虑唯一的 $K$-线性映射 $V_i \to K$，
它把每个基向量 $s \in S_i$ 送到 $1$。令 $W_i \subset V_i$
为其核。则
$$
0 \to (W_i) \to (V_i) \to (K) \to 0
$$
是以有向集 $I$ 为指标的向量空间逆系统的不分裂短正合列。因此
% P12-ERR-0315
$W_i$ 是转移映射均为满射、极限为零而
$R^1\lim$ 非零的 $K$-向量空间有向系统。



\section{拟紧空间的非拟紧逆极限}
\label{examples-section-lim-not-quasi-compact}

\noindent
以 $\mathbf{N}$ 表示自然数集。
对每个整数 $n$，令 $I_n$ 表示所有自然数 $> n$ 所成的集合。
把 $T_n$ 定义为 $\mathbf{N}$ 上以
$\{1\}, \ldots , \{n\}, I_n$ 为基的唯一拓扑。以 $X_n$ 表示拓扑空间
$(\mathbf{N}, T_n)$。对每个 $m < n$，恒等映射
$$
f_{n, m} : X_n \longrightarrow X_m
$$
连续。显然，当 $m < n < p$ 时，复合
$f_{p, n} \circ f_{n, m}$ 等于 $f_{p, m}$。因此 $((X_n), (f_{n,m}))$
是由拟紧拓扑空间构成的有向逆系统。

\medskip\noindent
设 $T$ 是 $\mathbf{N}$ 上的离散拓扑，并令 $X$ 为
$(\mathbf{N}, T)$。则对每个整数 $n$，恒等映射
$$
f_n : X \longrightarrow X_n
$$
连续。我们断言，这是上述有向系统的逆极限。设 $(Y, S)$ 是任意拓扑空间。
对每个整数 $n$，设
$$
g_n : (Y, S) \longrightarrow (\mathbf{N}, T_n)
$$
为连续映射。假设对每个 $m < n$ 都有
$f_{n,m} \circ g_n = g_m$，即系统 $(g_n)$ 与上述有向系统相容。
特别地，所有集合映射 $g_n$ 都等于同一个集合映射
$$
g : Y \longrightarrow \mathbf{N}.
$$
此外，对每个整数 $n$，由于 $\{n\}$ 在 $X_n$ 中开，
$g^{-1}(\{n\}) = g_n^{-1}(\{n\})$ 在 $Y$ 中也开。
因此，集合映射 $g$ 对拓扑 $S$（定义在 $Y$ 上）与拓扑 $T$
（定义在 $\mathbf{N}$ 上）连续。故 $(X, (f_n))$ 是上述有向系统的逆极限。

\medskip\noindent
然而，$X$ 显然不是拟紧的，因为由单点集组成的无限开覆盖没有逆极限。
% P12-ERR-0316

\begin{lemma}
\label{examples-lemma-lim-not-quasi-compact}
存在一个以 $\mathbf{N}$ 为指标的拟紧拓扑空间逆系统，其极限并不拟紧。
\end{lemma}

\begin{proof}
见上面的讨论。
\end{proof}





\section{纤维积上的结构层}
\label{examples-section-silly}

\noindent
设 $X, Y, S, a, b, p, q, f$ 如《概形的导出范畴》引言
第 \ref{perfect-section-kunneth} 节中所述。图示为
$$
\xymatrix{
& X \times_S Y \ar[ld]^p \ar[rd]_q \ar[dd]^f \\
X \ar[rd]_a & & Y \ar[ld]^b \\
& S
}
$$
于是有典范映射
$$
can :
p^{-1}\mathcal{O}_X \otimes_{f^{-1}\mathcal{O}_S} q^{-1}\mathcal{O}_Y
\longrightarrow
\mathcal{O}_{X \times_S Y}
$$
它一般不是同构。

\medskip\noindent
例如，令 $S = \Spec(\mathbf{R})$、$X = \Spec(\mathbf{C})$，并令
$Y = \Spec(\mathbf{C})$。则
$X \times_S Y = \Spec(\mathbf{C}) \amalg \Spec(\mathbf{C})$
是由两个点构成的离散空间，而层 $p^{-1}\mathcal{O}_X$、
$q^{-1}\mathcal{O}_Y$ 与 $f^{-1}\mathcal{O}_S$ 分别是取值于
$\mathbf{C}$、$\mathbf{C}$ 与 $\mathbf{R}$ 的常值层。
因此，$can$ 的源是这个两点离散空间上取值于
$\mathbf{C} \otimes_\mathbf{R} \mathbf{C}$ 的常值层。
故其整体截面作为 $\mathbf{R}$-向量空间的维数为 $8$；
另一方面，对 $can$ 的目标取整体截面得到
$\mathbf{C} \times \mathbf{C}$，
它作为 $\mathbf{R}$-向量空间的维数为 $4$。

\medskip\noindent
另一个例子如下。设 $k$ 是代数闭域。考虑
$S = \Spec(k)$、$X = \mathbf{A}^1_k$ 与 $Y = \mathbf{A}^1_k$。
若 $U \subset X \times_S Y = \mathbf{A}^2_k$ 是非空开集，
则像 $p(U) \subset X = \mathbf{A}^1_k$ 与 $q(U) \subset \mathbf{A}^1_k$
均为开集；读者可以证明
$$
\left(
p^{-1}\mathcal{O}_X \otimes_{f^{-1}\mathcal{O}_S} q^{-1}\mathcal{O}_Y
\right)(U) = \mathcal{O}_X(p(U)) \otimes_k \mathcal{O}_Y(q(U))
$$
若 $U$ 是 $X \times_S Y = \mathbf{A}^2_k$ 中一条不可约曲线 $C$
的补集，并且 $p|_C$ 与 $q|_C$ 都非常值，
则这不等于 $\mathcal{O}_{X \times_S Y}(U)$。

\medskip\noindent
回到一般情形，设 $z = (x, y, s, \mathfrak p)$ 是
$X \times_S Y$ 的一点，如《概形》中的引理
\ref{schemes-lemma-points-fibre-product} 所述。
则映射 $can$ 在 $z$ 处的茎上给出映射
$$
can_z :
\mathcal{O}_{X, x} \otimes_{\mathcal{O}_{S, s}} \mathcal{O}_{Y, y}
\longrightarrow
\mathcal{O}_{X \times_S Y, z}
$$
这是平坦环同态，因为目标是源的一个局部化
（细节从略；提示：约化到 $X$、$Y$ 与 $S$ 均为仿射的情形）。
注意，源一般不是局部环；这也从另一个角度说明 $can$ 一般不是同构。

\medskip\noindent
更一般地，设给定一个 $\mathcal{O}_X$-模 $\mathcal{F}$ 与一个
$\mathcal{O}_Y$-模 $\mathcal{G}$。则有典范映射
% P12-ERR-0317
\begin{align*}
& p^{-1}\mathcal{F} \otimes_{f^{-1}\mathcal{O}_S} q^{-1}\mathcal{G} \\
& =
p^{-1}(\mathcal{F} \otimes_{\mathcal{O}_X} \mathcal{O}_X)
\otimes_{f^{-1}\mathcal{O}_S}
q^{-1}(\mathcal{O}_Y \otimes_{\mathcal{O}_Y} \mathcal{G}) \\
& =
p^{-1}\mathcal{F} \otimes_{p^{-1}\mathcal{O}_X} p^{-1}\mathcal{O}_X
\otimes_{f^{-1}\mathcal{O}_S}
q^{-1}\mathcal{O}_Y \otimes_{q^{-1}\mathcal{O}_Y} q^{-1}\mathcal{G} \\
& \xrightarrow{can}
p^{-1}\mathcal{F} \otimes_{q^{-1}\mathcal{O}_X}
\mathcal{O}_{X \times_S Y}
\otimes_{q^{-1}\mathcal{O}_Y} q^{-1}\mathcal{G}  \\
& =
p^{-1}\mathcal{F} \otimes_{q^{-1}\mathcal{O}_X}
\mathcal{O}_{X \times_S Y}
\otimes_{\mathcal{O}_{X \times_S Y}}
\mathcal{O}_{X \times_S Y}
\otimes_{q^{-1}\mathcal{O}_Y} q^{-1}\mathcal{G}  \\
& =
p^*\mathcal{F} \otimes_{\mathcal{O}_{X \times_S Y}} q^*\mathcal{G}
\end{align*}
它很少是同构。





\section{一个非整的连通概形，其局部环均为整环}
\label{examples-section-connected-locally-integral-not-integral}

\noindent
我们给出仿射概形 $X = \Spec(A)$ 的一个例子：它是连通的，
其所有局部环都是整环，但它不是整概形。
$X$ 的连通性意味着 $A$ 没有非平凡幂等元；见
《代数》中的引理 \ref{algebra-lemma-disjoint-decomposition}。
若每当 $A$ 中 $fg = 0$ 时，$X$ 的每一点都有一个邻域，
使得 $f$ 或 $g$ 之一在其上为零，则 $X$ 的局部环均为整环。
只要 $A$ 不是整环，$X$ 就不是整概形
（《性质》中的定义 \ref{properties-definition-integral}）。

\medskip\noindent
粗略地说，构造如下：令 $X_0$ 为仿射平面中的十字形
（即两条坐标轴的并）。然后，在平面的吹起上取 $X_0$ 的
赋予约化结构的全逆像 $X_1$（$X_1$ 由三条排成链的有理分支组成）。
再在所得曲面上沿 $X_1$ 的两个奇点吹起，并令 $X_2$ 为
$X_1$ 的约化逆像（它有五条有理分支），如此继续。
取 $X$ 为逆极限。这个构造唯一的问题是，吹起会粘入一条射影直线，
所以 $X_1$ 不是仿射的。我们改为粘入一条仿射直线以修正这一点
（因而我们的概形将是上述概形的一个开子集）。

\medskip\noindent
下面给出完全代数的构造：对每个 $k \ge 0$，令 $A_k$
为如下环：其元素是多项式族
$p_i \in \mathbf{C}[x]$，其中 $i = 0, \ldots, 2^k$，并且满足
$p_i(1) = p_{i + 1}(0)$。令 $X_k = \Spec(A_k)$。注意，$X_k$
是 $2^k + 1$ 条仿射直线的并，它们横截相交并排成一条链。
如下定义环同态 $A_k \to A_{k + 1}$：
$$
(p_0, \ldots, p_{2^k})
\longmapsto
(p_0, p_0(1), p_1, p_1(1), \ldots, p_{2^k}),
$$
换言之，每隔一个多项式便是常值的。这个映射把
$A_k$ 认作 $A_{k + 1}$ 的子环。令 $A$ 为 $A_k$ 的直极限
（本质上即它们的并）。令 $X = \Spec(A)$。对每个 $k$，都有
自然嵌入 $A_k \to A$，也就是映射 $X\to X_k$。
每个 $A_k$ 都连通但不整；这蕴含 $A$ 连通但不整。
尚需证明 $A$ 的局部环均为整环。

\medskip\noindent
取 $f, g \in A$，满足 $fg = 0$，并取 $x \in X$。我们来构造
$x$ 的一个邻域，使 $f$ 与 $g$ 之一在其上为零。选取 $k$，
使得 $f, g \in A_{k - 1}$（注意指标为 $k - 1$）。
令 $y$ 为 $x$ 在 $X_k$ 中的像。只需证明 $y$ 有一个邻域，
使得视为 $\mathcal{O}_{X_k}$ 截面的 $f$ 或 $g$ 之一在其上为零。
若 $y$ 是 $X_k$ 的光滑点，即它只位于 $2^k + 1$ 条直线中的一条上，
则这很显然。因此，可以假设 $y$ 是 $2^k$ 个奇点之一，
从而 $X_k$ 的两条分支经过 $y$。然而，在这两条分支中的一条
（指标为奇数的那条）上，$f$ 与 $g$ 都是常值的，因为它们是
$X_{k - 1}$ 上函数的拉回。由于处处都有 $fg = 0$，
$f$ 或 $g$ 之一（不妨设为 $f$）在另一条分支上为零。
这蕴含 $f$ 在两条分支上都为零，正合所需。

\section{非完备的完备化}
\label{examples-section-noncomplete-completion}

\noindent
设 $R$ 是环，$\mathfrak m$ 是极大理想。考虑完备化
$$
R^\wedge = \lim R/\mathfrak m^n.
$$
注意，$R^\wedge$ 是局部环，其极大理想为
$\mathfrak m' = \Ker(R^\wedge \to R/\mathfrak m)$。
事实上，若 $x = (x_n) \in R^\wedge$ 不属于 $\mathfrak m'$，则
$y = (x_n^{-1}) \in R^\wedge$ 满足 $xy = 1$，于是由
《代数》中的引理 \ref{algebra-lemma-characterize-local-ring}，
$R^\wedge$ 是局部环。现在，$R^\wedge$ 关于其极限拓扑总是完备的
% P12-ERR-0318
（见《代数进阶》第 \ref{more-algebra-section-topological-ring} 节中的讨论）。
但除此之外，还有如下问题：
\begin{enumerate}
\item 是否有 $\mathfrak m R^\wedge = \mathfrak m'$？
\item 把 $R^\wedge$ 视为 $R^\wedge$-模时，它是否 $\mathfrak m'$-进完备？
\item 把 $R^\wedge$ 视为 $R$-模时，它是否 $\mathfrak m$-进完备？
\end{enumerate}
事实证明，这些问题的答案都是否定的。
下面的例子取自 Bart de Smit 与 Hendrik Lenstra 的一份未发表笔记。
另见 \cite[Exercise III.2.12]{Bourbaki-CA} 与
\cite[Example 1.8]{Yekutieli}

\medskip\noindent
设 $k$ 是域，$R = k[x_1, x_2, x_3, \ldots]$，且
$\mathfrak m = (x_1, x_2, x_3, \ldots)$。
我们把 $R^\wedge$ 的元素 $f$ 看作一个（可能）无限和
$$
f = \sum a_I x^I
$$
（采用多重指标记号），并要求对每个 $d \geq 0$，
满足 $|I| = d$ 的非零 $a_I$ 只有有限多个。极大理想
$\mathfrak m' \subset R^\wedge$ 由常数项为零的 $f$ 组成。
特别地，元素
$$
f = x_1 + x_2^2 + x_3^3 + \ldots
$$
属于 $\mathfrak m'$，但不属于 $\mathfrak m R^\wedge$，
这说明本例中的 (1) 不成立。然而，若 (1) 不成立，
则 (3) 必然不成立，因为
$\mathfrak m' = \Ker(R^\wedge \to R/\mathfrak m)$，
并且可以对 $n = 1$ 应用
《代数》中的引理 \ref{algebra-lemma-hathat}。

\medskip\noindent
最后，我们证明 $R^\wedge$ 不是 $\mathfrak m'$-进完备的。
对 $n \geq 1$，令 $K_n = \Ker(R^\wedge \to R/\mathfrak m^n)$。于是有短正合列
$$
0 \to K_n/(\mathfrak m')^n \to R^\wedge/(\mathfrak m')^n \to
R/\mathfrak m^n \to 0
$$
投影映射 $R^\wedge \to R/\mathfrak m^{n + 1}$ 把
$(\mathfrak m')^n$ 满射到 $\mathfrak m^n/\mathfrak m^{n + 1}$。
由此，$K_{n + 1} \to K_n/(\mathfrak m')^n$ 是满射。
因此，逆系统 $\left(K_n/(\mathfrak m')^n\right)$
的转移映射均为满射；取逆极限，并应用
《代数》中的引理 \ref{algebra-lemma-Mittag-Leffler}，得到正合列
$$
0 \to \lim K_n/(\mathfrak m')^n \to
\lim R^\wedge/(\mathfrak m')^n \to
\lim R/\mathfrak m^n \to 0
$$
于是可见，$R^\wedge$ 关于 $\mathfrak m'$ 完备，当且仅当
对所有 $n \geq 1$ 都有 $K_n = (\mathfrak m')^n$。

\medskip\noindent
为说明本例中的 $R^\wedge$ 不是 $\mathfrak m'$-进完备的，
我们证明 $K_2 = \Ker(R^\wedge \to R/\mathfrak m^2)$
不等于 $(\mathfrak m')^2$。
注意，$(\mathfrak m')^2$ 的元素可以写成有限和
\begin{equation}
\label{examples-equation-sum}
\sum\nolimits_{i = 1, \ldots, t} f_i g_i
\end{equation}
其中 $f_i, g_i \in R^\wedge$ 的常数项为零。
为构造一个例子，我们将选取如下形式的 $z \in K_2$
% P12-ERR-0319
$$
z = z_1 + z_2 + z_3 + \ldots
$$
并使它满足下列性质：
\begin{enumerate}
\item 存在序列 $1 < d_1 < d_2 < d_3 < \ldots $ 与
$0 < n_1 < n_2 < n_3 < \ldots$。

其中
$z_i \in k[x_{n_i}, x_{n_i + 1}, \ldots, x_{n_{i + 1} - 1}]$
是次数为 $d_i$ 的齐次元；并且
\item 在环 $k[[x_{n_i}, x_{n_i + 1}, \ldots, x_{n_{i + 1} - 1}]]$
中，元素 $z_i$ 不能写成 $t \leq i$ 的和式 (\ref{examples-equation-sum})。
\end{enumerate}
这显然蕴含 $z$ 不属于 $(\mathfrak m')^2$，因为当 $i$ 足够大时，
关系式 (\ref{examples-equation-sum}) 在环
$k[[x_{n_i}, x_{n_i + 1}, \ldots, x_{n_{i + 1} - 1}]]$
中的像会导出矛盾。因此，只需证明：对每个 $t > 0$，
都存在 $d \gg 0$ 与整数 $n$，使我们能够找到
次数为 $d$ 的齐次元 $z \in k[x_1, \ldots, x_n]$，
它在 $k[[x_1, \ldots, x_n]]$ 中不能针对给定的 $t$
写成和式 (\ref{examples-equation-sum})。取 $n > 2t$，再取任意
与 $k$ 的特征互素的 $d > 1$，并令
$z = \sum_{i = 1, \ldots, n} x_i^d$。则理想
$$
(\frac{\partial z}{\partial x_1}, \ldots, \frac{\partial z}{\partial x_n})
=
(dx_1^{d - 1}, \ldots, dx_n^{d - 1})
$$
的零点集只含一个点。另一方面，由 Leibniz 法则，
$$
\frac{\partial ( \sum\nolimits_{i = 1, \ldots, t} f_i g_i ) }{\partial x_j}
\in (f_1, \ldots, f_t, g_1, \ldots, g_t)
$$
故这些偏导数的零点集至少包含
$$
V(f_1, \ldots, f_t, g_1, \ldots, g_t) \subset
\Spec(k[[x_1, \ldots, x_n]]).
$$
这便产生矛盾，因为
$V(f_1, \ldots, f_t, g_1, \ldots, g_t)$ 的维数至少为
$n - 2t \geq 1$。

\begin{lemma}
\label{examples-lemma-noncomplete-completion}
存在局部环 $R$ 与极大理想 $\mathfrak m$，使得 $R$ 关于
$\mathfrak m$ 的完备化 $R^\wedge$ 具有下列性质：
\begin{enumerate}
\item $R^\wedge$ 是局部环，但其极大理想不等于
$\mathfrak m R^\wedge$；
\item $R^\wedge$ 不是完备局部环；并且
\item 作为 $R$-模，$R^\wedge$ 不是 $\mathfrak m$-进完备的。
\end{enumerate}
\end{lemma}

\begin{proof}
由上面的讨论即得，因为（取 $R = k[x_1, x_2, x_3, \ldots]$）
局部化 $R_{\mathfrak m}$ 的完备化等于 $R$ 的完备化。
\end{proof}


\section{非完备商}
\label{examples-section-noncomplete-quotient}

\noindent
设 $k$ 是域。令
$$
R = k[t, z_1, z_2, z_3, \ldots, w_1, w_2, w_3, \ldots, x]/
(z_it - x^iw_i, z_i w_j)
$$
特别注意，在此环中 $z_iz_jt = 0$。$R$ 的任意元素 $f$
都能唯一地写成有限和
$$
f = \sum\nolimits_{i = 0, \ldots, d} f_i x^i
$$
其中每个 $f_i \in k[t, z_i, w_j]$ 都不含涉及乘积
$z_it$ 或 $z_iw_j$ 的项。此外，若 $f$ 以此方式写出，
则 $f \in (x^n)$ 当且仅当对 $i < n$ 有 $f_i = 0$。
所以 $x$ 是非零因子，且 $\bigcap (x^n) = 0$。
令 $R^\wedge$ 为 $R$ 关于理想 $(x)$ 的完备化。
注意，$R^\wedge$ 是 $(x)$-进完备的；见
《代数》中的引理 \ref{algebra-lemma-hathat-finitely-generated}。
由上可见，$R^\wedge$ 的元素能唯一地写成无限和
$$
f = \sum\nolimits_{i = 0}^\infty f_i x^i
$$
其中每个 $f_i \in k[t, z_i, w_j]$ 都不含涉及乘积
$z_it$ 或 $z_iw_j$ 的项。考虑元素
$$
f = \sum\nolimits_{i = 1}^\infty x^i w_i =
xw_1 + x^2w_2 + x^3w_3 + \ldots
$$
也就是说，$f_n = w_n$。注意，对每个 $n$ 都有
$f \in (t , x^n)$，因为对所有 $m$ 都有 $x^mw_m \in (t)$。
我们断言 $f \not \in (t)$。为证明这一点，假设
$tg = f$，其中 $g = \sum g_lx^l$ 采用上述典范形式。
由于 $tz_iz_j = 0$，不妨假设没有任何 $g_l$ 含涉及乘积
$z_iz_j$ 的项。考察把 $tg$ 化为典范形式的过程，可得如下规则：
给定 $g_l$ 的任意一项 $c m$，其中 $c \in k$，而 $m$ 是
$t, z_i, w_j$ 中的单项式，作如下替换：
\begin{enumerate}
\item 若单项式 $m$ 不含任何 $z_i$，则 $ctm$ 是 $f_l$ 的一项；
\item 若单项式 $m$ 含有某个 $z_i$，则它等于
$m = z_i$，并且 $cw_i$ 是 $f_{l + i}$ 中的一项。
% P12-ERR-0320
\end{enumerate}
由于 $g_0$ 是多项式，其中只出现有限多个变量 $z_i$。
选取 $n$，使 $z_n$ 不出现在 $g_0$ 中。
于是上述规则表明 $w_n$ 不出现在 $f_n$ 中，矛盾。
因此，$R^\wedge/(t)$ 并不完备；见
《代数》中的引理 \ref{algebra-lemma-quotient-complete}。

\begin{lemma}
\label{examples-lemma-noncomplete-quotient}
存在一个关于主理想 $I$ 完备的环 $R$ 以及一个主理想 $J$，
使得 $R/J$ 不是 $I$-进完备的。
\end{lemma}

\begin{proof}
见上面的讨论。
\end{proof}


\section{完备化不正合}
\label{examples-section-completion-not-exact}

\noindent
一个简明的例子如下。设 $R = k[t]$。以 $M^\wedge$
表示 $R$-模关于 $I = (t)$ 的完备化。令
$P = K = \bigoplus_{n \in \mathbf{N}} R$，并令
$M = \bigoplus_{n \in \mathbf{N}} R/(t^n)$。则有短正合列
$0 \to K \to P \to M \to 0$，其中第一个映射在第 $n$ 个直和项上
由乘以 $t^n$ 给出。我们断言
$0 \to K^\wedge \to P^\wedge \to M^\wedge \to 0$ 在中间不正合。
事实上，$\xi = (t^2, t^3, t^4, \ldots) \in P^\wedge$
在 $M^\wedge$ 中的像为零，但它不在 $K^\wedge \to P^\wedge$ 的像中，
因为若在像中，它将是 $(t, t, t, \ldots)$ 的像，而后者并非
$K^\wedge$ 的元素。

\medskip\noindent
另一个“更小”的例子如下。在引理
\ref{examples-lemma-noncomplete-quotient} 的情形下，短正合列
$0 \to J \to R \to R/J \to 0$ 在关于 $I$ 完备化后不再正合。
事实上，若 $f \in J$ 是生成元，则
$f : R \to J$ 是满射，故 $R \to J^\wedge$ 是满射，因而
$J^\wedge \to R$ 的像为 $(f) = J$；但 $R/J$ 不完备意味着满射
$R \to (R/J)^\wedge$ 的核严格大于 $J$，见
《代数》中的引理 \ref{algebra-lemma-completion-generalities} 与
\ref{algebra-lemma-quotient-complete}。
同理，序列
$R \to R \to R/(f) \to 0$ 在完备化后也不再正合。

\begin{lemma}
\label{examples-lemma-completion-not-exact}
\begin{slogan}
一般而言，完备化既非左正合，也非右正合。
\end{slogan}
一般而言，完备化不是正合函子；甚至一般也不是右正合函子。
即便 $I$ 有限生成，并把它限制在有限表示模的范畴上，这仍然成立。
\end{lemma}

\begin{proof}
见上面的讨论。
\end{proof}



\section{完备模范畴不是 Abel 范畴}
\label{examples-section-non-abelian}

\noindent
设 $R$ 是环，$I \subset R$ 是有限生成理想。考虑由 $I$-进完备
$R$-模组成的范畴 $\mathcal{A}$；见
《代数》中的定义 \ref{algebra-definition-complete}。
设 $\varphi : M \to N$ 是 $\mathcal{A}$ 中的态射。
$\varphi$ 在 $\mathcal{A}$ 中的余核是通常余核的 $I$-进完备化
$(\Coker(\varphi))^\wedge$
（由于 $I$ 有限生成，这个完备化是完备的；见
《代数》中的引理 \ref{algebra-lemma-hathat-finitely-generated}）。
令 $K = \Ker(\varphi)$；由于 $\varphi$ 连续，它是 $M$ 的闭子模。
由下面的引理 \ref{examples-lemma-closed-is-complete}，$K$ 是 $I$-进完备的，
因而就是 $\varphi$ 在 $\mathcal{A}$ 中的核。

\begin{lemma}
\label{examples-lemma-closed-is-complete}
设 $I \subset R$ 是环的理想，$M$ 是 $I$-进完备的 $R$-模，
$N \subset M$ 是子模。下列条件等价：
\begin{enumerate}
\item $N$ 在 $M$ 中闭；
\item $N = \bigcap_{n \geq 1} (N + I^n M)$；
\item $M/N$ 是 $I$-进完备的。
\end{enumerate}
若 $I$ 有限生成，则这些条件蕴含 $N$ 是 $I$-进完备的。
\end{lemma}

\begin{proof}
(1) 与 (2) 的等价性来自
《形式空间》中的引理 \ref{formal-spaces-lemma-closed}。
(2) 与 (3) 的等价性是
《代数》中的引理 \ref{algebra-lemma-quotient-complete}。
假设 $I$ 有限生成且 (1) 成立。令 $N^\wedge$ 为 $N$ 的 $I$-进完备化。
由于 $M$ 是 $I$-进完备的，包含映射具有分解
$$
N \to N^\wedge \to M
$$
由于 $N^\wedge \to M$ 连续，$N$ 在 $N^\wedge$ 中稠密，
而 $N$ 在 $M$ 中闭，可见 $N^\wedge \to M$ 经 $N$ 分解。
于是，对某个 $A$-模 $C$，有 $N^\wedge = N \oplus C$。
% P12-ERR-0321
因此，$N$ 是 $I$-进完备模 $N^\wedge$ 的直和项；见
《代数》中的引理 \ref{algebra-lemma-hathat-finitely-generated}
（此处用到了 $I$ 有限生成），故 $N$ 是 $I$-进完备的。
\end{proof}

\noindent
我们给出一个例子，说明一般在范畴 $\mathcal{A}$ 中
$\Im \not = \Coim$。取
$R = \mathbf{Z}_p = \lim_n \mathbf{Z}/p^n\mathbf{Z}$
为 $p$-进整数环，并取 $I = (p)$。考虑映射
$$
\text{diag}(1, p, p^2, \ldots) :
\left(\bigoplus\nolimits_{n \geq 1} \mathbf{Z}_p\right)^\wedge
\longrightarrow
\prod\nolimits_{n \geq 1} \mathbf{Z}_p
$$
其中左侧是直和的 $p$-进完备化。因此，左侧的元素是向量
$(x_1, x_2, x_3, \ldots)$，其中 $x_i \in \mathbf{Z}_p$，且当
$i \to \infty$ 时其 $p$-进赋值 $v_p(x_i) \to \infty$。
它被送到 $(x_1, px_2, p^2x_3, \ldots)$。由此可见，
$(1, p, p^2, \ldots)$ 属于像的闭包，但不属于像。
由上面对核与余核的描述，显然对该映射有
$\Im \not = \Coim$。

\begin{lemma}
\label{examples-lemma-complete-modules-not-abelian}
设 $R$ 是环，$I \subset R$ 是有限生成理想。
$I$-进完备 $R$-模的范畴具有核与余核，但一般不是 Abel 范畴，
即使 $R$ 是 Noether 环亦然。
\end{lemma}

\begin{proof}
见上文。
\end{proof}



\section{导出完备模范畴}
\label{examples-section-derived-complete-modules}

\noindent
阅读本节之前，请先阅读《代数进阶》第
\ref{more-algebra-section-derived-complete-modules} 节。

\medskip\noindent
设 $A$ 是环，$I$ 是 $A$ 的理想，并以 $\mathcal{C}$ 表示
《代数进阶》定义 \ref{more-algebra-definition-derived-complete}
所定义的导出完备模范畴。

\medskip\noindent
设 $T$ 是集合，$M_t$（$t \in T$）是一族导出完备模。
我们断言，一般而言，$\bigoplus M_t$ 不是导出完备模。
作为一个具体例子，令 $A = \mathbf{Z}_p$、$I = (p)$，并考虑
$\bigoplus_{n \in \mathbf{N}} \mathbf{Z}_p$。从
$\bigoplus_{n \in \mathbf{N}} \mathbf{Z}_p$ 到其 $p$-进完备化的映射
不是满射。这意味着 $\bigoplus_{n \in \mathbf{N}} \mathbf{Z}_p$
不可能导出完备，否则将蕴含该映射为满射；见
《代数进阶》中的引理 \ref{more-algebra-lemma-complete-derived-complete}。
因此，包含函子 $\mathcal{C} \to \text{Mod}_A$
既不与直和交换，也不与（滤过）余极限交换。

\medskip\noindent
假设 $I$ 有限生成。由《代数进阶》第
\ref{more-algebra-section-derived-complete-modules} 节中的讨论，
范畴 $\mathcal{C}$ 具有任意余极限。然而，我们断言，
滤过余极限在范畴 $\mathcal{C}$ 中并不正合。事实上，设
$A = \mathbf{Z}_p$ 且 $I = (p)$。存在 $p$-进完备 $A$-模的嵌入
$f_n : \mathbf{Z}_p/p\mathbf{Z}_p \to \mathbf{Z}_p/p^n\mathbf{Z}_p$，
由乘以 $p^{n - 1}$ 给出。存在交换图
$$
\xymatrix{
\mathbf{Z}_p/p\mathbf{Z}_p \ar[r]_{f_n} \ar[d]^1 &
\mathbf{Z}_p/p^n\mathbf{Z}_p \ar[d]_p \\
\mathbf{Z}_p/p\mathbf{Z}_p \ar[r]^{f_{n + 1}} &
\mathbf{Z}_p/p^{n + 1}\mathbf{Z}_p
}
$$
我们断言：这些嵌入在范畴 $\mathcal{C}$ 中的余极限给出映射
$\mathbf{Z}_p/p\mathbf{Z}_p \to 0$。事实上，右侧系统在
$\text{Mod}_A$ 中的余极限是 $\mathbf{Q}_p/\mathbf{Z}_p$。
因此，其在 $\mathcal{C}$ 中的余极限为
$$
H^0((\mathbf{Q}_p/\mathbf{Z}_p)^\wedge) =
H^0(\mathbf{Z}_p[1]) = 0
$$
这是《代数进阶》第
\ref{more-algebra-section-derived-complete-modules} 节的结论，
其中 ${}^\wedge$ 表示导出完备化。这证明了我们的断言。

\begin{lemma}
\label{examples-lemma-derived-complete-modules}
设 $A$ 是环，$I \subset A$ 是理想。
导出完备模的范畴 $\mathcal{C}$ 是 Abel 范畴，而包含函子
$F : \mathcal{C} \to \text{Mod}_A$ 正合且与任意极限交换。
若 $I$ 有限生成，则 $\mathcal{C}$ 具有任意直和与余极限，
但一般而言，$F$ 不与它们交换。最后，滤过余极限一般在
$\mathcal{C}$ 中并不正合，故 $\mathcal{C}$ 不是 Grothendieck Abel 范畴。
\end{lemma}

\begin{proof}
见《代数进阶》中的引理
\ref{more-algebra-lemma-derived-complete-modules} 以及上面的讨论。
\end{proof}



\section{非平坦完备化}
\label{examples-section-nonflat}

\noindent
与 Noether 情形不同，环关于理想的完备化并不总是平坦的。
我们已在《代数进阶》的例
\ref{more-algebra-example-not-glueing-pair} 中见过这种现象的两个例子。
本节再给出两个例子。

\begin{lemma}
\label{examples-lemma-countable-fg-tensor}
设 $R$ 是环，$M$ 是可数的 $R$-模。则 $M$ 是有限 $R$-模，
当且仅当 $M \otimes_R R^\mathbf{N} \to M^\mathbf{N}$ 是满射。
\end{lemma}

\begin{proof}
若 $M$ 是有限模，则由《代数》中的命题
\ref{algebra-proposition-fg-tensor}，该映射是满射。反之，假设该映射满射。
把 $M$ 的元素枚举为 $m_1, m_2, m_3, \ldots$。
设张量积中的元素 $\sum_{j = 1, \ldots, m} x_j \otimes a_j$
被送到 $(m_n) \in M^\mathbf{N}$。于是，正如《代数》命题
\ref{algebra-proposition-fg-tensor} 的证明中那样，可见
$x_1, \ldots, x_m$ 在 $R$ 上生成 $M$。
\end{proof}

\begin{lemma}
\label{examples-lemma-countable-fp-tensor}
设 $R$ 是可数环，$M$ 是可数的 $R$-模。则 $M$ 有限表示，
当且仅当典范映射
$M \otimes_R R^\mathbf{N} \to M^\mathbf{N}$ 是同构。
\end{lemma}

\begin{proof}
若 $M$ 是有限表示模，则由《代数》中的命题
\ref{algebra-proposition-fp-tensor}，该映射是同构。反之，假设该映射是同构。
由引理 \ref{examples-lemma-countable-fg-tensor}，模 $M$ 是有限的。
选取以 $K$ 为核的满射 $R^{\oplus m} \to M$。则 $K$ 作为
$R^{\oplus m}$ 的子模是可数的。仿照《代数》命题
\ref{algebra-proposition-fp-tensor} 的证明，可见
$K \otimes_R R^\mathbf{N} \to K^\mathbf{N}$ 是满射。
于是由引理 \ref{examples-lemma-countable-fg-tensor}，$K$ 是有限 $R$-模。
故 $M$ 有限表示。
\end{proof}

\begin{lemma}
\label{examples-lemma-countable-coherent}
设 $R$ 是可数环。则 $R$ 是凝聚环，当且仅当
$R^\mathbf{N}$ 是平坦 $R$-模。
\end{lemma}

\begin{proof}
若 $R$ 是凝聚环，则由《代数》中的命题
\ref{algebra-proposition-characterize-coherent}，$R^\mathbf{N}$ 是平坦模。
假设 $R^\mathbf{N}$ 平坦。设 $I \subset R$ 是有限生成理想。
为证明该引理，我们证明 $I$ 作为 $R$-模有限表示。事实上，
由于 $R^\mathbf{N}$ 平坦，映射
$I \otimes_R R^\mathbf{N} \to R^\mathbf{N}$ 是单射；
由引理 \ref{examples-lemma-countable-fg-tensor}，其像为 $I^\mathbf{N}$。
于是由引理 \ref{examples-lemma-countable-fp-tensor} 得出结论。
\end{proof}

\noindent
设 $R$ 是可数环。注意，作为 $R$-模，$R[[x]]$ 同构于
$R^\mathbf{N}$。由引理 \ref{examples-lemma-countable-coherent} 可见，
$R \to R[[x]]$ 平坦，当且仅当 $R$ 是凝聚环。
存在许多非凝聚的可数环，例如
$$
R = k[y, z, a_1, b_1, a_2, b_2, a_3, b_3, \ldots]/
(a_1 y + b_1 z, a_2 y + b_2 z, a_3 y + b_3 z, \ldots)
$$
其中 $k$ 是可数域。该环不是凝聚环，因为 $R$ 的理想
$(y, z)$ 不是有限表示 $R$-模。注意，$R[[x]]$ 是
$R[x]$ 关于主理想 $(x)$ 的完备化。

\begin{lemma}
\label{examples-lemma-completion-polynomial-ring-not-flat}
存在一个环，使得 $R[x]$ 在 $(x)$ 处的完备化 $R[[x]]$
在 $R$ 上不平坦，因而更不在 $R[x]$ 上平坦。
\end{lemma}

\begin{proof}
见上面的讨论。
\end{proof}

\noindent
事实证明，存在环 $R$，使得 $R[[x]]$ 在 $R$ 上平坦，
但 $R[[x]]$ 在 $R[x]$ 上不平坦。见 Badam Baplan 的
\href{https://math.stackexchange.com/users/164860/badam-baplan}{这篇帖子}。
具体而言，令 $R$ 为赋值环。则 $R$ 是凝聚环
（《代数》中的例 \ref{algebra-example-valuation-ring-coherent}），
因而由《代数》中的命题
\ref{algebra-proposition-characterize-coherent}，$R[[x]]$ 在 $R$ 上平坦。
另一方面，有如下引理。

\begin{lemma}
\label{examples-lemma-almost-integral-when-powerseries-flat}
设 $R$ 是分式域为 $K$ 的整环。若 $R[[x]]$ 在 $R[x]$ 上平坦，
则 $R$ 正规，当且仅当 $R$ 完全正规
（《代数》中的定义 \ref{algebra-definition-almost-integral}）。
\end{lemma}

\begin{proof}
假设存在 $\alpha \in K$ 与非零元 $r \in R$，使得对所有
$n \geq 1$ 都有 $r \alpha^n \in R$。在 $R[[x]]$ 中考虑
$f = \sum r \alpha^{n - 1} x^n$。写 $\alpha = a/b$，
其中 $a, b \in R$ 且 $b$ 非零。于是 $(a x - b)f = -rb$。
故 $rb$ 属于理想 $(ax - b)R[[x]]$。
令 $S = \{h \in R[x] : h(0) = 1\}$。这是乘法子集，
而 $R[x] \to R[[x]]$ 的平坦性蕴含
$S^{-1}R[x] \to R[[x]]$ 忠实平坦
（细节从略；提示：使用《代数》中的引理
\ref{algebra-lemma-ff-rings}）。因此
$$
S^{-1}R[x]/(ax - b)S^{-1}R[x] \to R[[x]]/(ax - b)R[[x]]
$$
是单射。由此存在 $h \in S$ 与 $g \in R[x]$，使得
$h rb = (ax - b) g$。写成
$h = 1 + h_1 x + \ldots + h_d x^d$，便得到
$$
1 + h_1 x + \ldots + h_d x^d = (1/r)(\alpha x - 1)g
$$
这个在 $K[x]$ 中的分解给出 $K[x^{-1}]$ 中相应的分解，
从而说明 $\alpha$ 是一个系数在 $R$ 中的首一多项式的根，正合所需。
\end{proof}

\begin{lemma}
\label{examples-lemma-completion-polynomial-ring-not-flat-bis}
若 $R$ 是维数 $> 1$ 的赋值环，则 $R[[x]]$ 在 $R$ 上平坦，
但在 $R[x]$ 上不平坦。
\end{lemma}

\begin{proof}
上面的论证表明，只要证明 $R$ 不是完全正规的，结论便成立
（赋值环是正规的；见《代数》中的引理
\ref{algebra-lemma-valuation-ring-normal}）。
取素理想链 $\mathfrak p \subset \mathfrak m \subset R$。
选取非零元 $x \in \mathfrak p$ 与
$y \in \mathfrak m \setminus \mathfrak p$。
则对所有 $n \geq 1$，都有 $x y^{-n} \in R$
（否则 $y^n/x \in R$，这与 $y \not \in \mathfrak p$ 矛盾）。
因此，$1/y$ 在 $R$ 上几乎整，但不属于 $R$。
\end{proof}

\noindent
下面构造一个局部化的完备化不平坦的例子。为此，考虑环
$$
R = k[y, z, a_1, a_2, a_3, \ldots]/(ya_i, a_i a_j)
$$
以 $f \in R$ 表示 $z$ 的剩余类。我们断言环映射
\begin{equation}
\label{examples-equation-nonflat}
R[[x]] \longrightarrow R_f[[x]]
\end{equation}
不平坦。令 $I$ 为 $y : R[[x]] \to R[[x]]$ 的核。
$I$ 的典型元素 $g$ 形如 $g = \sum g_{n, m} a_mx^n$，
其中 $g_{n, m} \in k[z]$，并且对给定的 $n$，
非零的 $g_{n, m}$ 只有有限多个。令 $J$ 为
$y : R_f[[x]] \to R_f[[x]]$ 的核。我们断言
$J \not = I R_f[[x]]$。事实上，若二者相等，则有
$$
\sum z^{-n} a_n x^n = \sum\nolimits_{i = 1, \ldots, m} h_i g_i
$$
其中某个 $m \geq 1$、$g_i \in I$，且 $h_i \in R_f[[x]]$。
设
$h_i = \bar h_i \bmod (y, a_1, a_2, a_3, \ldots)$，
其中 $\bar h_i \in k[z, 1/z][[x]]$。考察 $a_n$ 的系数，
并使用上面对元素 $g_i$ 的描述，将得到
$$
z^{-n} x^n = \sum \bar h_i \bar g_{i, n}
$$
其中某些 $\bar g_{i, n} \in k[z][[x]]$。这意味着所有
$z^{-n}x^n$ 都包含在由元素 $\bar h_i$ 生成的有限
$k[z][[x]]$-模中。由于 $k[z][[x]]$ 是 Noether 环，
这蕴含 $k[z, 1/z][[x]]$ 中由
$1, z^{-1}x, z^{-2}x^2, \ldots$ 生成的
$R[z][[x]]$-子模是有限的。
% P12-ERR-0322
由《代数》中的引理 \ref{algebra-lemma-characterize-integral-element}，
可得 $z^{-1}x$ 在 $k[z][[x]]$ 上整，矛盾。另一方面，
若 (\ref{examples-equation-nonflat}) 平坦，则把正合列
$0 \to I \to R[[x]] \xrightarrow{y} R[[x]]$ 与 $R_f[[x]]$ 张量，
将得到 $J = IR_f[[x]]$。

\begin{lemma}
\label{examples-lemma-nonflat-completion-localization}
存在关于主理想 $I$ 完备的环 $A$ 以及元素 $f \in A$，
使得 $A_f$ 的 $I$-进完备化 $A_f^\wedge$ 在 $A$ 上不平坦。
\end{lemma}

\begin{proof}
令 $A = R[[x]]$、$I = (x)$，并注意 $R_f[[x]]$
是 $R[[x]]_f$ 的完备化。
\end{proof}



\section{拟凝聚模的非 Abel 范畴}
\label{examples-section-nonabelian-QCoh}

\noindent
在《叠上的层》第 \ref{stacks-sheaves-section-quasi-coherent} 节中，
我们定义了 $\Sch$ 上群胚纤维化范畴上的拟凝聚模范畴。
虽然我们在《叠上的层》第
\ref{stacks-sheaves-section-quasi-coherent-algebraic-stacks} 节中证明，
对于代数叠，该范畴是 Abel 范畴；但本节将说明，
对于形式代数空间，情况并非如此。

\medskip\noindent
具体而言，把 $\mathbf{Z}_p$ 视为带有 $p$-进拓扑的拓扑环。
令 $X = \text{Spf}(\mathbf{Z}_p)$；见《形式空间》中的定义
\ref{formal-spaces-definition-affine-formal-spectrum}。
则 $X$ 是 $(\Sch/\mathbf{Z})_{fppf}$ 上的集合层，并给出一个
setoid 叠 $\mathcal{X}$；见《叠》中的引理
\ref{stacks-lemma-when-stack-in-sets}。因此，《叠上的层》第
\ref{stacks-sheaves-section-quasi-coherent-algebraic-stacks} 节的讨论适用。

\medskip\noindent
设 $\mathcal{F}$ 是 $\mathcal{X}$ 上的拟凝聚模。
由于 $X = \colim \Spec(\mathbf{Z}/p^n\mathbf{Z})$，由《叠上的层》中的引理
\ref{stacks-sheaves-lemma-quasi-coherent}，显然 $\mathcal{F}$
由序列 $(\mathcal{F}_n)$ 给出，其中
\begin{enumerate}
\item $\mathcal{F}_n$ 是 $\Spec(\mathbf{Z}/p^n\mathbf{Z})$ 上的拟凝聚模；并且
\item 转移映射给出同构
$\mathcal{F}_n = \mathcal{F}_{n + 1}/p^n\mathcal{F}_{n + 1}$。
\end{enumerate}
转化为模，可见 $\mathcal{F}$ 对应于系统 $(M_n)$，
其中每个 $M_n$ 都是被 $p^n$ 零化的 Abel 群，而转移映射诱导同构
$M_n = M_{n + 1}/p^n M_{n + 1}$。在此情形下，模
$M = \lim M_n$ 是 $p$-进完备模，且 $M_n = M/p^n M$；见
《代数》中的引理 \ref{algebra-lemma-limit-complete}。
由此可知，$X$ 上的拟凝聚模范畴等价于 $p$-进完备 Abel 群范畴。
该范畴不是 Abel 范畴；见第 \ref{examples-section-non-abelian} 节。

\begin{lemma}
\label{examples-lemma-quasi-coherent-not-abelian}
形式代数空间 $X$ 上的拟凝聚\footnote{此处采用上面所定义的拟凝聚模。
% P12-ERR-0323
鉴于 Stacks Project 中的设置方式，这确实是正确的定义；
如上所见，我们的定义与人们对 $\text{Spf}(A)$ 上拟凝聚模的朴素定义一致，
即完备 $A$-模。}模范畴一般不是 Abel 范畴，
即使 $X$ 是 Noether 仿射形式代数空间亦然。
\end{lemma}

\begin{proof}
见上面的讨论。
\end{proof}





\section{局部可分裂的不可分裂序列}
\label{examples-section-nonsplit-locally-split}

\noindent
考虑短正合序列
$$
0 \to M \to
\bigoplus\nolimits_{p\text{ 为素数}} \mathbf{Z}_{(p)} \to \mathbf{Q} \to 0
$$
其中 $M$ 是右侧映射的核，而该映射在每个直和项上均取嵌入
$\mathbf{Z}_{(p)} \to \mathbf{Q}$。这个 $\mathbf{Z}$-模序列不可分裂，
因为不存在非零同态 $\mathbf{Q} \to \mathbf{Z}_{(p)}$。另一方面，
若在任意素数 $p$ 处局部化，该序列便成为可分裂的。

\begin{lemma}
\label{examples-lemma-nonsplit-locally-split}
存在环 $R$ 以及一个不可分裂的模序列，使它在 Zariski 局部成为可分裂的。
\end{lemma}

\begin{proof}
见上述讨论。
\end{proof}












































\section{正则序列与基变换}
\label{examples-section-regular-base-change}

\noindent
我们将构造一个带有正则序列 $(x, y, z)$ 的环 $R$，使得存在非零元素
$\delta \in R/zR$ 满足 $x\delta = y\delta = 0$。

\medskip\noindent
为构造这个例，我们先在环 $k[x, y, z]$ 上构造一个特殊的模 $E$，
其中 $k$ 是任意域。具体地，$E$ 是下图中的推出：
$$
\xymatrix{
\frac{xk[x, y, z, y^{-1}]}{xyk[x, y, z]} \ar[r] \ar[d]^{z/x} &
\frac{k[x, y, z, x^{-1}, y^{-1}]}{yk[x, y, z, x^{-1}]} \ar[r] \ar[d] &
\frac{k[x, y, z, x^{-1}, y^{-1}]}{yk[x, y, z, x^{-1}] + xk[x, y, z, y^{-1}]}
\ar[d] \\
\frac{k[x, y, z, y^{-1}]}{yzk[x, y, z]} \ar[r] &
E \ar[r] &
\frac{k[x, y, z, x^{-1}, y^{-1}]}{yk[x, y, z, x^{-1}] + xk[x, y, z, y^{-1}]}
}
$$
其中各行都是短正合序列（由于排版上的困难，我们省略了两端的零）。
也可将 $E$ 描述为
$$
E = \{(f, g) \mid f \in k[x, y, z, x^{-1}, y^{-1}],
g \in k[x, y, z, y^{-1}] \}/\sim
$$
其中 $(f, g) \sim (f', g')$ 当且仅当存在
$h \in k[x, y, z, y^{-1}]$ 使得
$$
f = f' + xh \bmod yk[x, y, z, x^{-1}], \quad
g = g' - zh \bmod yzk[x, y, z]
$$
我们断言：(a) $x : E \to E$ 是单射；(b)
$y : E/xE \to E/xE$ 是单射；(c) $E/(x, y)E = 0$；(d) 存在
非零元素 $\delta \in E/zE$ 使得 $x\delta = y\delta = 0$。

\medskip\noindent
为证明 (a)，设 $(f, g)$ 是给出 $E$ 中一个元素的一对，并且
$(xf, xg) \sim 0$。则存在 $h \in k[x, y, z, y^{-1}]$，使得
$xf + xh \in yk[x, y, z, x^{-1}]$ 且
$xg - zh \in yzk[x, y, z]$。可以假设
$h = \sum a_{i, j, k}x^iy^jz^k$ 是仅出现 $j \leq 0$ 的单项式之和。
于是 $xg - zh \in yzk[x, y, z]$ 蕴含只出现 $i > 0$，即对某个
$h' \in k[x, y, z, y^{-1}]$ 有 $h = xh'$。
此时 $(f, g) \sim (f + xh', g - zh')$，故可设 $g = 0$ 且 $h = 0$。
在这种情形下，$xf \in yk[x, y, z, x^{-1}]$ 蕴含
$f \in yk[x, y, z, x^{-1}]$，从而 $(f, g) \sim 0$。
因此 $x : E \to E$ 是单射。

\medskip\noindent
由于乘以 $x$ 在
$\frac{k[x, y, z, x^{-1}, y^{-1}]}{yk[x, y, z, x^{-1}]}$ 上是同构，
可见 $E/xE$ 同构于
$$
\frac{k[x, y, z, y^{-1}]}{
yzk[x, y, z] + xk[x, y, z, y^{-1}] + zk[x, y, z, y^{-1}]}
=
\frac{k[x, y, z, y^{-1}]}{xk[x, y, z, y^{-1}] + zk[x, y, z, y^{-1}]}
$$
因而乘以 $y$ 在 $E/xE$ 上是同构。这显然蕴含 (b) 和 (c)。

\medskip\noindent
令 $e \in E$ 为 $(1, 0)$ 的等价类。假设 $e \in zE$。
则存在 $f \in k[x, y, z, x^{-1}, y^{-1}]$、
$g \in k[x, y, z, y^{-1}]$ 和 $h \in k[x, y, z, y^{-1}]$，使得
$$
1 + zf + xh \in yk[x, y, z, x^{-1}], \quad
0 + zg - zh \in yzk[x, y, z].
$$
这是不可能的：单项式 $1$ 既不可能出现在 $zf$ 中，也不可能出现在
$xh$ 中。另一方面，$ye = 0$，且
$xe = (x, 0) \sim (0, -z) = z(0, -1)$。因此，令 $\delta$ 等于
$e$ 在 $E/zE$ 中的同余类，即得到 (d)。

\begin{lemma}
\label{examples-lemma-strange-regular-sequence}
存在局部环 $R$ 和（位于极大理想中的）正则序列 $x, y, z$，
使得存在非零元素 $\delta \in R/zR$ 满足
$x\delta = y\delta = 0$。
\end{lemma}

\begin{proof}
令 $R = k[x, y, z] \oplus E$，其中把上述模 $E$ 视为平方零理想。
显然，$x, y, z$ 是 $R$ 中的正则序列，而元素
$\delta \in E/zE \subset R/zR$ 具有所需性质。
为得到局部的例，可将 $R$ 在极大理想 $\mathfrak m = (x, y, z, E)$ 处局部化。
序列 $x, y, z$ 仍为正则序列（因为局部化正合），而元素 $\delta$
仍非零，因为它的支撑位于 $\mathfrak m$。
\end{proof}

\begin{lemma}
\label{examples-lemma-base-change-regular-sequence}
存在局部环的局部同态 $A \to B$，以及 $B$ 的极大理想中的正则序列
$x, y$，使得 $B/(x, y)$ 在 $A$ 上平坦，但 $x, y$ 在
$B/\mathfrak m_AB$ 中的像 $\overline{x}, \overline{y}$ 不构成正则序列，
甚至不构成 Koszul 正则序列。
\end{lemma}

\begin{proof}
令 $A = k[z]_{(z)}$ 且
$B = (k[x, y, z] \oplus E)_{(x, y, z, E)}$。
由 $x, y, z$ 是 $B$ 中的正则序列（见引理
\ref{examples-lemma-strange-regular-sequence} 的证明），可知 $x, y$ 是 $B$ 中的
正则序列，并且 $B/(x, y)$ 是无挠 $A$-模，因而平坦。
另一方面，存在非零元素
$\delta \in B/\mathfrak m_AB = B/zB$ 被
$\overline{x}, \overline{y}$ 所零化。因此
$H_2(K_\bullet(B/\mathfrak m_AB, \overline{x}, \overline{y})) \not = 0$。
故 $\overline{x}, \overline{y}$ 不是 Koszul 正则的，特别地，它不是正则序列；
见 More on Algebra，引理
\ref{more-algebra-lemma-regular-koszul-regular}。
\end{proof}





\section{无限维 Noether 环}
\label{examples-section-Noetherian-infinite-dimension}

\noindent
正如我们在《代数》命题 \ref{algebra-proposition-dimension} 中所见，
Noether 局部环具有有限维数。不过，也存在无限维的 Noether 环。
见 \cite[附录，例 1]{Nagata}。

\medskip\noindent
具体地，设 $k$ 为域，并考虑环
$$
R = k[x_1, x_2, x_3, \ldots ].
$$
对 $i = 1, 2, \ldots$，令
$\mathfrak p_i = (x_{2^{i - 1}}, x_{2^{i - 1} + 1}, \ldots, x_{2^i - 1})$；
它们是 $R$ 的素理想。令 $S$ 为乘法子集
$$
S = \bigcap\nolimits_{i \geq 1} (R \setminus \mathfrak p_i).
$$
考虑环 $A = S^{-1}R$。我们断言：
\begin{enumerate}
\item 环 $A$ 的极大理想恰为理想 $\mathfrak m_i = \mathfrak p_iA$。
% P12-ERR-0324
\item 有 $A_{\mathfrak m_i} = R_{\mathfrak p_i}$，它是维数为 $2^i$ 的
Noether 局部环。
\item 环 $A$ 是 Noether 环。
\end{enumerate}
因此，这显然就是我们要找的例。细节从略。





\section{完备化非约化的局部环}
\label{examples-section-local-completion-nonreduced}

\noindent
在《代数》例 \ref{algebra-example-bad-dvr-char-p} 中，我们给出了一个
特征为 $p$、维数为 $1$ 的 Noether 局部整环 $R$，其完备化非约化。
本节给出 \cite[命题 3.1]{Ferrand-Raynaud} 中的例，它在特征零时给出类似的环。

\medskip\noindent
令 $\mathbf{C}\{x\}$ 为复数域 $\mathbf{C}$ 上的收敛幂级数环。
全体幂级数环 $\mathbf{C}[[x]]$ 是它的完备化。令
$K = \mathbf{C}\{x\}[1/x]$ 为收敛 Laurent 级数域。
$K$ 在 $\mathbf{C}$ 上的代数微分 $K$-模
$\Omega_{K/\mathbf{C}}$ 是无限维 $K$-向量空间（证明从略）。
我们可以选取 $f_n \in x\mathbf{C}\{x\}$（$n \geq 1$），使得
$
\text{d}x, \text{d}f_1, \text{d}f_2, \ldots
$
构成 $\Omega_{K/\mathbf{C}}$ 的某个基的一部分。因此，可以找到一个
$\mathbf{C}$-导子
$$
D : \mathbf{C}\{x\} \longrightarrow \mathbf{C}((x))
$$
% P12-ERR-0325
使得 $D(x) = 0$ 且 $D(f_i) = x^{-n}$。令
$$
A = \{f \in \mathbf{C}\{x\} \mid D(f) \in \mathbf{C}[[x]]\}
$$
我们断言：
\begin{enumerate}
\item $\mathbf{C}\{x\}$ 在 $A$ 上整；
\item $A$ 是局部整环；
\item $\dim(A) = 1$；
\item $A$ 的极大理想由 $x$ 和 $xf_1$ 生成；
\item $A$ 是 Noether 环；并且
\item $A$ 的完备化等于 $\mathbf{C}[[x]]$ 上的对偶数环。
\end{enumerate}
由于对偶数环非约化，环 $A$ 给出了所需的例。

\medskip\noindent
注意，若 $0 \not = f \in x\mathbf{C}\{x\}$，则对某个 $n \geq 0$
和 $h \in \mathbf{C}[[x]]$，可写成 $D(f) = h/f^n$。
于是 $D(f^{n + 1}/(n + 1)) \in \mathbf{C}[[x]]$ 且
$D(f^{n + 2}/(n + 2)) \in \mathbf{C}[[x]]$。因此可见
$f^{n + 1}, f^{n + 2} \in A$！特别地，(1) 成立。
我们还得出 $A$ 的分式域等于 $\mathbf{C}\{x\}$ 的分式域。
同时立即可见，$A \cap x\mathbf{C}\{x\}$ 是 $A$ 的非单位元集合，
故 $A$ 是维数为 $1$ 的局部整环。若能证明 (4)，便可推出 $A$ 是
Noether 环（证明从略）。设 $f \in A \cap x\mathbf{C}\{x\}$。
写成 $D(f) = h$，其中 $h \in \mathbf{C}[[x]]$。
再写成 $h = c + xh'$，其中 $c \in \mathbf{C}$、
$h' \in \mathbf{C}[[x]]$。则
$D(f - cxf_1) = c + xh' - c = xh'$。另一方面，
$f - cxf_1 = xg$，其中 $g \in \mathbf{C}\{x\}$；但由上述计算，
$D(g) = h' \in \mathbf{C}[[x]]$，因而 $g \in A$。
所以如所需，$f = cxf_1 + xg \in (x, xf_1)$。

\medskip\noindent
最后，为什么 $A$ 的完备化非约化？以 $\hat A$ 表示 $A$ 的完备化。
当然，它满射到 $\mathbf{C}\{x\}$ 的完备化 $\mathbf{C}[[x]]$，因为
$x \in A$。将此映射记为 $\psi : \hat A \to \mathbf{C}[[x]]$。
上面已经看到 $\mathfrak m_A = (x, xf_1)$，于是经简单计算有
$D(\mathfrak m_A^n) \subset (x^{n - 1})$。因此
$D : A \to \mathbf{C}[[x]]$ 连续，并在 $\psi$ 上诱导连续导子
$\hat D : \hat A \to \mathbf{C}[[x]]$。由此得到环映射
$$
\psi + \epsilon \hat D :
\hat A
\longrightarrow
\mathbf{C}[[x]][\epsilon].
$$
由于 $\hat A$ 是一维 Noether 完备局部环，只要证明该箭头满射，
便可推出 $\hat A$ 非约化。事实上该映射是同构，但我们省略对此的验证。
子环 $\mathbf{C}[x]_{(x)} \subset A$ 在完备化上诱导映射
$i : \mathbf{C}[[x]] \to \hat A$，使得
% P12-ERR-0326
$i \circ \psi = \text{id}$ 且 $D \circ i = 0$
（因为按构造 $D(x) = 0$）。考虑元素 $x^nf_n \in A$。有
$$
(\psi + \epsilon D)(x^nf_n) = x^n f_n + \epsilon
$$
对所有 $n \geq 1$ 均成立。由这些说明很容易推出满射性。





\section{另一个完备化非约化的局部环}
\label{examples-section-another-local-completion-nonreduced}

\noindent
本节构造一个维数为 $2$、关于某个主理想完备的 Noether 局部整环，
使得其一个局部化的再完备化非约化。

\medskip\noindent
设 $p$ 为素数。设 $k$ 为特征 $p$ 的域，并且 $k$ 在其 $p$ 次幂子域
$k^p$ 上的次数无限。例如 $k = \mathbf{F}_p(t_1, t_2, t_3, \ldots)$。
考虑环
$$
A =
\left\{
\begin{matrix}
\sum a_{i, j} x^iy^j \in k[[x, y]] \text{，满足对所有 }n \geq 0
\text{ 都有 } \\
[k^p(a_{n, n}, a_{n, n + 1}, a_{n + 1, n},
a_{n, n + 2}, a_{n + 2, n}, \ldots) : k^p] < \infty
\end{matrix}
\right\}
$$
作为集合，有
$$
k^p[[x, y]] \subset A  \subset k[[x, y]]
$$
$A$ 的每个元素 $f$ 都能唯一写成级数
$$
f = f_0 + f_1 xy + f_2 (xy)^2 + f_3 (xy)^3 + \ldots
$$
其中
$$
f_n = a_{n, n} + a_{n, n + 1} y + a_{n + 1, n} x +
a_{n, n + 2} y^2 + a_{n + 2, n} x^2 + \ldots
$$
而定义 $A$ 的公式中的条件意味着 $f_n$ 的系数生成 $k^p$ 的有限扩张。
从这一表示可清楚看出，$A$ 是 $k[[x, y]]$ 的一个
$k^p[[x, y]]$-子代数，并且关于理想 $xy$ 完备。此外，显然有
$$
A/xy A = C \times_k D
$$
其中 $k^p[[x]] \subset C \subset k[[x]]$ 和
$k^p[[y]] \subset D \subset k[[y]]$ 是《代数》例
\ref{algebra-example-bad-dvr-char-p} 中的幂级数子环。
因此，$C$ 和 $D$ 都是离散赋值环，从而 $A/ xy A$ 是 Noether 环。
由《代数》引理 \ref{algebra-lemma-completion-Noetherian}，可得 $A$ 是
Noether 环。由于 $\dim(k[[x, y]]) = 2$，再用《代数》引理
\ref{algebra-lemma-integral-sub-dim-equal}，可得 $\dim(A) = 2$。

\medskip\noindent
设 $f = \sum a_i x^i$ 为幂级数，使得
$k^p(a_0, a_1, a_2, \ldots)$ 在 $k^p$ 上的次数无限。
则 $f \not \in A$，但 $f^p \in A$。令
$$
B = A[f] \subset k[[x, y]]
$$
由于 $B$ 在 $A$ 上有限，$B$ 是 Noether 环。又因为 $B$ 关于由 $xy$
生成的理想完备；见《代数》引理 \ref{algebra-lemma-completion-tensor}。
事实上，$B$ 是以 $1, f, f^2, \ldots, f^{p - 1}$ 为基的自由 $A$-模；
证明从略。

\medskip\noindent
我们断言环
$$
(B_y)^\wedge = (B[1/y])^\wedge = \lim B[1/y]/(xy)^n B[1/y] =
\lim B[1/y] / x^n B[1/y]
$$
非约化。具体地，该环在
$$
(A_y)^\wedge = (A[1/y])^\wedge = \lim A[1/y]/(xy)^n A[1/y] =
\lim A[1/y] / x^n A[1/y]
$$
上自由，基为 $1, f, \ldots, f^{p - 1}$。然而，存在元素
$g \in (A_y)^\wedge$ 使得 $f^p = g^p$。事实上，只需取
$g = \sum a_i x^i$（与定义 $f$ 时使用相同的表达式），
该表达式在 $(A_y)^\wedge$ 中有意义。因此
$$
(B_y)^\wedge = (A_y)^\wedge[f]/(f^p - g^p) \cong
(A_y)^\wedge[\tau]/(\tau^p)
$$
非约化。事实上，这个例还说明稍强的结论。注意，$(A_y)^\wedge$
是以 $x$ 为一致化参数的离散赋值环，其剩余域是上述离散赋值环 $D$
的分式域。因此，甚至
$$
(B_y)^\wedge[1/(xy)] = ((B_y)^\wedge)_{xy}
$$
也非约化。由此得到如下类型的例。

\begin{lemma}
\label{examples-lemma-nonreduced-recompletion}
存在一个关于主理想 $I = (b)$ 完备的维数为 $2$ 的 Noether 局部整环
$(B, \mathfrak m)$，以及元素 $f \in \mathfrak m$、$f \not \in I$，
使得主元局部化 $B_f$ 的 $I$-进完备化
$C = (B_f)^\wedge$ 非约化，而且
$C_b = C[1/b] = (B_f)^\wedge[1/b]$ 也非约化。
\end{lemma}

\begin{proof}
见上述讨论。
\end{proof}





\section{一个非链状 Noether 局部环}
\label{examples-section-non-catenary-Noetherian-local}

\noindent
尽管 Noether 局部环有一套成功的维数理论，仍然存在非链状的 Noether
局部环。一个例见 \cite[附录，例 2]{Nagata}。事实上，我们将在最简单的
情形下给出这个例。具体地，我们将构造一个维数为 $2$、但并非泛链状的
Noether 局部整环 $A$。（注意，$A$ 自动是链状的；见《练习》中的练习
\ref{exercises-exercise-Noetherian-local-domain-dim-2-catenary}。）
一个并非泛链状的 Noether 局部环的存在性蕴含非链状 Noether 局部环的
存在性；在本节末尾，我们将对当前这个具体例说明这一点。

\medskip\noindent
设 $k$ 为域，并考虑 $k$ 上一个变量的形式幂级数环 $k[[x]]$。令
$$
z = \sum\nolimits_{i = 1}^\infty a_i x^i
$$
为形式幂级数。我们假设，作为 Laurent 级数域
$k((x)) = k[[x]][1/x]$ 的元素，$z$ 在 $k(x)$ 上超越。令
$$
z_j
=
x^{-j}(z - \sum\nolimits_{i = 1, \ldots, j - 1} a_i x^i)
=
\sum\nolimits_{i = j}^\infty a_i x^{i - j}
\in k[[x]].
$$
注意 $z = xz_1$。令 $R$ 为 $k[[x]]$ 中由 $x$、$z$ 以及所有 $z_j$
生成的子环，换言之，
$$
R = k[x, z_1, z_2, z_3, \ldots ] \subset k[[x]].
$$
考虑 $R$ 的理想 $\mathfrak m = (x)$ 和
$\mathfrak n = (x - 1, z_1, z_2 + a_1, z_3 + a_1 + a_2, \ldots)$。

\medskip\noindent
有 $xz_{j + 1} + a_j = z_j$。因此 $R/\mathfrak m = k$，且
$\mathfrak m$ 是极大理想。此外，$R$ 中不属于 $\mathfrak m$ 的任意元素
在 $k[[x]]$ 中映为单位元，故 $R_{\mathfrak m} \subset k[[x]]$。
% P12-ERR-0327
事实上，很容易推出 $R_{\mathfrak m}$ 是离散赋值环，且剩余域为 $k$。

\medskip\noindent
我们断言
$$
R/(x - 1) =
k[x, z_1, z_2, z_3, \ldots ]/(x - 1)
\cong
k[z].
$$
事实上，上述关系蕴含
$z_{j + 1} = z_j - a_j - (x - 1)z_{j + 1}$，因此在商
$R/(x - 1)$ 中，可以用 $z_j$ 表示 $z_{j + 1}$ 的类。
按构造，$R$ 的分式域在 $k$ 上的超越次数为 $2$，所以在
$R/(x - 1)$ 中 $z$ 在 $k$ 上超越，从而得到所需同构。
于是 $\mathfrak n = (x - 1, z)$，并且它是极大理想。事实上，映射
$$
k[x, x^{-1}, z]_{(x - 1, z)} \longrightarrow R_{\mathfrak n}
$$
是同构（因为 $x^{-1}$ 在 $R_{\mathfrak n}$ 中可逆，并且
$z_{j + 1} = x^{-1}z_j - a_j = \ldots = f_j(x, x^{-1}, z)$）。
这说明 $R_{\mathfrak n}$ 是维数为 $2$、剩余域为 $k$ 的正则局部环。

\medskip\noindent
令 $S$ 为乘法子集
$$
S =
(R \setminus \mathfrak m) \cap (R \setminus \mathfrak n) =
R \setminus (\mathfrak m \cup \mathfrak n)
$$
并令 $B = S^{-1}R$。我们断言：
\begin{enumerate}
\item 环 $B$ 是 $k$-代数。
\item 环 $B$ 的极大理想恰为两个理想 $\mathfrak mB$ 和 $\mathfrak nB$。
\item 在这两个极大理想处的剩余域均为 $k$。
\item 有 $B_{\mathfrak mB} = R_{\mathfrak m}$ 和
$B_{\mathfrak nB} = R_{\mathfrak n}$；它们分别是维数为 $1$ 和 $2$ 的
Noether 正则局部环。
\item 环 $B$ 是 Noether 环。
\end{enumerate}
验证细节从略。

\medskip\noindent
每当给定一个具有上述性质的 $k$-代数 $B$ 时，我们可以如下得到一个例。
取 $A = k + \text{rad}(B) \subset B$，其中
$\text{rad}(B) = \mathfrak mB \cap \mathfrak nB$ 是 Jacobson 根。
容易看出 $B$ 在 $A$ 上有限，因而由 Eakin 定理，$A$ 是 Noether 环
（见 \cite{Eakin}，或 \cite[附录 A1]{Nagata}，或
% P12-ERR-0328
在此插入将来的参考文献）。又因为 $A$ 是与 $B$ 具有相同分式域和剩余域
$k$ 的局部整环，而 $B$ 的维数为 $2$，由《代数》引理
\ref{algebra-lemma-integral-sub-dim-equal} 可知 $A$ 的维数也为 $2$。

\medskip\noindent
若 $A$ 是泛链状的，则维数公式《代数》引理
\ref{algebra-lemma-dimension-formula} 将给出
$\dim(B_{\mathfrak mB}) = 2$，矛盾。

\medskip\noindent
注意，$B$ 由一个元素在 $A$ 上生成。因此，对 $A[x]$ 的某个素理想
$\mathfrak p$，有 $B = A[x]/\mathfrak p$。令
$\mathfrak m' \subset A[x]$ 为对应于 $\mathfrak mB$ 的极大理想。
则一方面 $\dim(A[x]_{\mathfrak m'}) = 3$，另一方面
$$
(0)
\subset \mathfrak pA[x]_{\mathfrak m'}
\subset \mathfrak m'A[x]_{\mathfrak m'}
$$
是素理想的极大链。因此，$A[x]_{\mathfrak m'}$ 是一个非链状
Noether 局部环的例。





\section{不良 Noether 局部环的存在性}
\label{examples-section-bad}

\noindent
设 $(A, \mathfrak m, \kappa)$ 为 Noether 完备局部环。
文献 \cite{Lech} 证明：若 $\text{depth}(A) \geq 1$，并且 $A$ 包含
$\mathbf{Q}$ 或 $\mathbf{F}_p$ 作为子环，或者包含 $\mathbf{Z}$ 作为子环
且 $A$ 作为 $\mathbf{Z}$-模无挠，则 $A$ 是某个 Noether 局部整环的完备化。
这产生了许多具有“奇异”性质的 Noether 局部整环的例。

\medskip\noindent
例如，将此应用于 $A = \mathbf{C}[[x, y]]/(y^2)$，便得到一个完备化
非约化的 Noether 局部整环。请与第
\ref{examples-section-local-completion-nonreduced} 节比较。

\medskip\noindent
文献 \cite{LLPY} 给出了刻画 $A$ 何时为某个约化 Noether 局部环之
完备化的条件。

\medskip\noindent
文献 \cite{Heitmann-completion-UFD} 证明：若
$\text{depth}(A) \geq 2$，并且 $A$ 包含 $\mathbf{Q}$ 或
$\mathbf{F}_p$ 作为子环，或者包含 $\mathbf{Z}$ 作为子环且 $A$ 作为
$\mathbf{Z}$-模无挠，则 $A$ 是某个 Noether 局部唯一分解整环 $R$ 的
完备化。特别地，$R$ 正规（《代数》引理
\ref{algebra-lemma-UFD-normal}），因而 $R$ 的 Hensel 化也是正规整环
（More on Algebra，引理 \ref{more-algebra-lemma-henselization-normal}）。
所以，上述 $A$ 是某个 Hensel Noether 局部正规整环的完备化
（因为 $R$ 与其 Hensel 化的完备化相同；见 More on Algebra，引理
\ref{more-algebra-lemma-henselization-noetherian}）。

\medskip\noindent
应用这一结果，可找到一个 Noether 局部唯一分解整环 $R$，使得
$R^\wedge \cong \mathbf{C}[[x, y, z, w]]/(wx, wy)$。
注意，$\Spec(R^\wedge)$ 是一个维数为 $2$ 的正则分支和一个维数为 $3$ 的正则分支的并。
环 $R$ 不可能泛链状：令
$$
X \longrightarrow \Spec(R)
$$
为极大理想的爆破。于是 $X$ 是整概形。存在闭点 $x \in X$，使得
$\dim(\mathcal{O}_{X, x}) = 2$；具体地，在完备局部环的层面上，
我们选取 $x$ 位于维数为 $2$ 的分支的严格变换上，而不位于维数为 $3$ 的分支的严格变换上。
由《态射》引理 \ref{morphisms-lemma-dimension-formula}，可见 $R$ 并非
泛链状。请与第 \ref{examples-section-non-catenary-Noetherian-local} 节比较。

\medskip\noindent
上述环是链状的（因为它是维数为 $3$ 的 Noether 局部唯一分解整环）。
不过，在 \cite{Ogoma-example} 中，作者构造了一个 Noether 局部正规整环
$R$，满足 $R^\wedge \cong \mathbf{C}[[x, y, z, w]]/(wx, wy)$，
而 $R$ 不是链状的。另见 \cite{Heitmann-Ogoma} 和 \cite{Lech-YAPO}。

\medskip\noindent
文献 \cite{Heitmann-isolated} 证明：若 $A$ 包含 $\mathbf{Q}$ 或
$\mathbf{F}_p$ 作为子环，或者 $A$ 的剩余特征为 $p > 0$，且 $p$ 在
$A$ 的任意真局部化中都不可能映为非零零因子，则 $A$ 是某个具有
孤立奇点的 Noether 局部环 $R$ 的完备化。这里称 Noether 局部环 $R$
具有孤立奇点，是指对所有非极大素理想 $\mathfrak p \subset R$，
$R_\mathfrak p$ 都是正则局部环。

\medskip\noindent
论文 \cite{Nishimura-few} 和 \cite{Nishimura-few-II} 列出了许多具有
给定完备化的“不良” Noether 局部环。特别地，
% P12-ERR-0329
其中构造了一个维数为 $2$ 的 Nagata 局部正规整环，其完备化为
$\mathbf{C}[[x, y, z]]/(yz)$；还构造了一个完备化为
$\mathbf{C}[[x, y, z]]/(y^2 - z^3)$ 的例。

\medskip\noindent
顺便一提，文献 \cite{Loepp} 证明：若 $A$ 约化、等维，且 $A$ 中没有
整数是零因子，则 $A$ 是某个优秀 Noether 局部整环的完备化。
不过，这不会产生“不良” Noether 局部环，因为得到的是优秀环！





\section{Noether Jacobson 环中的维数}
\label{examples-section-noetherian-jacobson}

\noindent
设 $k$ 为有限域的代数闭包。令 $A = k[x, y]$ 且
$X = \Spec(A)$。令 $C = V(x)$ 为 $y$-轴（也可取 $X$ 中任意其他
维数为 $1$ 的整闭子概形）。把 $X$ 中其他维数为 $1$ 的整闭子概形
枚举为 $C_1, C_2, C_3, \ldots$，并把 $C$ 的闭点枚举为
$p_1, p_2, p_3, \ldots$。

\medskip\noindent
断言：对每个 $n$，存在维数为 $1$ 的不可约闭子集
$Z_n \subset X$，使得
$$
\{p_n\} = Z_n \cap (C \cup C_1 \cup \ldots \cup C_n)
$$
在集合论意义下成立。为此，令
$Y = C \cup C_1 \cup C_2 \cup \ldots \cup C_n$。
这是 $k$ 上维数为 $1$ 的约化仿射代数概形。只需找到
$f \in k[x, y]$，使得 $V(f) \cap  Y = \{p_n\}$ 在集合论意义下成立；
因为这时可取 $Z_n$ 为 $V(f)$ 的一个适当不可约分支。由于限制映射
$$
k[x, y] \longrightarrow \Gamma(Y, \mathcal{O}_Y)
$$
满射，只需在 $Y$ 上找到一个零点集在集合论意义下为 $\{p_n\}$ 的正则
函数 $g$。为说明这是可能的，选取支撑为 $p_n$ 的有效 Cartier 除子
$D \subset Y$（这可由《簇》引理
\ref{varieties-lemma-complement-codim-1-closed-points} 得到）。
% P12-ERR-0330
因此，只需证明对某个 $N > 0$ 有
$\mathcal{O}_X(ND) \cong \mathcal{O}_X$。不过，有限域代数闭包上的仿射
$1$-维代数概形的 Picard 群是挠群
% P12-ERR-0331
（在此插入将来的参考文献），故断言成立。

\medskip\noindent
对所有 $n$ 如上选取 $Z_n$。由于 $k[x, y]$ 是唯一分解整环，
可对某个不可约元素 $f_n \in A$ 写成 $Z_n = V(f_n)$。
令 $S \subset k[x, y]$ 为由 $f_1, f_2, f_3, \ldots$ 生成的乘法子集。
考虑 Noether 环 $B = S^{-1}A$。

\medskip\noindent
显然，环映射 $A \to B$ 识别局部环，并诱导单射
$\Spec(B) \to \Spec(A)$。此外，考察曲线 $C_1$ 可见，在从
$\Spec(A)$ 过渡到 $\Spec(B)$ 时，仅删去了 $C \cap C_1$ 中的点。
特别地，$\Spec(B)$ 有无穷多个对应于 $A$ 的极大理想的极大理想。
另一方面，$xB$ 是极大理想，因为 $B/xB$ 的谱仅含一个素理想：
我们删去了 $C = V(x)$ 的所有闭点（但没有删去泛点）。
最后，对 $i \geq 1$ 考虑曲线 $C_i$。取不可约元 $g_i \in A$，使得
$C_i = V(g_i)$。若对某个 $n$ 有 $C_i = Z_n$，则 $g_iB$ 是单位理想。
否则，在从 $A$ 过渡到 $B$ 时，$C_i$ 的闭点除有限多个外都保留下来：
具体地，从 $C_i$ 中仅删去了
$(Z_1 \cup \ldots \cup Z_{i - 1} \cup C) \cap C_i$ 中的点。

\medskip\noindent
上述 $B$ 的素谱结构及《代数》引理
\ref{algebra-lemma-noetherian-dim-1-Jacobson} 表明 $B$ 是 Jacobson 环。
极大理想包括位于 $\Spec(B)$ 中的 $A$ 的极大理想
% P12-ERR-0332
（而这样的理想有无穷多个），再加上极大理想 $xB$。
因此可见，出现了维数为 $1$ 和 $2$ 的局部环。

\begin{lemma}
\label{examples-lemma-Noetherian-Jacobson}
存在一个 Jacobson、泛链状的 Noether 整环 $B$，它有极大理想
$\mathfrak m_1, \mathfrak m_2$，并且
$\dim(B_{\mathfrak m_1}) = 1$ 且 $\dim(B_{\mathfrak m_2}) = 2$。
\end{lemma}

\begin{proof}
$B$ 的构造如上。只需指出，由《代数》引理
\ref{algebra-lemma-localization-catenary} 和《态射》引理
\ref{morphisms-lemma-ubiquity-uc}，$B$ 是泛链状的。
\end{proof}





\section{底空间为 Noether 而概形非 Noether}
\label{examples-section-noetherian-not-noetherian}

\noindent
我们给出两个例，说明底拓扑空间为 Noether 空间的概形未必是 Noether 概形。

\begin{example}
\label{examples-example-many-variables}
设 $k$ 为域，并令 $A = k[x_1, x_2, x_3, \dots] /
(x_1^2, x_2^2, x_3^2, \dots)$。$A$ 的任意素理想都包含幂零元
$x_1, x_2, x_3, \dots$，所以
$\mathfrak p = (x_1, x_2, x_3, \dots)$ 是 $A$ 唯一的素理想。
因此，$\operatorname{Spec} A$ 的底拓扑空间是单点，特别地是 Noether
空间。然而，$\mathfrak p$ 显然不是有限生成的。
\end{example}

\begin{example}
\label{examples-example-many-mononials}
设 $k$ 为域，并令 $A \subseteq k[x, y]$ 为由 $k$ 和单项式
$\{xy^i\}_{i \ge 0}$ 生成的子环。$A$ 中不包含 $x$ 的素理想与
$A_x \cong k[x, x^{-1}, y]$ 的素理想一一对应。若 $\mathfrak p$ 是
一个包含 $x$ 的素理想，则它包含每个 $xy^i$（$i \ge 0$），因为
$(xy^i)^2 = x(xy^{2i}) \in \mathfrak p$ 且 $\mathfrak p$ 是根理想。
因此，$\mathfrak p = (\{xy^i\}_{i \ge 0})$。所以，
$\operatorname{Spec} A$ 的底拓扑空间是 Noether 空间，因为它由
Noether 概形 $\Spec(A[x, x^{-1}, y])$ 的点以及素理想 $\mathfrak p$
组成。但是，环 $A$ 非 Noether，因为 $\mathfrak p$ 不是有限生成的。
注意，在这个例中，$A$ 还是整环。
\end{example}





\section{正规化拟仿射的非拟仿射簇}
\label{examples-section-nonquasi-affine}

\noindent
文献 \cite[II 注 6.6.13]{EGA} 提到了这类例的存在性。该文献为此类例
援引 EGA 第五卷，但第五卷并未出版。

\medskip\noindent
设 $k$ 为域。令 $Y = \mathbf{A}^2_k \setminus \{(0, 0)\}$。
我们将构造有限满双有理态射
$\pi : Y \longrightarrow X$，其中 $X$ 是 $k$ 上的簇，且 $X$ 非拟仿射。
具体地，考虑 $Y$ 中的下列曲线：
$$
\begin{matrix}
C_1 & : & x = 0 \\
C_2 & : & y = 0
\end{matrix}
$$
注意 $C_1 \cap C_2 = \emptyset$。选取同构
$\varphi : C_1 \to C_2$，$(0, y) \mapsto (y^{-1}, 0)$。
我们断言，存在唯一的上述态射 $\pi : Y \to X$，使得
$$
\xymatrix{
C_1
\ar@<1ex>[rr]^{\text{id}} \ar@<-1ex>[rr]_{\varphi}
& &
Y \ar[r]^\pi & X
}
$$
是簇范畴（甚至概形范畴）中的余等化子图。暂且接受这一点；我们来证明
这样的 $X$ 不可能拟仿射。事实上，显然有
$$
\Gamma(X, \mathcal{O}_X) =
\{ f \in k[x, y] \mid f(0, y) = f(y^{-1}, 0)\} =
k \oplus (xy) \subset k[x, y].
$$
特别地，这些函数不能分离点 $(1, 0)$ 和 $(-1, 0)$；它们在 $X$ 中的像
（如下所示）彼此不同（若 $k$ 的特征不是 $2$）。

\medskip\noindent
为证明 $X$ 存在，考虑 $Y$ 的 Zariski 开集 $D(x + y) \subset Y$。
它是环 $k[x, y, 1/(x + y)]$ 的谱，而曲线 $C_1$、$C_2$ 完全包含在
$D(x + y)$ 中。此外，态射
$$
C_1 \amalg C_2
\longrightarrow
D(x + y) \cap Y = \Spec(k[x, y, 1/(x + y)])
$$
是闭浸入。由 More on Algebra，引理
\ref{more-algebra-lemma-fibre-product-finite-type}，环
$$
A =
\{f \in k[x, y, 1/(x + y)] \mid f(0, y) = f(y^{-1}, 0)\}
$$
在 $k$ 上有限型。另一方面，$Y$ 的开集 $D(xy) \subset Y$ 与曲线
$C_1$、$C_2$ 不交。它是环
$$
B = k[x, y, 1/xy].
$$
% P12-ERR-0333
注意 $A_{xy} \cong B_{x + y}$（因为 $A$ 显然包含对任意多项式 $P$ 的元素
$xyP(x, y)$，以及元素 $xy/(x + y)$）。概形 $X$ 由仿射概形
$\Spec(A)$ 和 $\Spec(B)$ 通过同构 $A_{xy} \cong B_{x + y}$ 粘合而得，
因而显然在 $k$ 上有限型。为说明它是分离的，需要证明环映射
$A \otimes_k B \to B_{x + y}$ 满射。为此，使用
$A \otimes_k B$ 包含元素 $xy/(x + y) \otimes 1/xy$，它映到
$1/(x + y)$。态射 $Y \to X$ 由自然映射
$D(x + y) \to \Spec(A)$ 和 $D(xy) \to \Spec(B)$ 给出。
由于两者均有限，可得 $Y \to X$ 如所需是有限的。我们省略验证 $X$
确为上面所示图的余等化子；不过，见
% P12-ERR-0334
（在此插入关于概形范畴中推出的将来参考文献）。注意，态射
$\pi : Y \to X$ 确实把点 $(1, 0)$ 和 $(-1, 0)$ 映到 $X$ 中不同的点，
因为函数 $(x + y^3)/(x + y)^2 \in A$ 在这两点的值分别为
$1/1$ 与 $-1/(-1)^2 = -1$，它们总是不同
（除非特征为 $2$；在特征 $2$ 时请自行寻找两个点）。
我们把上述讨论总结为一个引理。

\begin{lemma}
\label{examples-lemma-quasi-affine-normalization-not-quasi-affine}
设 $k$ 为域。存在一个簇 $X$，它的正规化拟仿射，但它自身非拟仿射。
\end{lemma}

\begin{proof}
见上述讨论以及
% P12-ERR-0335
（在此插入关于正规化的将来参考文献）。
\end{proof}




\section{取概形论像}
\label{examples-section-scheme-theoretic-image}

\noindent
设 $k$ 为域，$t$ 为变量。令 $Y = \Spec(k[t])$ 且
$X = \coprod_{n \geq 1} \Spec(k[t]/(t^n))$。以 $f : X \to Y$ 表示如下态射：
对每个 $n \geq 1$ 使用闭浸入
$\Spec(k[t]/(t^n)) \to \Spec(k[t])$。在此情形下：
\begin{enumerate}
\item $f$ 的概形论像（《态射》定义
\ref{morphisms-definition-scheme-theoretic-image}）是 $Y$。
另一方面，$f$ 的像是 $Y$ 中由 $t = 0$ 给出的闭点。
因此，$f$ 的概形论像的底闭子集不等于 $f$ 的像的闭包。
\item 概形论像的构造不与到开子概形
$V = \Spec(k[t, 1/t]) \subset Y$ 上的限制可交换。
事实上，$V$ 在 $X$ 中的原像为空，故
$f|_{f^{-1}(V)} : f^{-1}(V) \to V$ 的概形论像是空概形。
它不等于 $Y \cap V$。
\end{enumerate}





\section{局部闭子集的像}
\label{examples-section-non-chevalley}

\noindent
Chevalley 定理断言，仿射概形之间有限表示态射下可构造集的像仍然可构造；
见《代数》定理 \ref{algebra-theorem-chevalley} 和《态射》第
\ref{morphisms-section-constructible} 节。我们将看到，对局部闭子集的像，
同样的结论并不成立。

\medskip\noindent
设 $k$ 为特征 $0$ 的域。考虑投影态射
$$
f :
X = \Spec(k[t, x_1, x_2, \ldots, y_1, y_2, \ldots])
\longrightarrow
\Spec(k[x_1, x_2, \ldots, y_1, y_2, \ldots]) = Y
$$
这是有限表示态射。令 $Z$ 为 $X$ 中由下列等式定义的闭子集：
$$
x_1(t - 1) = 0,\quad
x_2(t - 1)(t - 2) = 0,\quad
x_3(t - 1)(t - 2)(t - 3) = 0,\quad \ldots
$$
令 $U = \bigcup_{j \geq 1} U_j$ 为 $X$ 的开集，其中
$$
U_j = \text{使 }y_j(t - 1)(t - 2) ... (t - j)\text{ 非零的点}
$$
则有
$$
f(Z \cap U_j) =
\text{使 }x_1, \ldots, x_j\text{ 均为零而 }y_j\text{ 非零的点}
$$
我们断言，$B = f(Z \cap U) = \bigcup_{j \geq 1} f(Z \cap U_j)$
不是 $Y$ 中有限多个局部闭子集的并。

\medskip\noindent
证明该断言。假设 $B = A_1 \cup \cdots \cup A_m$ 是由 $Y$ 中局部闭
子集组成的 $B$ 的有限覆盖。我们将对 $n$ 归纳证明 $m \geq n$。
基始情形 $n = 1$ 成立，因为 $B$ 非空。假设 $n > 1$，且对 $n - 1$
归纳假设成立。由于 $B$ 的闭包是 $(x_1 = 0)$，某个 $A_i$ 必须包含
$(x_1 = 0)$ 的某个非空开子集。于是 $A_i$ 必须在 $(x_1 = 0)$ 中开。
但这样的开子集不可能包含满足 $y_1 = 0$ 的点；事实上，对 $B$ 中的点，
$y_1 = 0$ 迫使 $x_2 = 0$，这表明 $B$ 不包含 $(x_1 = 0)$ 内点
$(x, y)$ 的任何邻域。因此，余下 $m - 1$ 个集合在限制后构成
$B \cap (y_1 = 0)$ 的可构造覆盖。然而，注意右移映射
$x_i \mapsto x_{i + 1}$、$y_i \mapsto y_{i + 1}$ 将 $B$ 与
$B \cap (y_1 = 0)$ 识别！故由归纳假设，$m - 1 \geq n - 1$，
% P12-ERR-0336
从而 $m \geq n$。这完成归纳步骤的证明，也就证明了该断言。

\begin{lemma}
\label{examples-lemma-no-chevalley}
存在仿射概形之间的有限表示态射 $f : X \to Y$，以及 $X$ 的局部闭子集
$T$，使得 $f(T)$ 不是 $Y$ 中有限多个局部闭子集的并。
\end{lemma}

\begin{proof}
见上述讨论。
\end{proof}





\section{不在闭子概形中开的局部闭子概形}
\label{examples-section-strange-immersion}

\noindent
这是《态射》例 \ref{morphisms-example-thibaut} 的副本。
下面给出一个不能表示为先开浸入、后闭浸入之复合的浸入的例。
设 $k$ 为域。令 $X = \Spec(k[x_1, x_2, x_3, \ldots])$，并令
$U = \bigcup_{n = 1}^{\infty} D(x_n)$。于是 $U \to X$ 是开浸入。
考虑理想
$$
I_n =
(x_1^n, x_2^n, \ldots, x_{n - 1}^n, x_n - 1, x_{n + 1}, x_{n + 2}, \ldots)
\subset
k[x_1, x_2, x_3, \ldots][1/x_n].
$$
注意，对任意 $m \not = n$，有
$I_n k[x_1, x_2, x_3, \ldots][1/x_nx_m] = (1)$。
因此，$D(x_n)$ 上的拟凝聚理想 $\widetilde I_n$ 在 $D(x_nx_m)$ 上
彼此相同；即当 $n \not = m$ 时，
$\widetilde I_n|_{D(x_nx_m)} = \mathcal{O}_{D(x_n x_m)}$。
故这些理想粘合为拟凝聚理想层
$\mathcal{I} \subset \mathcal{O}_U$。令 $Z \subset U$ 为对应于
$\mathcal{I}$ 的闭子概形。于是 $Z \to X$ 是浸入。

\medskip\noindent
我们断言，不能把 $Z \to X$ 分解为
$Z \to \overline{Z} \to X$，其中 $\overline{Z} \to X$ 是闭浸入，
而 $Z \to \overline{Z}$ 是开浸入。事实上，$\overline{Z}$ 必须由理想
$I \subset k[x_1, x_2, x_3, \ldots]$ 定义，并满足
$I_n = I k[x_1, x_2, x_3, \ldots][1/x_n]$。
然而，在所有 $I_n$ 内都有像的元素
$f \in k[x_1, x_2, x_3, \ldots]$ 只有 $0$！所以这样的 $I$ 不存在。

\medskip\noindent
态射 $Z \to X$ 还给出了浸入的概形论像表现不良的例。
事实上，上述论证表明浸入 $Z \to X$ 的概形论像是 $X$。
另一方面，可见：
\begin{enumerate}
\item $Z$ 在 $X$ 中并非拓扑稠密；并且
\item $Z = Z \cap U \to U$ 的概形论像就是 $Z$。
这不等于 $U \cap X = U$，因此在此情形下，概形论像的构造不与到
开集上的限制可交换。
\end{enumerate}





\section{合适开集的不存在性}
\label{examples-section-nonexistence-opens}

\noindent
本节补充《性质》第 \ref{properties-section-finding-affine-opens} 节的结果。

\medskip\noindent
设 $k$ 为域，并令 $A = k[z_1, z_2, z_3, \ldots]/I$，其中 $I$ 是由
所有两两乘积 $z_iz_j$（$i \not = j$，$i, j \in \mathbf{N}$）生成的理想。
令 $S = \Spec(A)$。令 $s \in S$ 为对应于极大理想 $(z_i)$ 的闭点。
我们断言，不存在在 $S \setminus \{s\}$ 中稠密的拟紧开集
$V \subset S \setminus \{s\}$。注意
$S \setminus \{s\} = \bigcup D(z_i)$。每个 $D(z_i)$ 都是以
$\eta_i$ 为泛点的不可约开集。由此，对所有 $i$ 都有 $\eta_i \in V$。
然而，$S \setminus \{s\}$ 的主仿射开集形如 $D(f)$，其中
$f \in (z_1, z_2, \ldots)$。于是对某个 $n$ 有
$f \in (z_1, \ldots, z_n)$，从而可见 $D(f)$ 只包含有限多个点
$\eta_i$。所以 $V$ 不可能拟紧。

\medskip\noindent
设 $k$ 为域，并令 $B = k[x, z_1, z_2, z_3, \ldots]/J$，其中 $J$
由乘积 $xz_i$（$i \in \mathbf{N}$）以及所有两两乘积
$z_iz_j$（$i \not = j$，$i, j \in \mathbf{N}$）生成。
令 $T = \Spec(B)$。考虑主开集 $U = D(x)$。我们断言，不存在拟紧开集
$V \subset T$，使得 $V \cap U = \emptyset$ 且 $V \cup U$ 在 $T$ 中稠密。
令 $t \in T$ 为对应于极大理想 $(x, z_i)$ 的闭点。
$U$ 在 $T$ 中的闭包是 $\overline{U} = U \cup \{t\}$。
因此，$V \subset \bigcup_i D(z_i)$ 是拟紧开集。由上一段的论证，
可见 $V$ 不可能在 $\bigcup D(z_i)$ 中稠密。

\begin{lemma}
\label{examples-lemma-complement-of-affine-does-not-contain-qc-dense-open}
% P12-ERR-0337
仿射概形中拟紧开集的不存在性：
\begin{enumerate}
\item 存在仿射概形 $S$ 和仿射开集 $U \subset S$，使得不存在拟紧开集
$V \subset S$ 满足 $U \cap V = \emptyset$ 且 $U \cup V$ 在 $S$ 中稠密。
\item 存在仿射概形 $S$ 和闭点 $s \in S$，使得
$S \setminus \{s\}$ 不包含拟紧稠密开集。
\end{enumerate}
\end{lemma}

\begin{proof}
见上述讨论。
\end{proof}

\noindent
令 $X$ 为上述仿射概形 $T$ 的两个副本沿仿射开集 $U$ 的粘合。
于是存在态射 $\pi : X \to T$ 以及 $X = U_1 \cup U_2$，使得
$\pi$ 把 $U_i$ 同构地映到 $T$，并把 $U_1 \cap U_2$ 同构地映到 $U$。
注意，$X$ 拟分离（由《概形》引理
\ref{schemes-lemma-characterize-quasi-separated}）且拟紧。
我们断言，不存在分离、稠密的拟紧开集 $W \subset X$。
事实上，考虑映到上述闭点 $t \in T$ 的两个闭点
$x_1 \in U_1$、$x_2 \in U_2$。令
$\tilde \eta \in U_1 \cap U_2$ 为映到 $U$ 的（唯一）泛点 $\eta$ 的泛点。
注意，$\tilde\eta \leadsto x_1$ 且 $\tilde\eta \leadsto x_2$，
它们位于特化 $\eta \leadsto t$ 之上。由于
$\pi|_W : W \to T$ 是分离的，由分离性的赋值判据《概形》引理
\ref{schemes-lemma-valuative-criterion-separatedness}，不能同时有
$x_1$ 和 $x_2 \in W$。不妨设 $x_1 \not \in W$。
则 $W \cap U_1$ 是 $U_1$ 中不含 $x_1$ 的拟紧（因为 $X$ 拟分离）
稠密开集。现在注意，存在概形的同构
$(T, t) \cong (S, s)$（把 $x$ 映到 $z_1$，并把 $z_i$ 映到
$z_{i + 1}$）。因此，由本节第一段得到矛盾。

\begin{lemma}
\label{examples-lemma-no-dense-separated-quasi-compact-open-in-qcqs}
存在一个拟紧且拟分离的概形 $X$，它不包含分离的拟紧稠密开集。
\end{lemma}

\begin{proof}
见上述讨论。
\end{proof}





\section{拟紧稠密开子概形的不存在性}
\label{examples-section-nonexistence-qc-dense-open-subscheme}

\noindent
设 $X$ 为域 $k$ 上拟紧且拟分离的代数空间。我们知道概形轨迹
$X' \subset X$ 是稠密开子空间；见《空间的性质》命题
\ref{spaces-properties-proposition-locally-quasi-separated-open-dense-scheme}。
事实上，当 $X$ 合理时，这一结果也成立；见《良态空间》命题
\ref{decent-spaces-proposition-reasonable-open-dense-scheme}。
一个自然问题是，能否找到 $X$ 的拟紧稠密开子概形。一般而言，
这是不可能的。

\medskip\noindent
假设 $k$ 的特征不是 2。令
$B = k[x, z_1, z_2, z_3, \ldots]/J$，其中 $J$ 由乘积
$xz_i$（$i \in \mathbf{N}$）以及所有两两乘积
$z_iz_j$（$i \not = j$，$i, j \in \mathbf{N}$）生成。
令 $U = \Spec(B)$。以 $0 \in U$ 表示所有坐标均为零的闭点。令
$$
j : R = \Delta \amalg \Gamma \longrightarrow U \times_k U
$$
其中 $\Delta$ 是 $U$ 在 $k$ 上的对角态射的像，并且
$$
\Gamma = \{((x, 0, 0, 0, \ldots), (-x, 0, 0, 0, \ldots))
\mid x \in \mathbf{A}^1_k, x \not = 0\}.
$$
显然，$s, t : R \to U$ 是 étale 的，因而 $j$ 是 étale 等价关系。
商 $X = U/R$ 是代数空间（《空间》定理
\ref{spaces-theorem-presentation}）。注意，$j$ 不是浸入，因为
$(0, 0) \in \Delta$ 位于 $\Gamma$ 的闭包中。所以 $X$ 不是概形。
另一方面，由于 $R$ 拟紧，$X$ 拟分离。以 $0_X$ 表示点 $0 \in U$
的像。我们断言 $X \setminus \{0_X\}$ 是概形，具体为
$$
X \setminus \{0_X\} =
\Spec\left(k[x^2, x^{-2}]\right) \amalg
\Spec\left(k[z_1, z_2, z_3, \ldots]/(z_iz_j)\right) \setminus \{0\}
$$
（细节从略）。另一方面，我们已在第
\ref{examples-section-nonexistence-opens} 节看到，右侧的概形不包含拟紧稠密开集。

\begin{lemma}
\label{examples-lemma-nonexistence-qc-dense-open-subscheme}
存在一个拟紧且拟分离的代数空间，它不包含拟紧稠密开子概形。
\end{lemma}

\begin{proof}
见上述讨论。
\end{proof}

\noindent
使用《空间》例 \ref{spaces-example-non-representable-descent} 的构造，
其方式与我们上面使用《空间》例
\ref{spaces-example-affine-line-involution} 构造的方式相同，可以得到一个拟紧、
拟分离且局部分离的代数空间，它不包含拟紧稠密开子概形。





\section{代数空间上的仿射概形}
\label{examples-section-embedding-affines}

\medskip\noindent
假设 $f : Y \to X$ 是概形之间的态射，$f$ 局部有限型，且 $Y$ 仿射。
则存在从 $Y$ 到 $X$ 上仿射 $n$-空间的浸入
$Y \to \mathbf{A}^n_X$。见稍为一般的《态射》引理
\ref{morphisms-lemma-quasi-affine-finite-type-over-S}。

\medskip\noindent
现在假设 $f : Y \to X$ 是代数空间之间的态射，$f$ 局部有限型，
而 $Y$ 是仿射概形。此时一般不能找到从 $Y$ 到 $X$ 上仿射
$n$-空间的浸入。

\medskip\noindent
第一个（不太友好的）反例是 $Y = \Spec(k)$ 和
$X = [\mathbf{A}^1_k/\mathbf{Z}]$，其中 $k$ 是特征零的域，
$\mathbf{Z}$ 通过平移 $(n, t) \mapsto t + n$ 作用在
$\mathbf{A}^1_k$ 上。事实上，对任意 $X$ 上的态射
$Y \to \mathbf{A}^n_X$，可以拉回到 $X$ 的覆盖 $\mathbf{A}^1_k$，
于是得到映入 $\mathbf{A}^{n + 1}_k$ 的无穷多个
% P12-ERR-0338
$\mathbf{A}^1_k$ 的不交并，而该映射不是浸入。

\medskip\noindent
第二个反例是 $Y = \mathbf{A}^1_k \to X = \mathbf{A}^1_k/R$，其中
$R = \{(t, t)\} \amalg \{(t, -t), t \not = 0\}$。
事实上，在此情形下，态射 $Y \to \mathbf{A}^n_X$ 将由 $Y$ 上的某些
正则函数 $f_1, \ldots, f_n$ 给出，因而 $Y$ 与覆盖
$\mathbf{A}^{n + 1}_k \to \mathbf{A}^n_X$ 的纤维积是概形
$$
\{(f_1(t), \ldots, f_n(t), t)\} \amalg
\{(f_1(t), \ldots, f_n(t), -t), t \not = 0\}
$$
它到 $\mathbf{A}^{n + 1}_k$ 的显然态射不是浸入。
注意，这给出了 $X$ 拟分离的反例。

\begin{lemma}
\label{examples-lemma-cannot-embed-into-affine}
存在代数空间之间的有限型态射 $Y \to X$，其中 $Y$ 仿射、$X$ 拟分离，
并且不存在 $X$ 上的浸入 $Y \to \mathbf{A}^n_X$。
\end{lemma}

\begin{proof}
见上述讨论。
\end{proof}





\section{拟凝聚模的推前}
\label{examples-section-push-quasi-coherent}

\noindent
我们在《概形》引理 \ref{schemes-lemma-push-forward-quasi-coherent} 中证明：
当 $f$ 拟紧且拟分离时，$f_*$ 把拟凝聚模变为拟凝聚模。
下面给出一些例，说明这两个条件都是必要的。

\medskip\noindent
假设 $Y = \Spec(A)$ 是仿射概形，并且
$X = \coprod_{n \in \mathbf{N}} Y$。令 $f : X \to Y$ 为显然的态射。
我们断言，$f_*\mathcal{O}_X$ 不是拟凝聚的。事实上，对 $a \in A$ 有
$$
f_*\mathcal{O}_X(D(a)) = \prod\nolimits_{n \in \mathbf{N}} A_a
$$
因此，要使 $f_*\mathcal{O}_X$ 拟凝聚，就必须对所有 $a \in A$ 有
$$
\prod\nolimits_{n \in \mathbf{N}} A_a
=
\left(\prod\nolimits_{n \in \mathbf{N}} A\right)_a
$$
这一般不成立。例如，若 $A = \mathbf{Z}$ 且 $a = 2$，则
$(1, 1/2, 1/4, 1/8, \ldots)$ 是左侧的元素，但不属于右侧。
注意，$f$ 是非拟紧的分离态射。

\medskip\noindent
设 $k$ 为域。令
$$
A = k[t, z, x_1, x_2, x_3, \ldots]/(tx_1z, t^2x_2^2z, t^3x_3^3z, \ldots)
$$
令 $Y = \Spec(A)$。令 $V \subset Y$ 为开子概形
$V = D(x_1) \cup D(x_2) \cup \ldots$。令 $X$ 为 $Y$ 的两个副本沿
$V$ 的粘合，并令 $f : X \to Y$ 为显然的态射。于是有正合序列
$$
0 \to f_*\mathcal{O}_X \to
\mathcal{O}_Y \oplus \mathcal{O}_Y \xrightarrow{(1, -1)} j_*\mathcal{O}_V
$$
其中 $j : V \to Y$ 是包含态射。由于
$$
A \longrightarrow \prod A_{x_n}
$$
是单射（细节从略），可见 $\Gamma(Y, f_*\mathcal{O}_X) = A$。
另一方面，映射
$$
A_t \longrightarrow \prod A_{tx_n}
$$
的核非零，因为它包含元素 $z$。所以
$\Gamma(D(t), f_*\mathcal{O}_X)$ 严格大于 $A_t$，因为它包含
$(z, 0)$。因此，$f_*\mathcal{O}_X$ 不是拟凝聚的。
注意，$f$ 拟紧但非拟分离。

\begin{lemma}
\label{examples-lemma-pushforward-quasi-coherent}
《概形》引理 \ref{schemes-lemma-push-forward-quasi-coherent} 的条件是最优的：
拟紧性假设与拟分离性假设都不能去掉。
\end{lemma}

\begin{proof}
见上述讨论。
\end{proof}





\section{非有限生成而茎为秩 1 有限自由的模}
\label{examples-section-nonfree}

\noindent
令 $R = \mathbf{Q}[x]$。在 $R$ 的分式域中取子模
$M = \sum_{n \in \mathbf{N}} \frac{1}{x - n}R$。
则 $M$ 不是有限生成的，但对 $R$ 的每个素理想 $\mathfrak p$，
作为 $R_{\mathfrak p}$-模有
$M_{\mathfrak p} \cong R_{\mathfrak p}$。

\medskip\noindent
一个风格类似的例是 $R = \mathbf{Z}$ 和
$M = \sum_{p \text{ 为素数}} \frac{1}{p} \mathbf{Z} \subset \mathbf{Q}$；
它等于分母 $b$ 非零且无平方因子的分数 $\frac{a}{b}$ 所成的集合。








\section{在茎中可逆但本身不可逆的理想}
\label{examples-section-locally-invertible-not-invertible}

\noindent
设 $A$ 为整环，$I \subset A$ 为非零理想。回顾一下，我们说 $I$ 可逆，
是指 $I$ 作为 $A$-模可逆。我们将构造这种情形的一个例：$I$ 不可逆，
但对每个素理想 $\mathfrak q \subset A$，
$I_\mathfrak q = (f) \subset A_\mathfrak q$ 都是（非零）主理想。
文献中有时把对所有素理想 $\mathfrak q$，$I_\mathfrak q$ 都是主理想的
性质表述为“$I$ 是局部主理想”。我们不能采用这一术语，因为我们的
“局部”总是指“在 Zariski 拓扑中局部”（或在当前所用的任何拓扑中局部）。

\medskip\noindent
令 $R = \mathbf{Q}[x]$，并令 $M = \sum \frac{1}{x - n}R$ 为第
\ref{examples-section-nonfree} 节构造的模。考虑环\footnote{环 $A$ 是一个局部环
均为 Noether 环的非 Noether 整环的例。}
$$
A = \text{Sym}^*_R(M)
$$
以及理想 $I = M A = \bigoplus_{d \geq 1} \text{Sym}^d_R(M)$。
由于 $M$ 作为 $R$-模不是有限生成的，$I$ 不可能作为 $A$ 中的理想由
有限多个元素生成。由于可逆模是有限生成的，这意味着 $I$ 不可逆。
另一方面，令 $\mathfrak p \subset R$ 为素理想。按构造，
$M_\mathfrak p \cong R_\mathfrak p$。因此
$$
A_\mathfrak p = \text{Sym}^*_{R_\mathfrak p}(M_\mathfrak p) \cong
\text{Sym}^*_{R_\mathfrak p}(R_\mathfrak p) =
R_\mathfrak p[T]
$$
作为分次 $R_\mathfrak p$-代数成立。由此，
$I_\mathfrak p \subset A_\mathfrak p$ 由非零因子 $T$ 生成。
所以，对任意素理想 $\mathfrak q \subset A$，$I_\mathfrak q$ 确由单个元素生成。

\begin{lemma}
\label{examples-lemma-locally-principal-not-invertible}
存在整环 $A$ 和非零理想 $I \subset A$，使得对所有素理想
$\mathfrak q \subset A$，$I_\mathfrak q \subset A_\mathfrak q$ 都是主理想，
但 $I$ 不是可逆 $A$-模。
\end{lemma}

\begin{proof}
见上述讨论。
\end{proof}





\section{非投射的有限平坦模}
\label{examples-section-finite-flat-not-projective}

\noindent
这是《代数》注 \ref{algebra-remark-warning} 的副本。
有限 $R$-模若在 $R$ 上平坦，并不自动投射。一个反例是：
$R = \mathcal{C}^\infty(\mathbf{R})$ 为 $\mathbf{R}$ 上无穷次可微函数环，
并且 $M = R_{\mathfrak m} = R/I$，其中
$\mathfrak m = \{f \in R \mid f(0) = 0\}$ 且
$I = \{f \in R \mid \exists \epsilon, \epsilon > 0 :
f(x) = 0\ \forall x, |x| < \epsilon\}$。

\medskip\noindent
态射 $\Spec(R/I) \to \Spec(R)$ 也是平坦但非开的闭浸入的一个例。

\begin{lemma}
\label{examples-lemma-finite-flat-non-projective}
特殊的平坦模：
\begin{enumerate}
\item 存在环 $R$ 和有限平坦 $R$-模 $M$，而该模不投射。
\item 存在平坦但非开的闭浸入。
\end{enumerate}
\end{lemma}

\begin{proof}
见上述讨论。
\end{proof}





\section{非局部自由的投射模}
\label{examples-section-projective-not-locally-free}

\noindent
我们给出两个例。一个例的秩介于 $0$ 与 $1$ 之间，另一个例的秩为
$\aleph_0$。

\begin{lemma}
\label{examples-lemma-ideal-generated-by-idempotents-projective}
设 $R$ 为环。设 $I \subset R$ 为由可数个幂等元生成的理想。
则 $I$ 作为 $R$-模是投射的。
\end{lemma}

\begin{proof}
设 $I = (e_1, e_2, e_3, \ldots)$，其中 $e_n$ 是 $R$ 的幂等元。
归纳地把 $e_{n + 1}$ 替换为 $e_n + (1 - e_n)e_{n + 1}$ 后，
可设 $(e_1) \subset (e_2) \subset (e_3) \subset \ldots$，从而
$I = \bigcup_{n \geq 1} (e_n) = \colim_n e_nR$。在此情形下，
$$
\Hom_R(I, M) = \Hom_R(\colim_n e_nR, M)
= \lim_n \Hom_R(e_nR, M) = \lim_n e_nM
$$
注意，转移映射 $e_{n + 1}M \to e_nM$ 由乘以 $e_n$ 给出，并且满射。
所以，由《代数》引理 \ref{algebra-lemma-ML-exact-sequence}，函子
$\Hom_R(I, M)$ 正合；即 $I$ 是投射 $R$-模。
\end{proof}

\noindent
假设 $P \subset Q$ 是 $R$-模的包含，$Q$ 是有限 $R$-模，而 $P$
局部自由；见《代数》定义 \ref{algebra-definition-locally-free}。
假设 $Q$ 作为 $R$-模可由 $N$ 个元素生成。由《代数》引理
\ref{algebra-lemma-map-cannot-be-injective} 可知，$P$ 有限局部自由
（各自由部分的秩至多为 $N$）。在这种情形下，$P$ 是有限 $R$-模；
见《代数》引理 \ref{algebra-lemma-finite-projective}。

\medskip\noindent
结合上述结果，可见由可数个幂等元生成的非有限生成理想是投射的，
但不是局部自由的。一个显式例是
$R = \prod_{n \in \mathbf{N}} \mathbf{F}_2$，而 $I$ 是由幂等元
$$
e_n = (1, 1, \ldots, 1, 0, \ldots )
$$
生成的理想，其中 $1$ 所成序列的长度为 $n$。

\begin{lemma}
\label{examples-lemma-ideal-projective-not-locally-free}
存在环 $R$ 和理想 $I$，使得 $I$ 作为 $R$-模投射，但作为
$R$-模不是局部自由的。
\end{lemma}

\begin{proof}
见上文。
\end{proof}

\begin{remark}
\label{examples-remark-projective-with-dual-finite-not-finite}
上面 $R = \prod_{n \in \mathbf{N}} \mathbf{F}_2$、$I$ 是由幂等元
$e_n = (1, 1, \ldots, 1, 0, \ldots )$ 生成的理想这一例，还给出了
$M = I$ 的一个例：它是非有限生成投射模，而其对偶有限生成，甚至有限自由。
事实上，使用引理 \ref{examples-lemma-ideal-generated-by-idempotents-projective}
证明中的策略，读者可证明 $\Hom_R(I, R) \cong R$。
\end{remark}

\begin{lemma}
\label{examples-lemma-chow-group-product}
设 $K$ 为域。设 $C_i$（$i = 1, \ldots, n$）为 $K$ 上光滑、投射、
几何不可约的曲线。设 $P_i \in C_i(K)$ 为有理点，并设
$Q_i \in C_i$ 为满足 $[\kappa(Q_i) : K] = 2$ 的点。
则 $[P_1 \times \ldots \times P_n]$ 在
$\CH_0(U_1 \times_K \ldots \times_K U_n)$ 中非零，其中
$U_i = C_i \setminus \{Q_i\}$。
\end{lemma}

\begin{proof}
存在次数映射
$\deg : \CH_0(C_1 \times_K \ldots \times_K C_n) \to \mathbf{Z}$。
由于每个 $Q_i$ 在 $K$ 上的次数为 $2$，可见支撑在“边界”
$$
C_1 \times_K \ldots \times_K C_n
\setminus
U_1 \times_K \ldots \times_K U_n
$$
上的任意零维循环的次数都能被 $2$ 整除。
\end{proof}

\noindent
利用上述引理，还可如下构造另一个投射但非局部自由的模的例。
设 $C_n$（$n = 1, 2, 3, \ldots$）为 $\mathbf{Q}$ 上光滑、投射、
几何不可约的曲线，每条曲线上都有一对点 $P_n, Q_n \in C_n$，满足
$\kappa(P_n) = \mathbf{Q}$，且 $\kappa(Q_n)$ 是 $\mathbf{Q}$ 的二次扩张。
令 $U_n = C_n \setminus \{Q_n\}$；这是仿射曲线。
令 $\mathcal{L}_n$ 为 $U_n$ 上点 $P_n$ 的理想层之逆。
注意，在零维循环群 $\CH_0(U_n)$ 中，
$c_1(\mathcal{L}_n) = [P_n]$。令
$A_n = \Gamma(U_n, \mathcal{O}_{U_n})$，并令
$L_n = \Gamma(U_n, \mathcal{L}_n)$；这是 $A_n$ 上秩为 $1$ 的局部自由模。
令
$$
B_n = A_1 \otimes_{\mathbf{Q}} A_2 \otimes_{\mathbf{Q}} \ldots
\otimes_{\mathbf{Q}} A_n
$$
于是 $\Spec(B_n) = U_1 \times \ldots \times U_n$，所有积都在
$\Spec(\mathbf{Q})$ 上取。对 $i \leq n$，令
$$
L_{n, i} =
A_1 \otimes_{\mathbf{Q}} \ldots \otimes_{\mathbf{Q}} L_i
\otimes_{\mathbf{Q}} \ldots \otimes_{\mathbf{Q}} A_n
$$
它是秩为 $1$ 的局部自由 $B_n$-模。注意，
% P12-ERR-0339
这也是 $\text{pr}_i^*\mathcal{L}_n$ 的整体截面。令
$$
B_\infty = \colim_n B_n
\quad\text{且}\quad
L_{\infty, i} = \colim_n L_{n, i}
$$
最后，令
$$
M = \bigoplus\nolimits_{i \geq 1} L_{\infty, i}.
$$
这是有限局部自由模的直和，因而投射。我们断言 $M$ 不是局部自由的。
事实上，假设 $f \in B_\infty$ 是非零函数，使得 $M_f$ 在
$(B_\infty)_f$ 上自由。设 $e_1, e_2, \ldots$ 为一组基。
选取 $n \geq 1$，使得 $f \in B_n$。再选取 $m \geq n + 1$，
使得 $e_1, \ldots, e_{n + 1}$ 属于
$$
\bigoplus\nolimits_{1 \leq i \leq m} L_{m, i}.
$$
由于元素 $e_1, \ldots, e_{n + 1}$ 在忠实平坦基变换后成为一组基的一部分，
可知 Chern 类
$$
c_i(\text{pr}_1^*\mathcal{L}_1 \oplus \ldots \oplus
\text{pr}_m^*\mathcal{L}_m), \quad i = m, m - 1, \ldots, m - n
$$
在
$$
D(f) \subset U_1 \times \ldots \times U_m
$$
的 Chow 群中为零。由于 $f$ 是 $U_1 \times \ldots \times U_n$ 上一个
函数的拉回，这特别蕴含
$$
c_{m - n}(\mathcal{O}_W^{\oplus n} \oplus
\text{pr}_1^*\mathcal{L}_{n + 1} \oplus \ldots
\oplus \text{pr}_{m - n}^*\mathcal{L}_m) = 0.
$$
在域 $K = \mathbf{Q}(C_1 \times \ldots \times C_n)$ 上的簇
$$
W = (C_{n + 1} \times \ldots \times C_m)_K
$$
上成立。换言之，循环
$$
[(P_{n + 1} \times \ldots \times P_m)_K]
$$
在 $W$ 上的零维循环 Chow 群中为零。这与上面的引理
\ref{examples-lemma-chow-group-product} 矛盾，因为点 $Q_i$（
$n + 1 \leq i \leq m$）在
% P12-ERR-0340
$(C_n)_K$ 上诱导相应的点 $Q_i'$，并且由于
% P12-ERR-0341
$K/\mathbf{Q}$ 几何不可约，有 $[\kappa(Q_i') : K] = 2$。

\begin{lemma}
\label{examples-lemma-projective-not-locally-free}
存在可数环 $R$ 和投射模 $M$，它是可数多个秩为 $1$ 的局部自由模的直和，
但 $M$ 不是局部自由的。
\end{lemma}

\begin{proof}
见上文。
\end{proof}





\section{具有非零平坦理想的零维局部环}
\label{examples-section-zero-dimensional-flat-ideal}

\noindent
文献 \cite{Lazard} 和 \cite{Autour} 给出了具有非零平坦理想的零维
局部环的例。构造如下。设 $k$ 为域。设 $X_i, Y_i$（$i \geq 1$）为变量。
取 $R = k[X_i, Y_i]/(X_i - Y_i X_{i + 1}, Y_i^2)$。
分别以 $x_i$ 和 $y_i$ 表示 $X_i$ 和 $Y_i$ 在此环中的像。注意，在此环中
$$
x_i = y_i x_{i + 1} = y_i y_{i +1} x_{i + 2} =
y_i y_{i + 1} y_{i + 2} x_{i + 3} = \ldots
$$
环 $R$ 只有一个素理想，即 $\mathfrak m = (x_i, y_i)$。
我们断言，理想 $I = (x_i)$ 作为 $R$-模是平坦的。

\medskip\noindent
注意，$x_i$ 在 $R$ 中的零化理想为
$(x_1, x_2, x_3, \ldots, y_i, y_{i + 1}, y_{i + 2}, \ldots)$。
考虑由元素 $e_i$（$i \geq 1$）和关系 $e_i = y_i e_{i + 1}$ 生成的
$R$-模 $M$。$M$ 是平坦的，因为它是 $R$ 的诸副本经转移映射
$$
R \xrightarrow{y_1} R \xrightarrow{y_2} R \xrightarrow{y_3} \ldots
$$
所得的余极限 $\colim_i R$。注意，$e_i$ 在 $M$ 中的零化理想为
$(x_1, x_2, x_3, \ldots, y_i, y_{i + 1}, y_{i + 2}, \ldots)$。
由于 $M$ 与 $I$ 的每个元素都可分别写成 $f e_i$ 与 $h x_i$，
其中某些 $f, h \in R$，可见映射 $M \to I$、$e_i \to x_i$ 是同构，
且 $I$ 平坦。

\begin{lemma}
\label{examples-lemma-zero-dimensional-flat-ideal}
\begin{slogan}
具有平坦理想的零维环。
\end{slogan}
存在仅有一个素理想的局部环 $R$，以及非零理想 $I \subset R$，
它是平坦 $R$-模。
\end{lemma}

\begin{proof}
见上述讨论。
\end{proof}





\section{非满射的零维环上态射}
\label{examples-section-epimorphism-not-surjective}

\noindent
文献 \cite{Lazard-deux} 和 \cite{Autour} 中可找到如下例。
设 $k$ 为域。考虑环同态
$$
k[x_1, x_2, \ldots, z_1, z_2, \ldots]/
(x_i^{4^i}, z_i^{4^i})
\longrightarrow
k[x_1, x_2, \ldots, y_1, y_2, \ldots]/
(x_i^{4^i}, y_i - x_{i + 1}y_{i + 1}^2)
$$
它把 $x_i$ 映到 $x_i$，把 $z_i$ 映到 $x_iy_i$。注意，
$y_i^{4^{i + 1}}$ 在右侧为零，而 $y_1$ 非零（细节从略）。
该映射不满射：令 $\deg(x_i) = -1$、$\deg(z_i) = 0$、
$\deg(y_i) = 1$，便可把上述映射看成 $\mathbf{Z}$-分次代数的映射；
此时显然 $y_1$ 不在像中。最后，该映射是上态射，因为
$$
y_{i - 1} \otimes 1 = x_i y_i^2 \otimes 1 = y_i \otimes x_i y_i =
x_i y_i \otimes y_i = 1 \otimes x_i y_i^2.
$$
% P12-ERR-0342
所以，目标在源上的张量积同构于目标。

\begin{lemma}
\label{examples-lemma-epi-not-surjective}
存在维数为 $0$ 的局部环之间的上态射，而它不是满射。
\end{lemma}

\begin{proof}
见上述讨论。
\end{proof}





\section{有限型、非有限表示、在素理想处平坦}
\label{examples-section-ft-not-fp-flat-at-prime}

\noindent
设 $k$ 为域。考虑局部环 $A_0 = k[x, y]_{(x, y)}$。
记 $\mathfrak p_{0, n} = (y + x^n + x^{2n + 1})$。这是素理想。令
$$
A = A_0[z_1, z_2, z_3, \ldots]/(z_n z_m, z_n(y + x^n + x^{2n + 1}))
$$
注意，$A \to A_0$ 是满射，其核是平方零理想。因此，$A$ 也是局部环，
而 $A$ 的素理想与 $A_0$ 的素理想一一对应。以 $\mathfrak p_n$ 表示
$A$ 中对应于 $\mathfrak p_{0, n}$ 的素理想。注意，$\mathfrak p_n$ 是
$z_n$ 在 $A$ 中的零化理想。令
$$
C = A[z]/(xz^2 + z + y)[\frac{1}{2zx + 1}]
$$
注意，$A \to C$ 是 étale 环映射；见《代数》例
\ref{algebra-example-make-standard-smooth}。令 $\mathfrak q \subset C$
为由 $x$、$y$、$z$ 以及所有 $z_n$ 生成的极大理想。由于
$A \to C$ 平坦，$z_n$ 在 $C$ 中的零化理想为 $\mathfrak p_nC$。
计算得
\begin{align*}
C/\mathfrak p_n C
& =
A_0[z]/(xz^2 + z + y, y + x^n + x^{2n + 1})[1/(2zx + 1)] \\
& =
k[x]_{(x)}[z]/(xz^2 + z - x^n - x^{2n + 1})[1/(2zx + 1)] \\
& =
k[x]_{(x)}[z]/(z - x^n) \times
k[x]_{(x)}[z]/(xz + x^{n + 1} + 1)[1/(2zx + 1)] \\
& =
k[x]_{(x)} \times k(x)
\end{align*}
因为 $(z - x^n)(xz + x^{n + 1} + 1) = xz^2 + z - x^n - x^{2n + 1}$。
因此，$\mathfrak p_nC = \mathfrak r_n \cap \mathfrak q_n$，其中
$\mathfrak r_n = \mathfrak p_nC + (z - x^n)C$ 且
$\mathfrak q_n = \mathfrak p_nC + (xz + x^{n + 1} + 1)C$。
由于 $\mathfrak q_n + \mathfrak r_n = C$，还得到
$\mathfrak p_nC = \mathfrak r_n \mathfrak q_n$。
由此，$\mathfrak q_n$ 是 $\xi_n = (z - x^n)z_n$ 的零化理想。
注意，一方面 $\mathfrak r_n \subset \mathfrak q$，另一方面
$\mathfrak q_n + \mathfrak q = C$。例如，这可由 $\mathfrak q_n$ 是
$C$ 中不同于 $\mathfrak q$ 的极大理想推出。类似地，当
$n \not = m$ 时，有 $\mathfrak q_n + \mathfrak q_m = C$。
现在令
$$
B = \Im(C \longrightarrow C_{\mathfrak q})
$$
注意，由于 $xz + x^{n + 1} + 1$ 不属于 $\mathfrak q$，元素 $\xi_n$
在 $B$ 中映为零。以 $\mathfrak q' \subset B$ 表示 $\mathfrak q$ 的像。
按构造，$B$ 是有限型 $A$-代数，且
$B_{\mathfrak q'} \cong C_{\mathfrak q}$。特别地，
$B_{\mathfrak q'}$ 在 $A$ 上平坦。

\medskip\noindent
我们断言，不存在 $g' \in B$、$g' \not \in \mathfrak q'$，使得
$B_{g'}$ 在 $A$ 上有限表示。我们简述该断言的证明。
选取映到 $g' \in B$ 的元素 $g \in C$。考虑映射
$C_g \to B_{g'}$。由《代数》引理
\ref{algebra-lemma-finite-presentation-independent}，
% P12-ERR-0343
$B_g$ 在 $A$ 上有限表示，当且仅当 $C_g \to B_{g'}$ 的核有限生成。
但是，元素 $g \in C$ 不属于 $\mathfrak q$，故在
$A_0[z]/(xz^2 + z + y)$ 中映为非零元素。因此，$g$ 只能属于有限多个
素理想 $\mathfrak q_n$，因为素理想
$(y + x^n + x^{2n + 1}, xz + x^{n + 1} + 1)$ 构成二维不可约 Noether
空间 $\Spec(k[x, y, z]/(xz^2 + z + y))$ 的无穷多个余维 1 的点。
映射
$$
\bigoplus\nolimits_{g \not \in \mathfrak q_n} C/\mathfrak q_n
\longrightarrow
C_g, \quad
(c_n) \longrightarrow \sum c_n \xi_n
$$
是单射，其像是 $C_g \to B_{g'}$ 的核。我们省略这一陈述的证明。
（提示：把 $A$ 作为 $A_0$-模写成 $A = A_0 \oplus I$，其中 $I$ 是
$A \to A_0$ 的核。类似地，写成 $C = C_0 \oplus IC$。再写成
$IC = \bigoplus Cz_n \cong \bigoplus (C/\mathfrak r_n \oplus C/\mathfrak q_n)$，
并研究乘以 $g$ 对各直和项的作用。）
这完成了该断言证明的略述。这还证明，对上述任意 $g'$，$B_{g'}$
在 $A$ 上都不平坦。事实上，若它平坦，则 $z_n$ 在 $B_{g'}$ 中之像的
零化理想将是 $\mathfrak p_nB_{g'}$，并且不会包含 $z - x^n$。

\medskip\noindent
由此，我们可以对 \cite[第一部分，注 (3.4.7) (\romannumeral 5)]{GruRay}
中提出的问题给出否定回答。精确陈述如下。

\begin{lemma}
\label{examples-lemma-example-raynaud-gruson}
存在局部环 $A$、有限型环映射 $A \to B$，以及位于
$\mathfrak m_A$ 上方的素理想 $\mathfrak q$，使得
$B_{\mathfrak q}$ 在 $A$ 上平坦；并且对任意
$g \in B$、$g \not \in \mathfrak q$，环 $B_g$ 在 $A$ 上既非有限表示，
也不平坦。
\end{lemma}

\begin{proof}
见上述讨论。
\end{proof}





\section{有限型、平坦但非有限表示}
\label{examples-section-finite-type-flat-not-finite-presentation}

\noindent
本节给出一些有限型且平坦、但非有限表示的环映射和态射的例。

\medskip\noindent
设 $R$ 为具有理想 $I$ 的环，使得 $R/I$ 是有限平坦但非投射的模；
显式例见第 \ref{examples-section-finite-flat-not-projective} 节。
注意，这意味着 $I$ 不是有限生成的；见《代数》引理
\ref{algebra-lemma-finitely-generated-pure-ideal}。还要注意，
$I = I^2$；见《代数》引理 \ref{algebra-lemma-pure}。
各例的基环将取为 $R$，相应的基概形为 $S = \Spec(R)$。

\medskip\noindent
考虑环映射 $R \to R \oplus R/I\epsilon$，其中约定
$\epsilon^2 = 0$。这是有限平坦但非有限表示的环映射。
所有纤维环都是完全交且几何不可约。

\medskip\noindent
令 $A = R[x, y]/(xy, ay; a \in I)$。注意，作为 $R$-模，有
% P12-ERR-0344
$A = \bigoplus_{i \geq 0} Ry^i \oplus \bigoplus_{j > 0} R/Ix^j$。
所以 $R \to A$ 是平坦有限型但非有限表示的环映射。
每个纤维环同构于 $\kappa(\mathfrak p)[x, y]/(xy)$ 或
$\kappa(\mathfrak p)[x]$。

\medskip\noindent
可通过取 $B = R[X_0, X_1, X_2]/(X_1X_2, aX_2; a \in I)$，
把上一个例变成投射态射。在此情形下，
$X = \text{Proj}(B) \to S$ 是非有限表示的固有平坦态射，并且对每个
$s \in S$，纤维 $X_s$ 同构于 $\mathbf{P}^1_s$，或同构于
$\mathbf{P}^2_s$ 中由 $X_1X_2$ 的零化定义的闭子概形
（这是一条算术亏格为 $0$ 的投射结点曲线）。

\medskip\noindent
令 $M = R \oplus R \oplus R/I$。令 $B = \text{Sym}_R(M)$ 为 $M$
上的对称代数，并令 $X = \text{Proj}(B)$。于是
$X \to S$ 是非有限表示的固有平坦态射，并且对 $s \in S$，几何纤维
同构于 $\mathbf{P}^1_s$ 或 $\mathbf{P}^2_s$。特别地，这些纤维光滑且
几何不可约。

\begin{lemma}
\label{examples-lemma-finite-type-flat-not-finitely-presented}
存在下列各类例：
\begin{enumerate}
\item 非有限表示的平坦有限型环映射，其完全交纤维环几何不可约；
\item 非有限表示的平坦有限型环映射，其纤维环几何连通、几何约化、
维数为 1 且为完全交；
\item 非有限表示的概形固有平坦态射 $X \to S$，其每个纤维同构于
$\mathbf{P}^1_s$，或同构于 $\mathbf{P}^2_s$ 中 $X_1X_2$ 的零化轨迹；以及
\item 非有限表示的概形固有平坦态射 $X \to S$，其每个纤维同构于
$\mathbf{P}^1_s$ 或 $\mathbf{P}^2_s$。
\end{enumerate}
\end{lemma}

\begin{proof}
见上述讨论。
\end{proof}





\section{有限型环映射的拓扑}
\label{examples-section-topology-finite-type}

\noindent
设 $A \to B$ 为局部整环之间的局部映射。若 $A$ 为 Noether 环、
$A \to B$ 本质有限型，且 $A/\mathfrak m_A \subset B/\mathfrak m_B$
有限，则存在素理想 $\mathfrak q \subset B$、
$\mathfrak q \not = \mathfrak m_B$，使得 $A \to B/\mathfrak q$ 是某个
拟有限环映射的局部化。见 More on Morphisms，引理
\ref{more-morphisms-lemma-quasi-finite-quasi-section-meeting-nearby-open}。

\medskip\noindent
% P12-ERR-0345
本节给出一个例，说明当 $A$ 不再是 Noether 环时，上述结果不成立。
具体地，设 $k$ 为域，并令
$$
A = \{a_0 + a_1 x + a_2 x^2 + \ldots \mid a_0 \in k, a_i \in k((y))
\text{，对 }i\geq 1\}
$$
以及
$$
C = \{a_0 + a_1 x + a_2 x^2 + \ldots \mid a_0 \in k[y], a_i \in k((y))
\text{，对 }i\geq 1\}.
$$
包含 $A \to C$ 有限型，因为 $C$ 由 $y$ 在 $A$ 上生成。
我们断言，$A$ 是以
$\mathfrak m = \{a_1 x + a_2 x^2 + \ldots \in A\}$ 为极大理想的局部环，
且除 $(0)$ 和 $\mathfrak m$ 外没有其他素理想。事实上，$A$ 的元素
$f = a_0 + a_1 x + a_2 x^2 + \ldots$ 只要满足 $a_0 \not = 0$ 就可逆。
若 $\mathfrak q \subset A$ 是非零素理想，而
$f = a_i x^i + \ldots \in \mathfrak q$，则利用幂级数的性质可见，
对任意 $g \in k((y))$，元素
$g^{i + 1} x^{i + 1} \in \mathfrak q$；即 $gx \in \mathfrak q$。
这证明 $\mathfrak q = \mathfrak m$。

\medskip\noindent
至于环 $C$ 的谱，用与上面相同的论证可见，任意非零素理想都包含
位于 $\mathfrak m$ 上方的素理想
$\mathfrak p = \{a_1 x + a_2 x^2 + \ldots \in C\}$。
因此，$C$ 中位于 $(0)$ 上方的素理想只有 $(0)$。令
$\mathfrak m_C = yC + \mathfrak p$ 且 $B = C_{\mathfrak m_C}$。
则 $A \to B$ 就是所需的例。

\begin{lemma}
\label{examples-lemma-topology-finite-type}
存在局部整环之间的局部同态 $A \to B$，它本质有限型，且
$A/\mathfrak m_A \to B/\mathfrak m_B$ 有限，并且对 $B$ 的每个素理想
$\mathfrak q \not = \mathfrak m_B$，环映射 $A \to B/\mathfrak q$
都不是某个拟有限环映射的局部化。
\end{lemma}

\begin{proof}
见上述讨论。
\end{proof}





\section{纯但非泛纯}
\label{examples-section-pure-not-universally}

\noindent
设 $k$ 为域。令
$$
R = k[[x, xy, xy^2, \ldots]] \subset k[[x, y]].
$$
换言之，幂级数 $f \in k[[x, y]]$ 属于 $R$，当且仅当
$f(0, y)$ 是常数。特别地，$R[1/x] = k[[x, y]][1/x]$，而
$R/xR$ 是一个局部环，其极大理想的平方为零。
以 $R[y] \subset k[[x, y]]$ 表示满足 $f(0, y)$ 是 $y$ 的多项式的
幂级数 $f \in k[[x, y]]$ 所成的集合。于是 $R \to R[y]$ 是有限型但
非有限表示的环映射，并且在逆转 $x$ 后诱导同构。
还有一个满射 $R[y]/xR[y] \to k[y]$，其核的平方为零。
考虑有限表示环映射 $R \to S = R[t]/(xt - xy)$。
同样，$R[1/x] \to S[1/x]$ 是同构；在此情形下，
$S/xS \cong (R/xR)[t]/(xy)$ 映满 $k[t]$，且核幂零。
存在满射 $S \to R[y]$、$t \longmapsto y$，它在逆转 $x$ 后诱导同构，
并在模 $x$ 后诱导核幂零的满射。因此，$S \to R[y]$ 的核局部幂零。
特别地，$S \to R[y]$ 在谱上诱导泛同胚。

\medskip\noindent
首先，我们断言，$S$ 作为 $S$-模在 $R$ 上相对纯。
由于逆转 $x$ 后得到同构，只需在极大理想 $\mathfrak m \subset R$ 处
验证这一点。由于 $R$ 关于其极大理想完备，它是 Hensel 环；因此只需
验证：对每个素理想 $\mathfrak p \subset R$、
$\mathfrak p \not = \mathfrak m$，$S$ 中位于 $\mathfrak p$ 上方的唯一
素理想 $\mathfrak q$ 满足 $\mathfrak mS + \mathfrak q \not = S$。
由于 $\mathfrak p \not = \mathfrak m$，它对应于
$k[[x, y]][1/x]$ 的唯一素理想。因此，或者 $\mathfrak p = (0)$，
或者对某个不与 $x$ 相伴的不可约元素 $f \in k[[x, y]]$，有
$\mathfrak p = (f)$（这里使用 $k[[x, y]]$ 是唯一分解整环——
% P12-ERR-0346
在此插入将来的参考文献）。在第一种情形下，$\mathfrak q = (0)$，
结论显然。在第二种情形下，可将 $f$ 乘以一个单位元，使得
$f \in R[y]$（Weierstrass 预备定理；细节从略）。
于是容易看出 $R[y]/fR[y] \cong k[[x, y]]/(f)$，故 $f$ 在 $R[y]$
中定义素理想，且 $\mathfrak mR[y] + fR[y] \not = R[y]$。
由于 $S \to R[y]$ 在谱上诱导泛同胚，也就得到 $S$ 所需的结论。

\medskip\noindent
其次，我们断言，$S$ 在 $R$ 上并非泛相对纯。事实上，为说明这一点，
只需找到赋值环 $\mathcal{O}$ 和局部环映射
$R \to \mathcal{O}$，使得
$\Spec(R[y] \otimes_R \mathcal{O}) \to \Spec(\mathcal{O})$
不命中 $\Spec(\mathcal{O})$ 的闭点。等价地，需要找到
$\varphi : R \to \mathcal{O}$，使得 $\varphi(x) \not = 0$ 且
$v(\varphi(x)) > v(\varphi(xy))$，其中 $v$ 是 $\mathcal{O}$ 的赋值。
（因为这意味着 $y$ 的赋值为负。）为此，考虑按显然方式定义的环映射
$$
R
\longrightarrow
\{a_0 + a_1 x + a_2 x^2 + \ldots \mid a_0 \in k[y^{-1}], a_i \in k((y))\}
$$
可以找到一个赋值环 $\mathcal{O}$，它支配右侧在极大理想
$(y^{-1}, x)$ 处的局部化，由此即得结论。

\begin{lemma}
\label{examples-lemma-pure-not-universally-pure}
存在仿射概形之间的有限表示态射 $X \to S$，以及有限表示
$\mathcal{O}_X$-模 $\mathcal{F}$，使得 $\mathcal{F}$ 相对于 $S$ 纯，
但相对于 $S$ 并非泛纯。
\end{lemma}

\begin{proof}
见上述讨论。
\end{proof}





\section{非平坦的形式光滑环映射}
\label{examples-section-formally-smooth-nonflat}

\noindent
设 $k$ 为域。考虑 $k$-代数 $k[\mathbf{Q}]$。这是以
$x_\alpha, \alpha \in \mathbf{Q}$ 为基、乘法由
$x_\alpha x_\beta = x_{\alpha + \beta}$ 决定的 $k$-代数。
（特别地，$x_0 = 1$。）考虑 $k$-代数同态
$$
k[\mathbf{Q}] \longrightarrow k, \quad
x_\alpha \longmapsto 1.
$$
它满射，核 $J$ 由元素 $x_\alpha - 1$ 生成。我们来计算 $J/J^2$。
注意，乘以 $x_\alpha$ 在 $J/J^2$ 上是恒等映射。以 $z_\alpha$ 表示
$x_\alpha - 1$ 模 $J^2$ 的类。这些类生成 $J/J^2$。由于
$$
(x_\alpha - 1)(x_\beta - 1) = x_{\alpha + \beta} - x_\alpha - x_\beta + 1 =
(x_{\alpha + \beta} - 1) - (x_\alpha - 1) - (x_\beta - 1)
$$
可见在 $J/J^2$ 中，$z_{\alpha + \beta} = z_\alpha + z_\beta$。
$J/J^2$ 的一般元素形如 $\sum \lambda_\alpha z_\alpha$，其中
$\lambda_\alpha \in k$（只有有限多个非零）。注意，若 $k$ 的特征为
$p > 0$，则
$$
0 = pz_{\alpha/p} = z_{\alpha/p} + \ldots + z_{\alpha/p} = z_\alpha
$$
从而 $J/J^2 = 0$。若 $k$ 的特征为零，则
$$
J/J^2 = \mathbf{Q} \otimes_{\mathbf{Z}} k \cong k
$$
（细节从略），它非零。

\medskip\noindent
我们断言，当 $k$ 的特征为正时，$k[\mathbf{Q}] \to k$ 是形式光滑
环映射。事实上，假设给定实线交换图
$$
\xymatrix{
k \ar[r] \ar@{..>}[rd] & A \\
k[\mathbf{Q}] \ar[u] \ar[r]^\varphi & A' \ar[u]
}
$$
其中 $A' \to A$ 是满射，其核 $I$ 的平方为零。
为证明 $k[\mathbf{Q}] \to k$ 形式光滑，需要证明 $\varphi$ 经由 $k$
分解。由于 $\varphi(x_\alpha - 1)$ 在 $A$ 中映为零，$\varphi$ 诱导映射
$\overline{\varphi} : J/J^2 \to I$；该映射的零化恰是所需分解的障碍。
当特征为 $p > 0$ 时，$J/J^2 = 0$，所以得到所需结果，即此时
$k[\mathbf{Q}] \to k$ 形式光滑。最后，该环映射不平坦；例如，
非零因子 $x_2 - 1$ 被映为零。

\begin{lemma}
\label{examples-lemma-formally-smooth-nonflat}
存在形式光滑但非平坦的环映射。
\end{lemma}

\begin{proof}
见上述讨论。
\end{proof}





\section{非平坦的形式 étale 环映射}
\label{examples-section-formally-etale-nonflat}

\noindent
本节给出 \cite[0，例 19.10.3(i)]{EGA} 最后一句的反例
（这不在其后来的勘误表所捕获的条目之中）。考虑局部环 $A$ 及满足
$J^2 = J$ 的非零真理想 $J$ 所给出的 $A \to A/J$
（所以 $J$ 不是有限生成的）；一个代数闭非 Archimedean 域的赋值环，
取 $J$ 为其极大理想，就是这类 $(A, J)$ 的来源。
这些非平坦商映射形式 étale。事实上，假设给定交换图
$$
\xymatrix{
A/J \ar[r] & R/I \\
A \ar[u] \ar[r]^\varphi & R \ar[u]
}
$$
其中 $I$ 是环 $R$ 中满足 $I^2 = 0$ 的理想。则 $A \to R$ 唯一地经由
$A/J$ 分解，因为
$$
% P12-ERR-0347
\varphi(J) = \varphi(J^2) \subset (\varphi(J)A)^2 \subset I^2 = 0.
$$
因此，这也给出了 \cite[IV，定理 18.4.6(i)]{EGA} 中关于局部有限型、
形式 étale 态射“结构定理”的形式 étale 情形的反例
（但不是第 (ii) 部分的反例，而人们在实践中实际使用的正是第 (ii) 部分）。
后一证明中的错误是，证明的最后一步援引了错误的
% P12-ERR-0348
\cite[0，例 19.3.10(i)]{EGA}，上述反例正是由此混入的。

\begin{lemma}
\label{examples-lemma-formally-etale-not-presented}
存在形式 étale 的非平坦环映射。
\end{lemma}

\begin{proof}
见上述讨论。
\end{proof}





\section{余切复形非平凡的形式 étale 环映射}
\label{examples-section-formally-etale-nontrivial-cotangent-complex}

\noindent
设 $k$ 为域。考虑环
$$
R = k[\{x_n\}_{n \geq 1}, \{y_n\}_{n \geq 1}]/(
x_1y_1, x_{nm}^m - x_n, y_{nm}^m - y_n)
$$
令 $A$ 为在由所有 $x_n, y_n$ 生成的极大理想处的局部化，并以
$J \subset A$ 表示该极大理想。令 $B = A/J$。按构造，$J^2 = J$，
因而 $A \to B$ 形式 étale（见第
\ref{examples-section-formally-etale-nonflat} 节）。我们断言，元素
$x_1 \otimes y_1$ 是映射
$$
J \otimes_A J \longrightarrow J.
$$
之核中的非零元素。事实上，$(A, J)$ 是环
$$
R_n = k[x_n, y_n]/(x_n^n y_n^n)
$$
在相应极大理想处之局部化 $(A_n, J_n)$ 的余极限。
于是 $x_1 \otimes y_1$ 对应于元素
$x_n^n \otimes y_n^n \in J_n \otimes_{A_n} J_n$，而该元素非零
（由一个从略的显式计算可知）。由于 $\otimes$ 与余极限可交换，
即得结论。由 \cite[III 第 3.3 节]{cotangent}，$J$ 不是弱正则的。
因此，由 \cite[III 命题 3.3.3]{cotangent}，余切复形 $L_{B/A}$ 非零。
事实上，可以说得更精确。因为 $J/J^2 = 0$，有
$H_0(L_{B/A}) = \Omega_{B/A}$ 且 $H_1(L_{B/A}) = 0$。
但由 Quillen 基本谱序列的五项正合序列
（见《余切》注 \ref{cotangent-remark-elucidate-ss} 或
\cite[推论 8.2.6]{Reinhard}），以及
$\text{Tor}_2^A(B, B) = \Ker(J \otimes_A J \to J)$ 的非零性，
可得 $H_2(L_{B/A})$ 非零。

\begin{lemma}
\label{examples-lemma-formally-etale-nontrivial-cotangent-complex}
存在形式 étale 的满射环映射 $A \to B$，而 $L_{B/A}$ 不等于零。
\end{lemma}

\begin{proof}
见上述讨论。
\end{proof}





\section{平坦且形式非分歧并不蕴含形式 étale}
\label{examples-section-flat-formally-unramified-not-formally-etale}

\noindent
More on Morphisms，引理
\ref{more-morphisms-lemma-unramified-flat-formally-etale} 证明了概形的
非分歧平坦态射 $X \to S$ 形式 étale。本节的目标是给出两个例，
说明不能把“非分歧”替换为“形式非分歧”。第一个例利用完美环的特殊性质，
第二个例则说明，即使对 Noether 正则环之间的映射，该结果也会失效。

\begin{lemma}
\label{examples-lemma-perfect-closure-polynomial-ring}
设 $A = \mathbb{F}_p[T]$ 为 $\mathbb{F}_p$ 上单变量多项式环。
以 $A_{perf}$ 表示 $A$ 的完美闭包。则
$A \rightarrow A_{perf}$ 平坦且形式非分歧，但并非形式 étale。
\end{lemma}

\begin{proof}
注意，在 Frobenius 映射 $F_A : A \to A$ 下，目标副本 $A$ 是域上的
自由模，其基为 $\{1, T, \dots, T^{p - 1}\}$。因此，$F_A$ 忠实平坦；
所以，作为忠实平坦映射的余极限，$A \to A_{perf}$ 也忠实平坦。
由于 $A_{perf}$ 是完美环，相对 Frobenius
$F_{A_{perf}/A}$ 是满射。换言之，
$A_{perf} = A[A_{perf}^p]$，这立即蕴含
$\Omega_{A_{perf}/A} = 0$。于是，由 More on Morphisms，引理
\ref{more-morphisms-lemma-formally-unramified-differentials}，
$A \rightarrow A_{perf}$ 形式非分歧。

\medskip\noindent
只需证明 $A \rightarrow A_{perf}$ 不是形式光滑的。
注意，由于 $A$ 是光滑 $\mathbb{F}_p$-代数，余切复形
% P12-ERR-0349
$L_{A/\mathbb{F}_P} \simeq \Omega_{A/\mathbb{F}_p}[0]$
集中在次数 $0$；见《余切》引理 \ref{cotangent-lemma-when-projective}。
此外，由《余切》引理 \ref{cotangent-lemma-perfect-zero}，
$L_{A_{perf}/\mathbb{F}_p} = 0$ 在 $D(A_{perf})$ 中成立。
考虑 $D(A_{perf})$ 中余切复形的特异三角
$$
L_{A/\mathbb{F}_p} \otimes_A A_{perf} \to
L_{A_{perf}/\mathbb{F}_p} \to
L_{A_{perf}/A} \to
(L_{A/\mathbb{F}_p} \otimes_A A_{perf})[1]
$$
见《余切》第 \ref{cotangent-section-triangle} 节。得到
$L_{A_{perf}/A} = \Omega_{A/\mathbb{F}_p} \otimes_A A_{perf}[1]$；
即 $L_{A_{perf}/A}$ 等于置于次数 $-1$ 的秩为 $1$ 的自由
$A_{perf}$-模。因此，由 More on Morphisms，引理
\ref{more-morphisms-lemma-NL-formally-smooth} 和《余切》引理
\ref{cotangent-lemma-relation-with-naive-cotangent-complex}，
$A \rightarrow A_{perf}$ 不是形式光滑的。
\end{proof}

\noindent
下一个例也涉及素特征环，但可能更令人意外。缺点是，与前一个例相比，
它需要更多关于特征 $p$ 现象的知识。回顾一下，若素特征环 $A$ 的
Frobenius 映射有限，则称 $A$ 为 $F$-有限。

\begin{lemma}
\label{examples-lemma-completion-etale}
设 $(A, \mathfrak m, \kappa)$ 为素特征 $p > 0$ 的 Noether 局部环，
且 $[\kappa : \kappa^p] < \infty$。则从 $A$ 到其完备化的典范映射
$A \to A^\wedge$ 平坦且形式非分歧。不过，若 $A$ 正则但不优秀，
则该映射并非形式 étale。
\end{lemma}

\begin{proof}
完备化的平坦性是《代数》引理
\ref{algebra-lemma-completion-flat}。为证明该映射形式非分歧，
只需证明 $\Omega_{A^\wedge/A} = 0$；见《代数》引理
\ref{algebra-lemma-characterize-formally-unramified}。

\medskip\noindent
我们略述一个证明。选取 $x_1, \ldots, x_r \in A$，使它们映到
$\kappa$ 的一个 $p$-基 $\overline{x}_1, \ldots, \overline{x}_r$；
即 $\kappa$ 由 $\overline{x}_i$ 在 $\kappa^p$ 上极小生成。
选取 $\mathfrak m$ 的一组极小生成元 $y_1, \ldots, y_s$。
对每个 $n$，元素 $x_1, \ldots, x_r, y_1, \ldots, y_s$ 经 Frobenius
在 $(A/\mathfrak m^n)^p$ 上生成 $A/\mathfrak m^n$。略去若干细节。
由此，$F : A^\wedge \to A^\wedge$ 有限。因此，
$\Omega_{A^\wedge/A}$ 是有限 $A^\wedge$-模。另一方面，对任意
$a \in A^\wedge$ 和 $n$，可以找到 $a_0 \in A$，使得
$a - a_0 \in \mathfrak m^nA^\wedge$。于是
$\text{d}(a) \in \bigcap \mathfrak m^n \Omega_{A^\wedge/A}$；
由《代数》引理
\ref{algebra-lemma-intersect-powers-ideal-module-zero}，这蕴含
$\text{d}(a)$ 为零。因此，$\Omega_{A^\wedge/A} = 0$。

\medskip\noindent
假设 $A$ 正则。则由 Cohen 结构定理，
$x_1, \ldots, x_r, y_1, \ldots, y_s$ 是环 $A^\wedge$ 的一个
$p$-基；即有
$$
A^\wedge = \bigoplus\nolimits_{I, J} (A^\wedge)^p
x_1^{i_1} \ldots x_r^{i_r} y_1^{j_1} \ldots y_s^{j_s}
$$
其中 $I = (i_1, \ldots, i_r)$、$J = (j_1, \ldots, j_s)$，且
$0 \leq i_a, j_b \leq p - 1$。细节从略。特别地，
$\Omega_{A^\wedge}$ 是自由 $A^\wedge$-模，其基为
$\text{d}(x_1), \ldots, \text{d}(x_r), \text{d}(y_1), \ldots, \text{d}(y_s)$。

\medskip\noindent
现在，若 $A \to A^\wedge$ 形式 étale，甚至只需形式光滑，则由《代数》
命题 \ref{algebra-proposition-characterize-formally-smooth}，
$\NL_{A^\wedge/A}$ 在次数 $-1, 0$ 的上同调均为零。
对环映射 $\mathbf{F}_p \to A \to A^\wedge$ 使用 Jacobi--Zariski 序列
（《代数》引理 \ref{algebra-lemma-exact-sequence-NL}），得到同构
$\Omega_A \otimes_A A^\wedge \cong \Omega_{A^\wedge}$。
所以，$\Omega_A$ 以
$\text{d}(x_1), \ldots, \text{d}(x_r), \text{d}(y_1), \ldots, \text{d}(y_s)$
为自由基。考察分式域并利用 $A$ 正规，可得 $a \in A$ 是 $p$ 次幂，
当且仅当它在 $A^\wedge$ 中的像是 $p$ 次幂
（细节从略；使用《代数》引理
\ref{algebra-lemma-derivative-zero-pth-power}）。另一个推论是算子
$\partial/\partial x_a$ 和 $\partial/\partial y_b$ 在 $A$ 上有定义。

\medskip\noindent
% P12-ERR-0350
我们将证明，上述事实蕴含 $A$ 在 $A^p$ 上有限，并以
$x_1, \ldots, x_r, y_1, \ldots, y_s$ 为 $p$-基。
由 Kunz 的一个结果，这将与 $A$ 的非优秀性矛盾；见
\cite[推论 2.6]{Kun76}。具体地，设 $a \in A$，并写成
$$
a =  \sum\nolimits_{I, J} (a_{I, J})^p
x_1^{i_1} \ldots x_r^{i_r} y_1^{j_1} \ldots y_s^{j_s}
$$
其中 $a_{I, J} \in A^\wedge$。为完成证明，只需说明
$a_{I, J} \in A$。在等式两侧作用算子
$$
(\partial/\partial x_1)^{p - 1} \ldots
(\partial/\partial x_r)^{p - 1}
(\partial/\partial y_1)^{p - 1} \ldots
(\partial/\partial y_s)^{p - 1}
$$
可得 $a_{I, J}^p \in A$，其中
$I = (p - 1, \ldots, p - 1)$ 且 $J = (p - 1, \ldots, p - 1)$。
由上面的说明，这也蕴含 $a_{I, J} \in A$。把 $a$ 替换为
$a' = a - a_{I, J}^p x^I y^J$ 后，
% P12-ERR-0351
可以使用低 $1$ 阶的微分算子，得到另一个系数 $a_{I, J}$ 属于 $A$。
依此类推。
\end{proof}

\begin{remark}
\label{examples-remark-reference-existence-regular-nonexcellent-rings}
剩余域具有有限 $p$-基的非优秀正则环，甚至可以在特征 $p$ 的代数闭域
$k = \overline{k}$ 上的 $\mathbb{P}^2_k$ 的函数域中构造。
见 \cite[$\mathsection 4.1$]{DS18}。
\end{remark}

\noindent
引理 \ref{examples-lemma-completion-etale} 的证明实际上还说明了稍强的结论。

\begin{lemma}
\label{examples-lemma-excellent-regular-local-rings}
设 $(A, \mathfrak m, \kappa)$ 为特征 $p > 0$ 的正则局部环。
假设 $[\kappa : \kappa^p] < \infty$。则 $A$ 优秀，当且仅当
$A \to A^\wedge$ 形式 étale。
\end{lemma}

\begin{proof}
反向蕴含由引理 \ref{examples-lemma-completion-etale} 得到。对于正向蕴含，
注意，我们已经由引理 \ref{examples-lemma-completion-etale} 知道
$A \to A^\wedge$ 形式非分歧；等价地，
$\Omega_{A^\wedge/A}$ 为零。因此，当 $A$ 优秀时，只需证明完备化映射
形式光滑。由 Néron--Popescu 解奇异，$A \to A^\wedge$ 可写成光滑
$A$-代数的滤过余极限（《平滑化环映射》定理
\ref{smoothing-theorem-popescu}）。因此，$\NL_{A^\wedge/A}$ 在次数
$-1$ 的上同调为零。故由《代数》命题
\ref{algebra-proposition-characterize-formally-smooth}，
$A \to A^\wedge$ 形式光滑。
\end{proof}



\section{由幂等元集合生成的理想与局部化}
\label{examples-section-ideal-locally-idempotents}

\noindent
设 $R$ 为环。考虑环
$$
B(R) = R[x_n; n \in \mathbf{Z}]/(x_n(x_n - 1), x_nx_m; n \not = m)
$$
容易证明，每个素理想 $\mathfrak q \subset B(R)$ 或者形如
$$
\mathfrak q = \mathfrak pB(R) + (x_n; n \in \mathbf{Z})
$$
或者形如
$$
\mathfrak q =
\mathfrak pB(R) + (x_n - 1) + (x_m; n \not = m, m \in \mathbf{Z}).
$$
因此可见
$$
\Spec(B(R)) =
\Spec(R) \amalg \coprod\nolimits_{n \in \mathbf{Z}} \Spec(R)
$$
但其拓扑并不只是无交并拓扑。它具有如下性质：由
$n \in \mathbf{Z}$ 标号的每一份都是开子概形，确切地说就是标准开集
$D(x_n)$。“中央”的那一份 $\Spec(R)$ 位于其余任意无穷多份
$\Spec(R)$ 之并的闭包中。注意，最后这一份 $\Spec(R)$ 由理想
$(x_n, n \in \mathbf{Z})$ 截出，而该理想由幂等元 $x_n$ 生成。
因此，若 $\Spec(R)$ 连通，则上述分解恰好是 $\Spec(B(R))$
分解为连通分支的分解。

\medskip\noindent
其次，令
$A = \mathbf{C}[x, y]/((y - x^2 + 1)(y + x^2 - 1))$。
$A$ 的谱由两个不可约分支
$C_1 = \Spec(A_1)$、$C_2 = \Spec(A_2)$ 组成，其中
$A_1 = \mathbf{C}[x, y]/(y - x^2 + 1)$，且
$A_2 = \mathbf{C}[x, y]/(y + x^2 - 1)$。注意，它们分别由
$(x, y) = (t, t^2 - 1)$ 和 $(x, y) = (t, -t^2 + 1)$ 参数化，
并相交于 $P = (-1, 0)$ 和 $Q = (1, 0)$。我们可以构造 $B(A)$
的一个扭曲版本，按如下方式把 $B(A_1)$ 与 $B(A_2)$ 粘合起来：
在 $P$ 上方，使
$x_n \in B(A_1) \otimes \kappa(P)$
对应于 $x_n \in B(A_2) \otimes \kappa(P)$；但在 $Q$ 上方，使
$x_n \in B(A_1) \otimes \kappa(Q)$
对应于 $x_{n + 1} \in B(A_2) \otimes \kappa(Q)$。
把所得的 $A$-代数记作 $B^{twist}(A)$。细节从略。由构造，
$B^{twist}(A)$ 在 $A$ 上 Zariski 局部同构于未扭曲的版本。
具体而言，在主开集 $\Spec(A) \setminus \{P\}$ 和主开集
$\Spec(A) \setminus \{Q\}$ 上都是如此。然而，我们所选的粘合
产生了足够的“单值作用”，使得 $\Spec(B^{twist}(A))$ 连通
（细节从略）。最后，存在一个中央副本
% P12-ERR-0357
$\Spec(A) \to \Spec(B^{twist}(A))$，它给出一个闭子概形；其理想在
$B^{twist}(A)$ 上 Zariski 局部由
幂等元生成的理想截出，但整体上并非如此（因为 $B^{twist}(A)$
没有非平凡幂等元）。

\begin{lemma}
\label{examples-lemma-not-generated-idempotents}
存在一个仿射概形 $X = \Spec(A)$ 及其闭子概形 $T \subset X$，使得
$T$ 在 $X$ 上 Zariski 局部由幂等元生成的理想截出，但 $T$
并非由某个幂等元生成的理想截出。
\end{lemma}

\begin{proof}
见上述讨论。
\end{proof}



% P12-ERR-0352
\section{识别局部环但并非 ind-étale 的环映射}
\label{examples-section-not-ind-etale}

\noindent
注意，第 \ref{examples-section-ideal-locally-idempotents} 节构造的环映射
$R \to B(R)$ 是有限个 $R$ 的副本之积的余极限。因此
$R \to B(R)$ 是 ind-Zariski 的；见《Pro-étale 上同调》定义
\ref{proetale-definition-ind-zariski}。其次，考虑第
\ref{examples-section-ideal-locally-idempotents} 节构造的环映射
$A \to B^{twist}(A)$。由于该环映射在 $\Spec(A)$ 上 Zariski 局部
同构于某个 ind-Zariski 环映射 $R \to B(R)$，可知它识别局部环
（见《Pro-étale 上同调》引理
\ref{proetale-lemma-ind-zariski-implies}）。第
\ref{examples-section-ideal-locally-idempotents} 节的讨论表明，存在一个截面
$B^{twist}(A) \to A$，其核并非由幂等元生成。现在，若
$A \to B^{twist}(A)$ 是 ind-étale 的，即
$B^{twist}(A) = \colim A_i$，其中 $A \to A_i$ 为 étale，则
$A_i \to A$ 的核将由一个幂等元生成
（《代数》引理 \ref{algebra-lemma-map-between-etale} 和
\ref{algebra-lemma-surjective-flat-finitely-presented}）。
这与上述结果矛盾。

\begin{lemma}
\label{examples-lemma-not-ind-etale}
存在一个识别局部环但并非 ind-étale 的环映射 $A \to B$。
尤其，它并非 ind-Zariski。
\end{lemma}

\begin{proof}
见上述讨论。
\end{proof}




\section{与内射模伴随的非松弛拟凝聚层}
\label{examples-section-nonflasque}

\noindent
关于这类例子的更多内容，见 \cite[Exposé II，附录 I]{SGA6}，
其中 Illusie 解释了 Verdier 给出的一些例子。

\medskip\noindent
考虑仿射概形 $X = \Spec(A)$，其中
$$
A = k[x, y, z_1, z_2, \ldots]/(x^nz_n)
$$
是《性质》例 \ref{properties-example-does-not-work-in-general}
中的环。令 $I = (x) \subset A$。考虑 $X$ 的拟紧开集 $U = D(x)$。
我们已经在上述引文中看到，存在一个截面
$s \in \mathcal{O}_X(U)$，它对任意 $n \geq 0$ 都不来自某个
$A$-模映射 $I^n \to A$。

\medskip\noindent
令 $\alpha : A \to J$ 为 $A$ 到某个内射 $A$-模中的嵌入。
令 $Q = J/\alpha(A)$，并记商映射为 $\beta : J \to Q$。
我们断言，映射
$$
\Gamma(X, \widetilde{J})
\longrightarrow
\Gamma(U, \widetilde{J})
$$
不是满射。确切地说，我们断言 $\alpha(s)$ 不在其像中。
为证明这一点，使用反证法。假设 $x \in J$ 是一个在 $U$
上限制为 $\alpha(s)$ 的元素。于是 $\beta(x) \in Q$ 在 $U$
上的限制为 $0$。因此，对某个 $n$ 有 $I^n\beta(x) = 0$；见
《性质》引理
\ref{properties-lemma-sections-over-quasi-compact-open-in-affine}。
这蕴含我们得到一个态射
$\varphi : I^n \to A$，$h \mapsto \alpha^{-1}(hx)$。容易看出，
经《性质》引理
\ref{properties-lemma-sections-over-quasi-compact-open-in-affine}
中的映射，这个态射 $\varphi$ 给出截面 $s$，矛盾。

\begin{lemma}
\label{examples-lemma-nonflasque}
存在一个仿射概形 $X = \Spec(A)$ 和一个内射 $A$-模 $J$，使得
$\widetilde{J}$ 不是 $X$ 上的松弛层。即使 $U$ 是标准开集，限制映射
$\Gamma(X, \widetilde{J}) \to \Gamma(U, \widetilde{J})$
也未必是满射。
\end{lemma}

\begin{proof}
见上述讨论。
\end{proof}

\noindent
事实上，可以用类似构造得到一个内射模的例子，使其伴随拟凝聚层在某个
拟紧开集上具有非零上同调。具体地，从环
$$
A = k[x, y, w_1, u_1, w_2, u_2, \ldots]/(x^nw_n, y^nu_n, u_n^2, w_n^2)
$$
出发，其中 $k$ 为域。选取一个内射映射 $A \to I$，其中 $I$
为内射 $A$-模。我们断言，$A_{xy} \subset I_{xy}$ 中的元素
$1/xy$ 不在 $I_x \oplus I_y \to I_{xy}$ 的像中。反证之，假设
$$
\frac{1}{xy} = \frac{i}{x^n} + \frac{j}{y^n}
$$
对某个 $n \geq 1$ 及 $i, j \in I$ 成立。清除分母，得到
$$
(xy)^{n + m - 1} = x^my^{n + m}i + x^{n + m}y^mj
$$
对某个 $m \geq 0$ 成立。乘以 $u_{n + m}w_{n + m}$，可见
% P12-ERR-0353
$u_{n + m}w_{n + m}(xy)^{n + m - 1} = 0$ 在 $A$ 中成立，
这就是所需的矛盾。
令 $U = D(x) \cup D(y) \subset X = \Spec(A)$。对任意 $A$-模
$M$，由 Mayer--Vietoris 序列有正合列
$$
0 \to H^0(U, \widetilde{M}) \to M_x \oplus M_y \to M_{xy} \to
H^1(U, \widetilde{M}) \to 0
$$
由此可知 $H^1(U, \widetilde{I})$ 非零。

\begin{lemma}
\label{examples-lemma-nonvanishing}
存在一个仿射概形 $X = \Spec(A)$，其底层拓扑空间是 Noetherian 的；
还存在一个内射 $A$-模 $I$，使得 $\widetilde{I}$ 在 $X$ 的某个拟紧
开集 $U$ 上具有非零的 $H^1$。
\end{lemma}

\begin{proof}
见上述讨论。注意，作为拓扑空间，
$\Spec(A) = \Spec(k[x, y])$。
\end{proof}








\section{非分离平坦群概形}
\label{examples-section-non-separated-group-scheme}

\noindent
域上的每个群概形都是分离的；见《群胚》引理
\ref{groupoids-lemma-group-scheme-over-field-separated}。
对一般基上的群概形，这不成立。

\medskip\noindent
设 $k$ 为域。令 $S = \Spec(k[x]) = \mathbf{A}^1_k$。
令 $G$ 为原点 $0$ 加倍的仿射直线（见《概形》例
\ref{schemes-example-affine-space-zero-doubled}），并把它视为 $S$
上的概形。因此，$G \to S$ 的一个纤维或者是单点集，或者是一个二元集
（其中一个元素在 $U$ 中，另一个在 $V$ 中）。于是，可以在这些纤维上
赋予群结构（令 $U$ 中的元素为群结构的零元）。更精确地说，$G$
具有两个开集 $U, V$，它们同构地映到 $S$，而 $U \cap V$
同构地映到 $S \setminus \{0\}$。于是
$$
G \times_S G = U \times_S U \cup V \times_S U \cup
U \times_S V \cup V \times_S V
$$
其中每一片都同构于 $S$。因此，可以把乘法
$m : G \times_S G \to G$ 定义为唯一的 $S$-态射：它把第一片和
最后一片映入 $U$，把中间两片映入 $V$。这与上述逐点描述相符。
我们略去它确实定义群概形结构的验证。

\begin{lemma}
\label{examples-lemma-non-separated-group-scheme}
仿射直线上存在一个有限型平坦群概形，它不是分离的。
\end{lemma}

\begin{proof}
见上述讨论。
\end{proof}

\begin{lemma}
\label{examples-lemma-non-quasi-separated-group-scheme}
无穷维仿射空间上存在一个有限型平坦群概形，它不是拟分离的。
\end{lemma}

\begin{proof}
上述同一构造可以用无穷维仿射空间
$S = \mathbf{A}^\infty_k = \Spec k[x_1, x_2, \ldots]$ 作为基来实施，
并把与极大理想 $(x_1, x_2, \ldots)$ 对应的原点 $0 \in S$
作为在 $G$ 中加倍的闭点。所得群概形 $G \rightarrow S$
并非拟分离的，正如《概形》例
\ref{schemes-example-not-quasi-separated} 中所解释的那样。
\end{proof}



\section{单位分支平坦的非平坦群概形}
\label{examples-section-non-flat-group-scheme}

\noindent
设 $X \to S$ 为概形的单态射。令 $G = S \amalg X$。
令 $m : G \times_S G \to G$ 为 $S$-态射
$$
G \times_S G = X \times_S X \amalg X \amalg X \amalg S
\longrightarrow G = X \amalg S
$$
它把分量 $X \times_S X$ 和 $S$ 映入 $S$，并以恒等态射把两个
分量 $X$ 映入 $X$。这定义了一个群律。为看清这一点，需要证明
$m \circ (m \times \text{id}_G) = m \circ (\text{id}_G \times m)$
是 $G \times_S G \times_S G \to G$ 的两个相同映射。像上面一样把
$G \times_S G \times_S G$ 分解为各分量，可见只需在这 $8$ 片中的
每一片上验证。每片分别同构于 $S$、$X$、$X \times_S X$ 或
$X \times_S X \times_S X$ 之一。此外，两个映射分别把这些片映到
$S$、$X$、$S$、$X$。另一方面，$X \to S$ 是单态射这一事实蕴含
$X \times_S X \cong X$ 和 $X \times_S X \times_S X \cong X$，
并且在每种情形下恰有一个 $S$-态射 $S \to S$ 或 $X \to X$。
因此
$m \circ (m \times \text{id}_G) = m \circ (\text{id}_G \times m)$。
于是，取 $X \to S$ 为任意非平坦概形单态射（例如闭浸入），
便得到基 $S$ 上的一个群概形例子：其单位分支为 $S$（因而平坦），
但它本身不平坦。

\begin{lemma}
\label{examples-lemma-non-flat-group-scheme}
存在基 $S$ 上的群概形 $G$，其单位分支在 $S$ 上平坦，但它本身在
$S$ 上不平坦。
\end{lemma}

\begin{proof}
见上述讨论。
\end{proof}




\section{域上的非分离群代数空间}
\label{examples-section-non-separated-group-space}

\noindent
域上的每个群概形都是分离的；见《群胚》引理
\ref{groupoids-lemma-group-scheme-over-field-separated}。
但域上的群代数空间并非如此（不过，正面结果见本节末尾）。

\medskip\noindent
设 $k$ 为特征零的域。考虑《空间》例
\ref{spaces-example-affine-line-translation} 中的代数空间
$G = \mathbf{A}^1_k/\mathbf{Z}$。由构造，$G$ 是 $k$ 上的概形范畴中的预层
$$
T \longmapsto \Gamma(T, \mathcal{O}_T) / \mathbf{Z}
$$
在 fppf 拓扑下的伴随层。预层上显然的加法法则诱导出加法
$m : G \times G \to G$，它使 $G$ 成为 $\Spec(k)$ 上的群代数空间。
注意，$G$ 不是分离的（甚至不是拟分离或局部分离的）。另一方面，
$G \to \Spec(k)$ 却是有限型的！

\begin{lemma}
\label{examples-lemma-non-separated-group-space}
域上存在一个有限型群代数空间，它不是分离的（甚至不是拟分离或
局部分离的）。
\end{lemma}

\begin{proof}
见上述讨论。
\end{proof}

\noindent
正面结果如下：若群代数空间 $G$ 拟分离、局部分离，或者更一般地是
适度代数空间，则 $G$ 实际上是分离的；见《空间中的更多群胚》引理
\ref{spaces-more-groupoids-lemma-group-scheme-over-field-separated}。
此外，域上的有限型分离群代数空间实际上是概形；见《空间中的更多群胚》
引理
\ref{spaces-more-groupoids-lemma-group-space-scheme-locally-finite-type-over-k}。
证明的思路是，概形轨迹是稠密开集；见《空间的性质》命题
\ref{spaces-properties-proposition-locally-quasi-separated-open-dense-scheme}
或《适度空间》定理
\ref{decent-spaces-theorem-decent-open-dense-scheme}。
平移这个开集，便可看出 $G$ 的每一点都有一个为概形的开邻域。




% P12-ERR-0354
\section{étale 态射纤维中各点之间的专化}
\label{examples-section-specializations-fibre-etale}

\noindent
若 $f : X \to Y$ 是概形的 étale 态射，或更一般地是局部拟有限态射，
则纤维的各点之间不存在专化；见《态射》引理
\ref{morphisms-lemma-locally-quasi-finite-fibres}。
然而，对代数空间的态射，这一般不成立。

\medskip\noindent
为给出一个例子，设 $k$ 为域。令
$$
P = k[u, u^{-1}, y, \{x_n\}_{n \in \mathbf{Z}}].
$$
考虑由自同构 $\tau$ 生成的 $\mathbf{Z}$ 在 $P$ 上通过 $k$-代数映射
给出的作用，其中
$\tau(u) = u$、$\tau(y) = uy$ 且 $\tau(x_n) = x_{n + 1}$。
对 $d \geq 1$，令
$I_d = ((1 - u^d)y, x_n - x_{n + d}, n \in \mathbf{Z})$。
于是，
% P12-ERR-0355
$V(I_d) \subset \Spec(P)$ 是 $\tau^d$ 的不动点轨迹。
令 $S \subset P$ 为由 $y$ 及所有
$1 - u^d$（$d \in \mathbf{N}$）生成的乘法子集。于是可见，
$\mathbf{Z}$ 自由作用于 $U = \Spec(S^{-1}P)$。令
$X = U/\mathbf{Z}$ 为商代数空间；见《空间》定义
\ref{spaces-definition-quotient}。

\medskip\noindent
考虑 $S^{-1}P$ 中的素理想
$\mathfrak p_n = (x_n, x_{n + 1}, \ldots)$。注意
$\tau(\mathfrak p_n) = \mathfrak p_{n + 1}$。因此，它们中的每一个都
% P12-ERR-0356
定义一点 $\xi_n \in U$，其在 $X$ 中的像都是同一点 $x$，而该点属于
$X$。
此外，有专化
$$
\ldots \leadsto \xi_n \leadsto \xi_{n - 1} \leadsto \ldots
$$
因此，$U \to X$ 就是所求类型的例子。

\begin{lemma}
\label{examples-lemma-specializations-fibre-etale}
存在代数空间的 étale 态射 $f : X \to Y$，以及 $|X|$ 中点的一个
非平凡专化 $x \leadsto x'$，使得 $|Y|$ 中有
$f(x) = f(x')$。
\end{lemma}

\begin{proof}
见上述讨论。
\end{proof}







\section{不是 fppf 挠子的挠子}
\label{examples-section-torsor-not-fppf}

\noindent
在《群胚》注 \ref{groupoids-remark-fun-with-torsors} 中，我们提出了如下问题：
任意 $G$-挠子是否都是 fppf 拓扑下的 $G$-挠子？本节说明，情况并非
总是如此。

\medskip\noindent
设 $k$ 为域。以下所有概形和叠均在 $k$ 上。令
$G \to \Spec(k)$ 为群概形
$$
G = (\mu_{2, k})^\infty =
\mu_{2, k} \times_k \mu_{2, k} \times_k \mu_{2, k} \times_k \ldots =
\lim_n (\mu_{2, k})^n
$$
其中 $\mu_{2, k}$ 是 $\Spec(k)$ 上的二次单位根群概形；见《群胚》例
% P12-ERR-0358
\ref{groupoids-example-roots-of-unity}。由于 $G$ 是仿射概形的逆极限，
可见它是仿射群概形。事实上，它是环
$k[t_1, t_2, t_3, \ldots]/(t_i^2 - 1)$ 的谱。乘法映射
$m : G \times_k G \to G$ 在代数层面由
$t_i \mapsto t_i \otimes t_i$ 给出。

\medskip\noindent
我们断言，$k$ 上的任意 $G$-挠子都形如
$$
P = \Spec(k[x_1, x_2, x_3, \ldots]/(x_i^2 - a_i))
$$
其中 $a_i \in k^*$，而 $G$-作用 $G \times_k P \to P$
在代数层面由 $x_i \to t_i \otimes x_i$ 给出。证明从略。
事实上，对本例只需知道 $P$ 是 $G$-挠子；这是清楚的，因为在
$k' = k(\sqrt{a_1}, \sqrt{a_2}, \ldots)$ 上，概形 $P$ 以
$G$-等变方式同构于 $G$。注意，$P$ 平凡当且仅当 $k' = k$，
因为若 $P$ 有 $k$-有理点，则所有 $a_i$ 都是平方。

\medskip\noindent
我们断言，$P$ 是 fppf 挠子，当且仅当域扩张
$k' = k(\sqrt{a_1}, \sqrt{a_2}, \ldots)/k$ 有限。若 $k'$ 在 $k$
上有限，则
$\{\Spec(k') \to \Spec(k)\}$
是平凡化 $P$ 的 fppf 覆盖，因此 $P$ 确实是 fppf 挠子。反之，
假设 $P$ 是 fppf $G$-挠子。这意味着存在 fppf 覆盖
$\{S_i \to \Spec(k)\}$，使得每个 $P_{S_i}$ 都平凡。选取一个
$i$，使 $S_i$ 非空，并令 $s \in S_i$ 为闭点。由《簇》引理
\ref{varieties-lemma-locally-finite-type-Jacobson}，域扩张
$\kappa(s)/k$ 有限；由构造，$P_{\kappa(s)}$ 有一个
$\kappa(s)$-有理点。因此
$k \subset k' \subset \kappa(s)$，且 $k'$ 在 $k$ 上有限。

\medskip\noindent
为得到一个显式例子，例如可取 $k = \mathbf{Q}$ 和 $a_i = i$
（或者，如果愿意，可令 $a_i$ 为第 $i$ 个素数）。

\begin{lemma}
\label{examples-lemma-torsors-principal-spaces-not-equal}
设 $S$ 为概形，$G$ 为 $S$ 上的群概形。分类主齐 $G$-空间的叠
$G\textit{-Principal}$（见《叠的例子》小节
\ref{examples-stacks-subsection-principal-homogeneous-spaces}）与分类 fppf
$G$-挠子的叠 $G\textit{-Torsors}$（见《叠的例子》小节
\ref{examples-stacks-subsection-fppf-torsors}）一般并不等价。
\end{lemma}

\begin{proof}
上述讨论表明，函子
$G\textit{-Torsors} \to G\textit{-Principal}$ 一般不是本质满的。
\end{proof}








\section{具有拟紧平坦覆盖但非代数的叠}
\label{examples-section-not-algebraic-stack}

\noindent
本节简要描述 Brian Conrad 给出的一个例子。可以在
\href{https://mathoverflow.net/questions/15082/fpqc-covers-of-stacks/15269#15269}{这里}
在线找到该例。我们的例子略有不同。

\medskip\noindent
设 $k$ 为代数闭域。以下所有概形和叠均在 $k$ 上。
令 $G \to \Spec(k)$ 为仿射群概形。在《叠的例子》引理
\ref{examples-stacks-lemma-classifying-stacks} 中，我们给出了若干种彼此
等价的方式，把 $\mathcal{X} = [\Spec(k)/G]$ 看成
$(\Sch/\Spec(k))_{fppf}$ 上的群胚叠。特别地，$\mathcal{X}$ 分类
fppf $G$-挠子。更精确地说，一个 $1$-态射
$T \to \mathcal{X}$ 对应于 $T$ 上的一个 fppf $G_T$-挠子 $P$，而
$2$-箭头对应于挠子的同构。由此可知，对角 $1$-态射
$$
\Delta :
\mathcal{X}
\longrightarrow
\mathcal{X} \times_{\Spec(k)} \mathcal{X}
$$
可表且仿射。确切地说，给定概形 $T/k$ 上任意一对 fppf
$G_T$-挠子 $P_1, P_2$，概形 $\mathit{Isom}(P_1, P_2)$ 在 $T$
上仿射。$\Spec(k)$ 上平凡的 $G$-挠子定义一个 $1$-态射
$$
f : \Spec(k) \longrightarrow \mathcal{X}.
$$
我们断言，这是满的 $1$-态射。理由很简单：由定义，对任意 $1$-态射
$T \to \mathcal{X}$，存在一个 fppf 覆盖 $\{T_i \to T\}$，使得
$P_{T_i}$ 同构于平凡的 $G_{T_i}$-挠子。因此，复合
$T_i \to T \to \mathcal{X}$ 经过 $f$ 分解。于是显然，投影
$T \times_\mathcal{X} \Spec(k) \to T$ 是满射（这正是我们定义
$f$ 为满射的方式；见《代数叠》定义
\ref{algebraic-definition-relative-representable-property}）。
类似地可证明 $f$ 拟紧且平坦（细节从略）。还要在此记录如下观察：
$$
\Spec(k) \times_\mathcal{X} \Spec(k) \cong G
$$
是 $k$ 上概形的同构。

\medskip\noindent
假设存在满光滑态射 $p : U \to \mathcal{X}$，其中 $U$ 为概形。
考虑纤维积
$$
\xymatrix{
W \ar[d] \ar[r] & U \ar[d] \\
\Spec(k) \ar[r] & \mathcal{X}
}
$$
于是，$W$ 是 $k$ 上非空光滑概形，因而有一个 $k$-点。
这意味着 $f$ 可经过 $U$ 分解。因此得到
$$
G \cong
\Spec(k) \times_\mathcal{X} \Spec(k) \cong
(\Spec(k) \times_k \Spec(k))
\times_{(U \times_k U)}
(U \times_\mathcal{X} U)
$$
由于投影 $U \times_\mathcal{X} U \to U$ 按假设是光滑的，可知
$U \times_\mathcal{X} U \to U \times_k U$ 是局部有限型的；见
《态射》引理 \ref{morphisms-lemma-permanence-finite-type}。
因此，在这种情形下，$G$ 在 $k$ 上局部有限型。综上，我们证明了
下述引理（它可以显著推广）。

\begin{lemma}
\label{examples-lemma-BG-algebraic}
设 $k$ 为域，$G$ 为 $k$ 上的仿射群概形。若叠
$[\Spec(k)/G]$ 有一个由概形给出的光滑覆盖，则 $G$ 在 $k$ 上有限型。
\end{lemma}

\begin{proof}
见上述讨论。
\end{proof}

\noindent
为得到本节标题所述的显式例子，例如取第
\ref{examples-section-torsor-not-fppf} 节中的群概形
$G = (\mu_{2, k})^\infty$；它在 $k$ 上不是局部有限型。由上述讨论可见，
$\mathcal{X} = [\Spec(k)/G]$ 具有《代数叠》定义
\ref{algebraic-definition-algebraic-stack} 中的性质 (1) 和 (2)，
但不具有性质 (3)。因此 $\mathcal{X}$ 不是代数叠。另一方面，确实存在
一个概形 $U$ 和一个满、平坦、拟紧态射 $U \to \mathcal{X}$；它就是
上面研究过的态射 $f : \Spec(k) \to \mathcal{X}$。






\section{在对象上保持极限，但不保持极限}
\label{examples-section-limit-preserving}

\noindent
设 $S$ 为非空概形，$\mathcal{G}$ 为 $(\Sch/S)_{fppf}$ 上的内射 Abel 层。
我们得到 $S$ 上的群胚叠
$$
\mathcal{G}\textit{-Torsors} \longrightarrow (\Sch/S)_{fppf}
$$
见《叠的例子》引理
\ref{examples-stacks-lemma-torsors-sheaf-stack-in-groupoids}。
这个叠在 $(\Sch/S)_{fppf}$ 上的对象层面保持极限（见《可表性判据》第
\ref{criteria-section-limit-preserving} 节），因为每个
$\mathcal{G}$-挠子都是平凡的。另一方面，
$\mathcal{G}\textit{-Torsors}$ 一般不保持极限（见《Artin 公理》定义
\ref{artin-definition-limit-preserving}），因为作为层，$\mathcal{G}$
未必保持极限。例如，取任意非零内射层 $\mathcal{I}$，并令
$\mathcal{G} = \prod_{n \in \mathbf{Z}} \mathcal{I}$，便得到一个例子。

\begin{lemma}
\label{examples-lemma-limit-preserving-on-objects-not-limit-preserving}
设 $S$ 为非空概形。存在群胚叠
$p : \mathcal{X} \to (\Sch/S)_{fppf}$，使得 $p$ 在对象层面保持极限，
但 $\mathcal{X}$ 不保持极限。
\end{lemma}

\begin{proof}
见上述讨论。
\end{proof}



\section{一个非代数分类叠}
\label{examples-section-non-algebraic}

\noindent
令 $S = \Spec(\mathbf{F}_p)$，并以 $\mu_p$ 表示 $S$ 上的
$p$ 次单位根群概形。在《空间中的群胚》第
\ref{spaces-groupoids-section-stacks} 节中，我们引入了商叠
$[S/\mu_p]$；在《叠的例子》第
\ref{examples-stacks-section-group-quotient-stacks} 节中，我们证明了
$[S/\mu_p]$ 是 fppf $\mu_p$-挠子的分类叠：给定 $S$ 上概形 $T$，
范畴 $\Mor_S(T, [S/\mu_p])$ 典范等价于 $T$ 上 fppf
$\mu_p$-挠子的范畴。最后，我们已经在《可表性判据》定理
\ref{criteria-theorem-flat-groupoid-gives-algebraic-stack} 中看到，
$[S/\mu_p]$ 是代数叠。

\medskip\noindent
现在可以提出如下问题：“分类 étale $\mu_p$-挠子的群胚纤维范畴
$\mathcal{S}$ 又如何？”（换言之，$\mathcal{S}$ 是 $\Sch/S$ 上的范畴，
其在概形 $T$ 上方的纤维范畴是 $T$ 上 étale $\mu_p$-挠子的范畴。）

\medskip\noindent
第一个问题是，它并非 fppf 拓扑下的叠，因为对象的下降不会成立。
例如，$T = \Spec(\mathbf{F}_p(T))$ 上的 $\mu_p$-挠子
% P12-ERR-0359
$\Spec(\mathbf{F}_p(t)[x]/(x^p - t))$ 在 fppf 局部平凡，
但并非 étale 局部平凡。

\medskip\noindent
解决第一个问题的一种办法是改用 étale 拓扑；在这种情形下，对象的下降
确实成立。事实上，$\mathcal{S}$ 是 $(\Sch/S)_\etale$ 上的群胚叠。
此外，对角
$\Delta : \mathcal{S} \to \mathcal{S} \times \mathcal{S}$
也是（由概形）可表的。这是因为，给定两个 $\mu_p$-挠子（无论它们是否
étale 局部平凡），它们之间的同构层均由概形表示。

\medskip\noindent
因此，最后可以问：是否存在概形 $U$ 和光滑满 $1$-态射
$U \to \mathcal{S}$？我们将以两种方式证明这是不可能的：一种是直接
论证（建议读者跳过），另一种使用一般结果。

\medskip\noindent
直接论证（梗概）：注意，$1$-态射
$\mathcal{S} \to \Spec(\mathbf{F}_p)$ 满足形式光滑性的无穷小提升判据。
这是因为，给定概形的一阶无穷小加厚 $T \to T'$，映射
$\mu_p(T') \to \mu_p(T)$ 的核同构于 $T$ 在 $T'$ 中的理想层的截面，
因而
$H^1_\etale(T, \mu_p) = H^1_\etale(T', \mu_p)$。
此外，$\mathcal{S}$ 是保持极限的叠。因此，若
$U \to \mathcal{S}$ 光滑，则 $U \to \Spec(\mathbf{F}_p)$ 保持极限并满足
% P12-ERR-0360
形式光滑性的无穷小提升判据。这蕴含 $U$ 在 $\mathbf{F}_p$ 上光滑。
特别地，$U$ 是约化的，故
$H^1_\etale(U, \mu_p) = 0$。
因此，$U \to \mathcal{S}$ 分解为
$U \to \Spec(\mathbf{F}_p) \to \mathcal{S}$，且第一支箭头光滑。
由光滑性的下降可知，$U \to \mathcal{S}$ 光滑将蕴含
$\Spec(\mathbf{F}_p) \to \mathcal{S}$ 光滑。然而，事实并非如此，因为
$\Spec(\mathbf{F}_p) \times_\mathcal{S} \Spec(\mathbf{F}_p)$
是 $\mu_p$，它在 $\Spec(\mathbf{F}_p)$ 上不光滑。


\medskip\noindent
结构论证：在《可表性判据》第 \ref{criteria-section-stacks-etale} 节中，
我们已经看到，可把代数叠理解为 étale 拓扑下如下的群胚叠：其对角由
代数空间表示，并具有光滑覆盖。因此，若存在光滑满射
$U \to \mathcal{S}$，则 $\mathcal{S}$ 是代数叠，特别是满足 fppf
拓扑下的下降。但上面已经看到，$\mathcal{S}$
% P12-ERR-0361
不满足 fppf 拓扑下的下降。

\medskip\noindent
粗略地说，上述论证表明：对基 $B$ 上的群概形 $G$，étale 拓扑下
分类其 étale 局部平凡挠子的叠是代数的，当且仅当 $G$ 在 $B$
上光滑。采用 fppf 拓扑的一个优点是，只需假设
$G \to B$ 平坦且局部有限表示。事实上，fppf 拓扑下的商叠
$[B/G]$ 是代数的，当且仅当 $G \to B$ 平坦且局部有限表示；见
《可表性判据》引理 \ref{criteria-lemma-BG-algebraic}。







\section{具有拟紧平坦覆盖但非代数的层}
% P12-ERR-0362
\label{examples-section-not-algebraic}

\noindent
考虑函子 $F = (\mathbf{P}^1)^\infty$；也就是说，对概形 $T$，
$F(T)$ 是所有 $f = (f_1, f_2, f_3, \ldots)$ 的集合，其中每个
$f_i : T \to \mathbf{P}^1$ 都是概形态射。注意，$\mathbf{P}^1$
对 fpqc 覆盖满足层性质；见《下降》引理
\ref{descent-lemma-fpqc-universal-effective-epimorphisms}。
层的积仍是层，故 $F$ 也对 fpqc 拓扑满足层性质。$F$ 的对角可表：
若 $f : T \to F$ 和 $g : S \to F$ 为态射，则 $T \times_F S$
是概形 $T \times S$ 内各闭子概形
$T \times_{f_i, \mathbf{P}^1, g_i} S$ 的概形论交。考虑群概形
$\text{SL}_2$；它带有一个满光滑仿射态射
$\text{SL}_2 \to \mathbf{P}^1$。其次，考虑
$U = (\text{SL}_2)^\infty$ 及其典范（积）态射 $U \to F$。
注意，$U$ 是仿射概形。我们断言，态射 $U \to F$ 平坦、满且泛开。
确切地说，假设 $f : T \to F$ 为态射。于是
$Z = T \times_F U$ 是各概形
$Z_i = T \times_{f_i, \mathbf{P}^1} \text{SL}_2$ 在 $T$ 上的无穷纤维积。
每个态射 $Z_i \to T$ 都满、光滑且仿射，这蕴含
$$
Z = Z_1 \times_T Z_2 \times_T Z_3 \times_T \ldots
$$
% P12-ERR-0363
是一个在 $Z$ 上平坦且仿射的概形。一个简单的极限论证还表明
$Z \to T$ 是开的。

\medskip\noindent
另一方面，我们断言 $F$ 不是代数空间。事实上，若 $F$ 是代数空间，
则它将是拟紧且分离的代数空间（由我们对 $F$ 上纤维积的描述）。
因此，拟凝聚层的上同调将在某个截断次数以上消失（见《空间的上同调》
命题 \ref{spaces-cohomology-proposition-vanishing} 及其前面的评注）。
但显然，经投影
$$
(\mathbf{P}^1)^\infty \longrightarrow (\mathbf{P}^1)^n
$$
（它有一个截面）拉回 $\mathcal{O}(-2, -2, \ldots, -2)$，便可得到一个
拟凝聚层，其 $n$ 次上同调非零。综上，我们回答了 Anton Geraschenko
在 MathOverflow 上提出的一个问题。

\begin{lemma}
\label{examples-lemma-not-algebraic}
存在一个满足 fpqc 拓扑的层条件并具有可表对角
$\Delta : F \to F \times F$ 的函子
$F : \Sch^{opp} \to \textit{Sets}$；还存在概形 $U$ 及满、平坦、泛开、
拟紧态射 $U \to F$，但 $F$ 不是代数空间。
\end{lemma}

\begin{proof}
见上述讨论。
\end{proof}







\section{étale 拓扑与 Zariski 覆盖及有限 étale 覆盖的比较}
\label{examples-section-etale-zariski-finite-etale}

\noindent
本节给出 B.~Bhatt 为回答 C.~Deninger 的一个问题而找到的例子；
它说明，$\mathbf{C}$ 上代数簇的 étale 拓扑严格细于由开浸入和有限
étale 态射生成的拓扑。这类例子无疑为专家所知，也许可以一直追溯到
Artin 和 Grothendieck\footnote{据轶闻，Grothendieck 最初曾试图用
Zariski 开覆盖和有限 étale 覆盖定义 étale 拓扑，后来 Artin 说服他
放弃这种做法（大概正是依据本文所记录的这类例子）。}。
若读者知道已发表的例子，请发送电子邮件至
\href{mailto:stacks.project@gmail.com}{stacks.project@gmail.com}。

\medskip\noindent
设 $X$ 为 $\mathbf{C}$ 上的光滑连通仿射曲线。选取来自光滑连通曲线
$Y$ 的次数为 $3$ 的有限态射 $g: Y \to X$，使得存在一点
$x \in X$，其纤维 $g^{-1}(x) = \{u, y\}$ 由一个非分歧点 $u$
（次数 1）和一个分歧点 $y$（次数 2）组成。把 $X$ 替换为 $x$
的一个 Zariski 开邻域后，可以假设 $g$ 除了在 $y$ 处以外处处非分歧。
令 $U = Y \setminus \{y\}$。于是
$f=g|_U: U \to X$ 是 étale 满射，且 $f^{-1}(x) = \{u\}$。

\begin{lemma}
\label{examples-lemma-etale-not-dominated-zariski-finite-etale}
étale 覆盖 $U \to X$ 不能由开浸入和有限 étale 态射的任何有限复合
加细。
\end{lemma}

\begin{proof}
反证之，假设存在态射塔
$$
W_m \to W_{m-1} \to \dots \to W_1 \to W_0 = X
$$
其中每个映射 $W_i \to W_{i - 1}$ 或为开浸入，或为有限 étale 态射；
还存在 $W_m \to X$ 沿 $f : U \to X$ 的一个提升 $W_m \to U$，
以及一点 $w \in W_m$，它映到 $u \in U$（因而映到 $x \in X$）。
把每个 $W_k$ 替换为含 $w$ 的像 $w_k \in W_k$ 的连通分支后，
可以假设每个 $W_k$ 都光滑且连通。

\medskip\noindent
对每个 $0 \le i \le m$，考虑纤维积
$Z_i = Y \times_X W_i$。由于 $W_i \to X$ 为 étale 且 $Y$ 光滑，
曲线 $Z_i$ 是若干光滑连通曲线的无交并，而由基变换，映射
$Z_i \to W_i$ 均为次数 $3$ 的有限态射。显然，$Z_i \to W_i$
在 $w_i$ 上的纤维由两点 $z_i, v_i$ 组成，它们在 $Y$ 中分别映到
$y, u$。态射 $Z_i \to W_i$ 在 $v_i$ 处非分歧，而在 $z_i$ 处分歧
次数为 $2$。令 $n_i = \#\pi_0(Z_i)$。由于 $Z_i \to W_i$ 的纤维有
$2$ 个点，可知 $n_i \in \{1, 2\}$。此外，因为转移映射
$Z_i \to Z_{i - 1}$ 是占优的，$n_i$ 随 $i$ 只能增加。显然
$n_0 = 1$。而且必有 $n_m \geq 2$：存在 $X$-态射
$W_m \to U \subset Y$，故基变换 $Z_m \to W_m$ 有截面。因此，
存在最小指标 $i \ge 1$，使得 $n_{i - 1} = 1$ 且 $n_i = 2$。
% P12-ERR-0364
由于 $n_j$ 由 $Z_i$ 是否不可约决定，映射 $W_i \to W_{i - 1}$
不可能是开浸入，所以它必为有限 étale 态射。因为 $n_i = 2$，
光滑曲线 $Z_i$ 有两个连通分支；又因到 $W_i$ 的映射有限，其中一个
分支必同构地映到 $W_i$。因此可写
$Z_i = W_i \sqcup Z'_i$，其中 $Z'_i \to W_i$ 的次数为 $2$。
诱导映射 $W_i \subset Z_i \to Z_{i - 1}$ 是连通曲线之间的占优映射，
而这两条曲线在 $W_{i - 1}$ 上都是有限的。任何这样的映射都必为满射。
所以，可以选取一点 $t \in W_i$，它映到 $z_{i - 1}$。得到离散赋值环的
扩张
$$
% P12-ERR-0365
\mathcal{O}_{W_{i - 1}, w_{i - 1}} \to
\mathcal{O}_{Z_i, z_i} \to
\mathcal{O}_{W_i, z}
$$
其中第一个扩张是分歧的，而复合扩张是非分歧的，因为
$W_i \to W_{i - 1}$ 为 étale；这是不可能的，故得到矛盾。
\end{proof}










\section{层与专化}
\label{examples-section-sheaves}

\noindent
以下固定《拓扑》定义 \ref{topologies-definition-big-etale-site}
所构造的大 étale 位点 $\Sch_\etale$。此外，概形均视为该位点的对象。
回忆，若 $x, x'$ 是概形 $X$ 的点，而
$x \in \overline{\{x'\}}$，则称 $x$ 是 $x'$ 的一个{\it 专化}，
或写成 $x' \leadsto x$。特别地，当 $x = x'$ 时也成立。

\medskip\noindent
考虑由下述规则定义的函子 $F : \Sch_\etale \to \textit{Ab}$：
$$
F(X) = \prod\nolimits_{x \in X}
\prod\nolimits_{x' \in X, x' \leadsto x, x' \not = x}
\mathbf{Z}/2\mathbf{Z}
$$
给定概形 $X$，以 $|X|$ 表示其底层点集。把元素 $a \in F(X)$
视为集合映射
$|X| \times |X| \to \mathbf{Z}/2\mathbf{Z}$，
$(x, x') \mapsto a(x, x')$；当 $x = x'$ 或 $x$ 不是 $x'$ 的专化时，
它取零。给定概形态射 $f : X \to Y$，定义
$$
F(f) : F(Y) \longrightarrow F(X)
$$
如下：对 $b \in F(Y)$，令
$$
F(f)(b)(x, x') =
\left\{
\begin{matrix}
0 & \text{若 }x\text{ 并非专化自 }x' \\
b(f(x), f(x')) & \text{否则。}
\end{matrix}
\right.
$$
注意，这确实定义了 $F(X)$ 的一个元素。我们断言，若
$f : X \to Y$ 与 $g : Y \to Z$ 是可复合态射，则
$F(f) \circ F(g) = F(g \circ f)$。确切地说，令 $c \in F(Z)$，并设
$x' \leadsto x$ 是 $X$ 中各点的一个专化，则
$$
% P12-ERR-0366
F(g \circ f)(x, x') = c(g(f(x)), g(f(x'))) = F(g)(F(f)(c))(x, x')
$$
因为 $f(x') \leadsto f(x)$。（即使 $f(x) = f(x')$，论证也成立。）

\medskip\noindent
令 $G$ 为 $F$ 在 étale 拓扑下的层化。

\medskip\noindent
我们断言，若 $X$ 是概形，而 $x' \leadsto x$ 是满足
$x' \not = x$ 的专化，则 $G(X) \not = 0$。确切地说，取
$a \in F(X)$，使得在把 $a$ 视为函数
$|X| \times |X| \to \mathbf{Z}/2\mathbf{Z}$ 时，它在 $(x, x')$
处非零。令 $\{f_i : U_i \to X\}$ 为 $X$ 的 étale 覆盖。
可以选取一个 $i$ 和一点 $u_i \in U_i$，使 $f_i(u_i) = x$。
因为一般化可沿平坦态射提升（见《态射》引理
\ref{morphisms-lemma-generalizations-lift-flat}），可以找到专化
$u'_i \leadsto u_i$，使 $f_i(u'_i) = x'$。由上述构造可见
$F(f_i)(a) \not = 0$。因此，$a$ 确定 $G(X)$ 中的一个非零元素。

\medskip\noindent
注意，若 $X = \Spec(k)$，其中 $k$ 为域（或更一般地，是所有素理想均为
极大理想的环），则 $F(X) = 0$；而对每个 étale 态射 $U \to X$，
也有 $F(U) = 0$，因为 étale 态射的纤维中的不同点之间没有专化。
因此 $G(X) = 0$。

\medskip\noindent
假设 $X \subset X'$ 是加厚；见《态射的更多内容》定义
\ref{more-morphisms-definition-thickening}。由基变换函子
$U' \mapsto U = U' \times_{X'} X$，$X'$ 上 étale 概形的范畴与
$X$ 上 étale 概形的范畴等价；见《Étale 上同调》定理
\ref{etale-cohomology-theorem-topological-invariance}。
由于在这种情形下总有 $F(U) = F(U')$，所以也有 $G(X) = G(X')$。

\medskip\noindent
作为变体，可以考虑预层 $F_n$：它把概形 $X$ 映到所有映射
$a : |X|^{n + 1} \to \mathbf{Z}/2\mathbf{Z}$ 的集合，其中
$a(x_0, \ldots, x_n)$ 仅在
$x_n \leadsto \ldots \leadsto x_0$ 是专化序列且
$x_n \not = x_{n - 1} \not = \ldots \not = x_0$ 时才可能非零。
令 $G_n$ 为与 $F_n$ 伴随的层。与上面完全相同的方法表明，
$G_n$ 在 $\dim(X) \geq n$ 时非零，而在 $\dim(X) < n$ 时为零。

\begin{lemma}
\label{examples-lemma-sheaf-zero-on-low-dimension}
$\Sch_\etale$ 上存在一个 Abel 群层 $G$，具有如下性质：
\begin{enumerate}
\item 只要 $\dim(X) < n$，就有 $G(X) = 0$；
\item 若 $\dim(X) \geq n$，则 $G(X)$ 非零；且
\item 若 $X \subset X'$ 是加厚，则 $G(X) = G(X')$。
\end{enumerate}
\end{lemma}

\begin{proof}
见上述讨论。
\end{proof}

\begin{remark}
\label{examples-remark-specialization}
下面是若干评注：
\begin{enumerate}
\item 预层 $F$ 和 $F_n$ 是分离预层。
\item 事实上，$F$、$F_n$ 都不是层。
\item 可以证明，$G$、$G_n$
% P12-ERR-0367
实际上是 fppf 拓扑下的层。
\end{enumerate}
若以后需要，我们将证明这些结果。
\end{remark}


\section{层与可构造函数}
\label{examples-section-constructible-functions}

\noindent
以下固定《拓扑》定义 \ref{topologies-definition-big-etale-site}
所构造的大 étale 位点 $\Sch_\etale$。此外，概形均视为该位点的对象。
在本节中，概形 $X$ 的一个{\it 可构造分划}是局部有限的无交并分解
$X = \coprod_{i \in I} X_i$，其中每个 $X_i \subset X$ 都是
$X$ 的局部可构造子集。所谓局部有限，是指对任意拟紧开集
$U \subset X$，只有有限多个 $i \in I$ 满足
$X_i \cap U$ 非空。注意，若 $f : X \to Y$ 为概形态射，而
$Y = \coprod Y_j$ 是可构造分划，则
$X = \coprod f^{-1}(Y_j)$ 是 $X$ 的可构造分划。给定集合 $S$ 和概形
$X$，称映射 $f : |X| \to S$ 为{\it 可构造函数}，若
$X = \coprod_{s \in S} f^{-1}(s)$ 是 $X$ 的可构造分划。
若 $G$ 是（抽象）群，而 $a, b : |X| \to G$ 是可构造函数，则
$ab : |X| \to G, x \mapsto a(x)b(x)$ 也是可构造函数。
理由是，任意两个可构造分划都有一个同时加细二者的第三个分划。

\medskip\noindent
设 $A$ 为任意 Abel 群。对任意概形 $X$，定义
$$
F(X) =
\frac{\{a : |X| \to A \mid a \text{ 是可构造函数}\}}{\{\text{局部常值函数 }|X| \to A\}}
$$
把 $F(X)$ 的一个元素 $a$ 简单地视为一个模加局部常值函数后定义的函数。
给定概形态射 $f : X \to Y$ 和元素 $b \in F(Y)$，定义
$F(f)(b) = b \circ f$。于是 $F$ 是 $\Sch_\etale$ 上的预层。

\medskip\noindent
注意，若 $\{f_i : U_i \to X\}$ 是 fppf 覆盖，而 $a \in F(X)$
满足 $F(f_i)(a) = 0$ 于 $F(U_i)$，则对每个 $i$，$a \circ f_i$
都是局部常值函数。反过来，因为态射 $f_i$ 是开的，这意味着 $a$
是局部常值函数。因此 $a = 0$ 于 $F(X)$。由此可见，$F$ 是分离预层
（对 fppf 拓扑如此，因而更对 étale 拓扑如此）。

\medskip\noindent
令 $G$ 为 $F$ 在 étale 拓扑下的层化。由于 $F$ 分离，而且例如当
$X$ 是离散赋值环的谱时 $F(X) \not = 0$，可见 $G$ 非零。

\medskip\noindent
令 $X = \Spec(k)$，其中 $k$ 为域。于是，$X$ 的任意 étale 覆盖都可由
某个覆盖 $\{\Spec(k') \to \Spec(k)\}$ 控制，其中 $k'/k$
是有限可分域扩张。由于 $F(\Spec(k')) = 0$，可知 $G(X) = 0$。

\medskip\noindent
假设 $X \subset X'$ 是加厚；见《态射的更多内容》定义
\ref{more-morphisms-definition-thickening}。由基变换函子
$U' \mapsto U = U' \times_{X'} X$，$X'$ 上 étale 概形的范畴与
$X$ 上 étale 概形的范畴等价；见《Étale 上同调》定理
\ref{etale-cohomology-theorem-topological-invariance}。
因为在这种情形下 $F(U) = F(U')$，所以也有 $G(X) = G(X')$。

\medskip\noindent
层 $G$ 保持极限；见《空间的极限》定义
\ref{spaces-limits-definition-locally-finite-presentation}。
确切地说，设环 $R$ 写成环的有向余极限 $R = \colim_i R_i$。
令 $X = \Spec(R)$ 且 $X_i = \Spec(R_i)$，于是
$X = \lim_i X_i$。则 $G(X) = \colim_i G(X_i)$。为证明这一点，
先证明 $\Spec(R)$ 的一个可构造分划来自某个 $\Spec(R_i)$ 的一个
% P12-ERR-0368
可构造分划，由此得到 $F$ 的结论。为得到层化的结论，利用如下事实：
任意 étale 环映射 $R \to R'$ 都来自某个 $i$ 的 étale 环映射
$R_i \to R_i'$。细节从略。

\begin{lemma}
\label{examples-lemma-weird-sheaf}
$\Sch_\etale$ 上存在一个 Abel 群层 $G$，具有如下性质：
\begin{enumerate}
\item 对任意域 $k$，都有 $G(\Spec(k)) = 0$；
\item $G$ 保持极限；
\item 若 $X \subset X'$ 是加厚，则 $G(X) = G(X')$；且
\item $G$ 非零。
\end{enumerate}
\end{lemma}

\begin{proof}
见上述讨论。
\end{proof}




\section{光滑-étale 位点并非函子性的}
\label{examples-section-lisse-etale-not-functorial}

\noindent
$X$ 的{\it 光滑-étale} 位点 $X_{lisse,\etale}$ 是 $X$ 上光滑概形的
范畴，赋以（通常的）étale 覆盖；见《叠的上同调》第
\ref{stacks-cohomology-section-lisse-etale} 节。令
$f : X  \to Y$ 为概形态射。有函子
$$
u : Y_{lisse,\etale} \longrightarrow X_{lisse,\etale},\quad
V/Y \longmapsto V \times_Y X
$$
且它连续。因此，得到一对伴随函子
$$
u^s :
\Sh(X_{lisse,\etale})
\longrightarrow
\Sh(Y_{lisse,\etale}),
\quad
u_s :
\Sh(Y_{lisse,\etale})
\longrightarrow
\Sh(X_{lisse,\etale}),
$$
见《位点》第 \ref{sites-section-continuous-functors} 节。
我们断言，一般而言，$u$ 并{\bf 不}定义位点态射；见《位点》定义
\ref{sites-definition-morphism-sites}。换言之，我们断言 $u_s$
一般不是左正合的。注意，可表预层在光滑-étale 位点上是层。因此，
由《位点》引理 \ref{sites-lemma-pullback-representable-sheaf}，有
$u_sh_V = h_{V \times_Y X}$。现在考虑 $Y$ 上光滑概形 $V_1, V_2$
之间的两个态射
$$
\xymatrix{
V_1 \ar[rd] \ar@<1ex>[rr]^a \ar@<-1ex>[rr]_b & & V_2 \ar[ld] \\
& Y
}
$$
若 $u_s$ 左正合，则应有
$$
u_s \text{Equalizer}(h_a, h_b : h_{V_1} \to h_{V_2})
=
\text{Equalizer}(h_{a \times 1}, h_{b \times 1} :
h_{V_1 \times_Y X} \to h_{V_2 \times_Y X})
$$
我们将选取态射 $a, b : V_1 \to V_2$，使得不存在从 $Y$ 上光滑概形
到 $(a = b) \subset V_1$ 的态射，即等式左侧为空层；但基变换到
$X$ 后，等化子非空且在 $X$ 上光滑。一个简单例子是取
$X = \Spec(\mathbf{F}_p)$、$Y = \Spec(\mathbf{Z})$，并令
$V_1 = V_2 = \mathbf{A}^1_\mathbf{Z}$，态射为
$a(x) = x$ 和 $b(x) = x + p$。注意，$a$ 与 $b$ 的等化子是
$\mathbf{A}^1_\mathbf{Z}$ 在 $(p)$ 上的纤维。

\begin{lemma}
\label{examples-lemma-lisse-etale-not-functorial}
光滑-étale 位点并非函子性的，即使对概形态射也如此。
\end{lemma}

\begin{proof}
见上述讨论。
\end{proof}





\section{Noetherian 概形范畴上的层}
\label{examples-section-sheaves-locally-Noetherian}

\noindent
设 $S$ 为局部 Noetherian 概形。像《Artin 公理》第
\ref{artin-section-noetherian} 节一样，考虑包含函子
$$
u : (\textit{Noetherian}/S)_{fppf} \longrightarrow (\Sch/S)_{fppf}
$$
它把 $S$ 上局部 Noetherian 概形的 fppf 位点嵌入 $S$ 的一个大 fppf
位点。正如所引小节解释的，这个函子连续。因此，得到一对伴随函子
$$
u^s :
\Sh((\Sch/S)_{fppf})
\longrightarrow
\Sh((\textit{Noetherian}/S)_{fppf})
$$
和
$$
u_s :
\Sh((\textit{Noetherian}/S)_{fppf})
\longrightarrow
\Sh((\Sch/S)_{fppf})
$$
见《位点》第 \ref{sites-section-continuous-functors} 节。
然而，我们断言，一般而言，$u$ 不定义位点态射；见《位点》定义
\ref{sites-definition-morphism-sites}。换言之，函子 $u_s$
一般不是左正合的。

\medskip\noindent
令 $p$ 为素数，并置 $S = \Spec(\mathbf{F}_p)$。考虑
$(\textit{Noetherian}/S)_{fppf}$ 上如下定义的层的单射
$$
a : \mathcal{F} \longrightarrow \mathcal{G}
$$
对 $S$ 上局部 Noetherian 概形 $U$，定义
$$
\mathcal{G}(U) = \Gamma(U, \mathcal{O}_U)^* =
\Mor_S(U, \mathbf{G}_{m, S})
$$
并取
$$
\mathcal{F}(U) = \{f \in \mathcal{G}(U) \mid \text{fppf 局部地，}f
\text{ 有任意 }p\text{ 次幂根}\}
$$
Noetherian $\mathbf{F}_p$-代数 $A$ 的幂零根理想 $I \subset A$ 是
幂零的；$A$ 中 $1$ 的 $p$ 次幂根形如 $1 + a$，其中 $a \in I$；
% P12-ERR-0369
% P12-ERR-0370
因此，$1$ 的非平凡 $p$ 次幂根中没有一个具有任意 $p$ 次幂根。
由此可知，对任意局部 Noetherian 概形 $U$，$\mathcal{F}(U)$
都是无 $p$-挠的 Abel 群；一些细节从略。于是
$p : \mathcal{F} \to \mathcal{F}$ 是
$(\textit{Noetherian}/S)_{fppf}$ 上 Abel 层的单射。

\medskip\noindent
% P12-ERR-0371
为得到矛盾，假设 $u_s$ 正合。于是
$p : u_s\mathcal{F} \to u_s\mathcal{F}$ 也是单射，并可知对 $S$
上的任意 $V$，$(u_s\mathcal{F})(V)$ 都是无 $p$-挠的 Abel 群。
由于可表预层在 fppf 位点上是层，由《位点》引理
\ref{sites-lemma-pullback-representable-sheaf} 可见，
$u_s\mathcal{G}$ 由 $\mathbf{G}_{m, S}$ 表示。利用
$u_s\mathcal{F} \to u_s\mathcal{G}$ 是单射这一点，对 $S$ 上每个概形
$V$，得到一个无 $p$-挠子群
$$
(u_s\mathcal{F})(V) \subset \Gamma(V, \mathcal{O}_V)^*
$$
它具有如下性质：对 $S$ 上概形的每个态射 $V \to U$，若 $U$
局部 Noetherian，则子群
$$
\mathcal{F}(U) \subset \Gamma(U, \mathcal{O}_U)^*
$$
经限制映射
$\Gamma(U, \mathcal{O}_U)^* \to \Gamma(V, \mathcal{O}_V)^*$
映入子群 $(u_s\mathcal{F})(V)$。

\medskip\noindent
现在如下得到实际的矛盾：令
$k = \bigcup_{n \geq 0} \mathbf{F}_p(t^{1/{p^n}})$，并置
$$
B = k \otimes_{\mathbf{F}_p(t)} k
$$
以及 $V = \Spec(B)$。因为与两个余投影 $k \to B$ 对应，有两个投影态射
$V \to \Spec(k)$；又因为 $\Spec(k)$ 是 Noetherian 的，所以子群
$$
(u_s\mathcal{F})(V) \subset B^*
$$
含有 $k^* \otimes 1$ 和 $1 \otimes k^*$。这与如下事实矛盾：
$$
(t^{1/p} \otimes 1) \cdot (1 \otimes t^{-1/p}) =
t^{1/p} \otimes t^{-1/p}
$$
是 $B$ 的一个非平凡 $p$-挠单位。

\begin{lemma}
\label{examples-lemma-not-a-morphism-of-sites-noetherian-to-all}
取 $S = \Spec(\mathbf{F}_p)$ 时，包含函子
$(\textit{Noetherian}/S)_{fppf} \to (\Sch/S)_{fppf}$
不定义位点态射。
\end{lemma}

\begin{proof}
见上述讨论。
\end{proof}





\section{拟凝聚模的导出推前}
\label{examples-section-derived-push-quasi-coherent}

\noindent
设 $k$ 为特征 $p > 0$ 的域，并令 $S = \Spec(k[x])$。
令 $G = \mathbf{Z}/p\mathbf{Z}$；既可把它视为抽象群，也可视为 $S$
上的常群概形。考虑代数叠 $\mathcal{X} = [S/G]$，其中 $G$
在 $S$ 上平凡作用；见《叠的例子》注
\ref{examples-stacks-remark-X-mod-G-group} 和《可表性判据》引理
\ref{criteria-lemma-BG-algebraic}。考虑结构态射
$$
f : \mathcal{X} \longrightarrow S
$$
这个态射拟紧且拟分离。因此得到函子
$$
Rf_{\QCoh, *} :
D^{+}_\QCoh(\mathcal{O}_\mathcal{X})
\longrightarrow
D^{+}_\QCoh(\mathcal{O}_S),
$$
见《叠的导出范畴》命题
\ref{stacks-perfect-proposition-derived-direct-image-quasi-coherent}。
下面计算 $Rf_{\QCoh, *}\mathcal{O}_\mathcal{X}$。因为
$D_\QCoh(\mathcal{O}_S)$ 等价于 $k[x]$-模的导出范畴（见《概形的
导出范畴》引理 \ref{perfect-lemma-affine-compare-bounded}），
这等价于计算 $R\Gamma(\mathcal{X}, \mathcal{O}_\mathcal{X})$。
为此，可以使用覆盖 $S \to \mathcal{X}$ 及谱序列
$$
% P12-ERR-0372
H^q(S \times_\mathcal{X} \ldots \times_\mathcal{X} S, O)
\Rightarrow H^{p + q}(\mathcal{X}, \mathcal{O}_\mathcal{X})
$$
见《叠的上同调》命题
\ref{stacks-cohomology-proposition-smooth-covering-compute-cohomology}。
注意
$$
S \times_\mathcal{X} \ldots \times_\mathcal{X} S = S \times G^p
$$
是仿射的。因此，复形
$$
k[x] \to \text{Map}(G, k[x]) \to \text{Map}(G^2, k[x]) \to \ldots
$$
计算 $R\Gamma(\mathcal{X}, \mathcal{O}_\mathcal{X})$。这里，对
$\varphi \in \text{Map}(G^{p - 1}, k[x])$，其微分是把
$(g_1, \ldots, g_p)$ 映到下式的映射：
$$
\varphi(g_2, \ldots, g_p) +
\sum\nolimits_{i = 1}^{p - 1}
(-1)^i\varphi(g_1, \ldots, g_i + g_{i + 1}, \ldots, g_p)
+ (-1)^p\varphi(g_1, \ldots, g_{p - 1}).
$$
% P12-ERR-0373
这正是计算 $G$ 在 $k[x]$ 上平凡作用之群上同调的复形
（在此插入将来的参考文献）。循环群 $G$ 在 $k[x]$ 上的上同调，
在每个上同调次数 $\geq 0$ 中恰好都是一份 $k[x]$
（在此插入将来的参考文献）。由此得到
$$
Rf_*\mathcal{O}_\mathcal{X} = \bigoplus\nolimits_{n \geq 0} \mathcal{O}_S[-n]
$$
现在考虑复形
$$
E = \bigoplus\nolimits_{m \geq 0} \mathcal{O}_\mathcal{X}[m]
$$
它是 $D_\QCoh(\mathcal{O}_\mathcal{X})$ 的对象。我们先中断讨论，
给出一个一般结果。

\begin{lemma}
\label{examples-lemma-is-limit}
设 $\mathcal{X}$ 为代数叠。设 $K$ 是 $D(\mathcal{O}_\mathcal{X})$
的对象，其上同调层局部拟凝聚（《叠上的层》定义
\ref{stacks-sheaves-definition-locally-quasi-coherent}），并满足平坦基变换
性质（《叠的上同调》定义
\ref{stacks-cohomology-definition-flat-base-change}）。则在
$D(\mathcal{O}_\mathcal{X})$ 中存在特异三角
$$
K \to
\prod\nolimits_{n \geq 0} \tau_{\geq -n} K \to
\prod\nolimits_{n \geq 0} \tau_{\geq -n} K \to K[1]
$$
换言之，$K$ 是其典范截断的导出极限。
\end{lemma}

\begin{proof}
回忆，我们在 $\mathcal{X}$ 的“大 fppf 位点”
$\mathcal{X}_{fppf}$ 上工作（这是《叠上的层》和《叠的上同调》各章中
关于 $\mathcal{O}_\mathcal{X}$-模层的约定）。令 $\mathcal{B}$ 为
$\mathcal{X}_{fppf}$ 中位于某个仿射概形 $U$ 上方的对象 $x$ 的集合。
综合《叠上的层》引理
\ref{stacks-sheaves-lemma-compare-fppf-etale}、
\ref{stacks-sheaves-lemma-cohomology-restriction}，《下降》引理
\ref{descent-lemma-quasi-coherent-and-flat-base-change}，以及《概形的上同调》
引理 \ref{coherent-lemma-quasi-coherent-affine-cohomology-zero}，可见若
$\mathcal{F}$ 局部拟凝聚且 $x \in \mathcal{B}$，则
$H^p(x, \mathcal{F}) = 0$。现在，取 $d = 0$，结论即由《位点上的
上同调》引理 \ref{sites-cohomology-lemma-is-limit-dimension} 得到。
\end{proof}

\begin{lemma}
\label{examples-lemma-sum-is-product}
设 $\mathcal{X}$ 为代数叠。若 $\mathcal{F}_n$ 是 $\mathcal{X}$
上一族满足平坦基变换性质的局部拟凝聚层，则
$\oplus_n \mathcal{F}_n[n] \to \prod_n \mathcal{F}_n[n]$
是 $D(\mathcal{O}_\mathcal{X})$ 中的同构。
\end{lemma}

\begin{proof}
这是因为，由引理 \ref{examples-lemma-is-limit}，直和同构于直积。
\end{proof}

\noindent
继续上述讨论。因为拟凝聚模是局部拟凝聚的，并满足平坦基变换性质
（《叠上的层》引理 \ref{stacks-sheaves-lemma-quasi-coherent}），得到
$$
E = \prod\nolimits_{m \geq 0} \mathcal{O}_\mathcal{X}[m]
$$
由于上同调与极限交换，可见
$$
Rf_*E = \prod\nolimits_{m \geq 0}
\left(\bigoplus\nolimits_{n \geq 0} \mathcal{O}_S[m - n]\right)
$$
注意，这个复形不是 $D_\QCoh(\mathcal{O}_S)$ 的对象，因为其次数
$0$ 的上同调层是无穷多个 $\mathcal{O}_S$ 副本的直积，甚至不是局部
拟凝聚的 $\mathcal{O}_S$-模。

\begin{lemma}
\label{examples-lemma-push-not-OK}
代数叠的拟紧拟分离态射
$f : \mathcal{X} \to \mathcal{Y}$ 未必诱导函子
$Rf_* : D_\QCoh(\mathcal{O}_\mathcal{X}) \to
D_\QCoh(\mathcal{O}_\mathcal{Y})$。
\end{lemma}

\begin{proof}
见上述讨论。
\end{proof}



\section{大 Abel 范畴}
\label{examples-section-big}

\noindent
本节的目的是给出一个“大”Abel 范畴 $\mathcal{A}$ 及对象 $M, N$ 的例子，
使得扩张的同构类之汇集
% P12-ERR-0374
$\Ext_\mathcal{A}(M, N)$ 不是集合。该例由 Freyd 给出；见
\cite[第 131 页，练习 A]{Freyd}。

\medskip\noindent
如下定义 $\mathcal{A}$。$\mathcal{A}$ 的对象是三元组
$(M, \alpha, f)$，其中 $M$ 为 Abel 群，$\alpha$ 为序数，且
$f : \alpha \to \text{End}(M)$ 为映射。态射
$(M, \alpha, f) \to (M', \alpha', f')$ 由 Abel 群同态
$\varphi : M \to M'$ 给出，并要求对{\it 任意}序数 $\beta$ 都有
$$
\varphi \circ f(\beta) = f'(\beta) \circ \varphi
$$
这里的约定是：若 $\beta$ 不属于 $\alpha$，则置 $f(\beta) = 0$；类似地，
若 $\beta$ 不是 $\alpha'$ 的元素，则把 $f'(\beta)$ 置为零。
我们略去如此定义的范畴为 Abel 范畴的验证。

\medskip\noindent
考虑对象 $Z = (\mathbf{Z}, \emptyset, f)$，即所有算子均为零。
关键观察是，在 $\mathcal{A}$ 中计算的
$\Ext^1_\mathcal{A}(Z, Z)$ 是真类而非集合。确切地说，对每个序数
$\alpha$，可以找到一个 $Z$ 被 $Z$ 扩张得到的对象
$(M, \alpha + 1, f)$，其底层群为
$M = \mathbf{Z} \oplus \mathbf{Z}$，而 $f$ 的值除下式以外总为零：
$$
f(\alpha) =
\left(
\begin{matrix}
0 & 1 \\
0 & 0
\end{matrix}
\right).
$$
显然，这产生一个由扩张的同构类组成的真类。特别地，
$\mathcal{A}$ 的导出范畴的态射汇集是一些真类；见《导出范畴》引理
\ref{derived-lemma-ext-1}。这意味着，在定义三角范畴的 Verdier 商时，
必须格外谨慎。

\begin{lemma}
\label{examples-lemma-big-abelian-category}
存在一个“大”Abel 范畴 $\mathcal{A}$，其 $\Ext$-群为真类。
\end{lemma}

\begin{proof}
见上述讨论。
\end{proof}





\section{弱伴随点与概形论稠密性}
\label{examples-section-weak-ass-dense}

\noindent
设 $k$ 为域。令 $R = k[z, x_i, y_i]/(z^2, zx_iy_i)$，其中 $i$
遍历 $\mathbf{N}$ 的元素。注意，$R = R_0 \oplus M_0$，其中
$R_0 = k[x_i, y_i]$ 是子环，而 $M_0$ 是平方为零的理想，并且作为
$R_0$-模有 $M_0 \cong R_0/(x_iy_i)$。素理想
$\mathfrak p = (z, x_i)$ 弱伴随于作为 $R$-模的 $R$
（《代数》定义 \ref{algebra-definition-weakly-associated}）。
事实上，$R_\mathfrak p$ 中的元素 $z$ 非零，但被
$\mathfrak pR_\mathfrak p$ 零化。另一方面，考虑开子概形
$$
U = \bigcup D(x_i) \subset \Spec(R) = S
$$
我们断言 $U \subset S$ 是概形论稠密的
（《态射》定义 \ref{morphisms-definition-scheme-theoretically-dense}）。
为证明这一点，只需说明 $\mathcal{O}_S \to j_*\mathcal{O}_U$ 是单射，
其中 $j : U \to S$ 是包含态射；见《态射》引理
\ref{morphisms-lemma-characterize-scheme-theoretically-dense}。
译成代数语言，我们必须说明，对所有 $g \in R$，映射
$$
R_g \longrightarrow \prod R_{x_ig}
$$
都是单射。写成 $g = g_0 + m_0$，其中 $g_0 \in R_0$ 且
$m_0 \in M_0$。于是 $R_g = R_{g_0}$（略去细节）。故可假设
$g \in R_0$。还可假设 $g$ 非零。此时
$R_g = (R_0)_g \oplus (M_0)_g$。因为 $R_0$ 是整环，映射
$(R_0)_g \to \prod (R_0)_{x_ig}$ 是单射。若
$g \in (x_iy_i)$，则 $(M_0)_g = 0$，无须证明任何事情。若
$g \not \in (x_iy_i)$，则因 $(x_iy_i)$ 是 $R_0$ 的根理想，
只须说明 $M_0 \to \prod (M_0)_{x_ig}$ 是单射。映射
$R_0 \to M_0 \to (M_0)_{x_n}$ 的核为 $(x_iy_i, y_n)$。由于
$(x_iy_i, y_n)$ 是根理想，若 $g \not \in (x_iy_i, y_n)$，则映射
$R_0 \to M_0 \to (M_0)_{x_ng}$ 的核为 $(x_iy_i, y_n)$。因为对所有
$n \gg 0$ 都有 $g \not \in (x_iy_i, y_n)$，可知该核包含于
$\bigcap_{n \gg 0} (x_iy_i, y_n) = (x_iy_i)$，这正是所需结论。

\medskip\noindent
第二个例由 Ofer Gabber 给出。设 $k$ 为域，并分别令 $R$ 和 $R'$ 为
所有函数 $\mathbf{N} \to k$ 所成的环和所有最终为常值的函数
$\mathbf{N} \to k$ 所成的环。于是 $\Spec(R)$ 和 $\Spec(R')$ 分别是
Stone-{\v C}ech 紧化\footnote{每个 $f \in R$ 都可写成 $ue$，其中
$u$ 是单位而 $e$ 是幂等元。于是《代数》引理
\ref{algebra-lemma-ring-with-only-minimal-primes} 表明 $\Spec(R)$
是 Hausdorff 的。另一方面，可把带离散拓扑的 $\mathbf{N}$ 看成稠密开子集。
给定到 Hausdorff 拟紧拓扑空间 $X$ 的集合映射 $\mathbf{N} \to X$，得到环映射
$\mathcal{C}^0(X; k) \to R$，其中 $\mathcal{C}^0(X; k)$ 是局部常值映射
$X \to k$ 所成的 $k$-代数。这给出
$\Spec(R) \to \Spec(\mathcal{C}^0(X; k)) = X$，从而证明泛性质。}
$\beta\mathbf{N}$ 和单点紧化\footnote{此处可论证 $R'$ 中确实只有一个额外的
极大理想。}$\mathbf{N}^* = \mathbf{N} \cup \{\infty\}$。
所有点都是弱伴随点，因为 $R$ 和 $R'$ 中的所有素理想都是极小素理想。

\begin{lemma}
\label{examples-lemma-example-schematically-dense-missing-weakly-associated-point}
存在约化概形 $X$ 和概形论稠密开子概形 $U \subset X$，使得某个弱伴随点
$x \in X$ 不属于 $U$。
\end{lemma}

\begin{proof}
在第一个例中，由构造即有 $\mathfrak p \not \in U$。在 Gabber 的例中，
概形 $\Spec(R)$ 或 $\Spec(R')$ 是约化的。
\end{proof}



\section{迹不具可加性的例}
\label{examples-section-non-additive}

\noindent
设 $k$ 为域，并令 $R = k[\epsilon]$ 为 $k$ 上的对偶数环。换言之，
$R = k[x]/(x^2)$，且 $\epsilon$ 是 $x$ 在 $R$ 中的同余类。考虑复形的短正合列
$$
\xymatrix{
0 \ar[d] \ar[r] &
R \ar[d]^\epsilon \ar[r]_1 &
R \ar[d] \\
R \ar[r]^1 & R \ar[r] & 0
}
$$
这里各列为复形，第一行置于次数 $0$，第二行置于次数 $1$。分别以
$A^\bullet$、$B^\bullet$ 和 $C^\bullet$ 表示第一个复形（即左列）、
第二个复形和第三个复形。我们断言，图表
\begin{equation}
\label{examples-equation-commutes-up-to-homotopy}
\vcenter{
\xymatrix{
A^\bullet \ar[d]_{1 + \epsilon} \ar[r] &
B^\bullet \ar[r] \ar[d]_1 &
C^\bullet \ar[d]_1 \\
A^\bullet \ar[r] & B^\bullet \ar[r] & C^\bullet
}
}
\end{equation}
在 $K(R)$ 中交换，即它是一个在同伦意义下交换的复形图表。事实上，
右边的方块交换，而左边方块与交换相差同伦 $1 : A^1 \to B^0$。
另一方面，
$$
\text{Tr}_{A^\bullet}(1 + \epsilon) + \text{Tr}_{C^\bullet}(1)
\not = \text{Tr}_{B^\bullet}(1).
$$

\begin{lemma}
\label{examples-lemma-nonadditivity-of-trace}
存在环 $R$、同伦范畴 $K(R)$ 中的特异三角
$(K, L, M, \alpha, \beta, \gamma)$，以及该特异三角的自同态
$(a, b, c)$，使得 $K$、$L$、$M$ 都是完美复形，并且
$\text{Tr}_K(a) + \text{Tr}_M(c) \not = \text{Tr}_L(b)$。
\end{lemma}

\begin{proof}
考虑上述例子。映射 $\gamma : C^\bullet \to A^\bullet[1]$ 在次数 $0$
由乘以 $\epsilon$ 给出；见《导出范畴》定义
\ref{derived-definition-distinguished-triangle}。因此同样有
$$
\xymatrix{
C^\bullet \ar[d] \ar[r]_\gamma & A^\bullet[1] \ar[d] \\
C^\bullet \ar[r]^\gamma & A^\bullet[1]
}
$$
在 $K(R)$ 中交换，因为 $\epsilon(1 + \epsilon) = \epsilon$。
故我们确实得到特异三角的一个态射。
\end{proof}




\section{射影概形结构层的非有限上同调}
\label{examples-section-nonfinite-h0}

\noindent
设 $R$ 为任意非 Noether 环，$I$ 为 $R$ 的非有限生成理想。令
$S = R[x, y]/xyI$，并把它看成分次环，其中 $x$ 和 $y$ 的次数均为 $1$。
则 $X = \text{Proj}(S)$ 是 $R$ 上的射影概形。我们断言
$H^0(X, \mathcal{O}_X)$ 不是有限生成 $R$-模。

\medskip\noindent
按双次数分解环 $S_x$、$S_y$ 和 $S_{xy}$，可写成
$$
S_x = \bigoplus_{\substack{(a, b) \in \mathbf{Z}^2 \\ b \geq 0}} R'_{a, b},
\quad\quad
S_y = \bigoplus_{\substack{(a, b) \in \mathbf{Z}^2 \\ a \geq 0}} R''_{a, b},
\quad\quad
S_{xy} = \bigoplus_{\substack{(a, b) \in \mathbf{Z}^2}} R/I,
$$
其中 $(a, b)$-直和项位于次数 $a + b$，并且
$$
R'_{a, b} =
\left\{
\begin{matrix}
R & b = 0, \\
R/I & b \ne 0,
\end{matrix}
\right.
\quad\quad
R''_{a, b} =
\left\{
\begin{matrix}
R & a = 0, \\
R/I & a \ne 0.
\end{matrix}
\right.
$$
$R$-模 $H^0(X, \mathcal{O}_X)$ 可计算为映射
$$
(S_x)_0 \times (S_y)_0 \to (S_{xy})_0
$$
的核，该映射由 $(F, G) \mapsto F - G$ 定义。这里记号 $A_0$ 表示
$\mathbf{Z}$-分次环 $A$ 的次数 $0$ 分量。由上面对
$S_x, S_y, S_{xy}$ 的显式描述，该核同构于由满足
$F - G \in I$ 的偶对 $(F, G)$ 组成的子模 $M \subset R^2$。由于
$R$-模映射 $M \to I$，$(F, G) \mapsto F - G$ 是满射，
$M$ 不是有限生成 $R$-模。

\medskip\noindent
满射 $R[x, y] \to S$ 对应于闭浸入
$i : X \to \mathbf{P}^1_R$。于是
$\mathcal{F} = i_*\mathcal{O}_X$ 是有限型拟凝聚
$\mathcal{O}_{\mathbf{P}^1_R}$-模，其 $H^0$ 等于上面算得的非有限生成模
$M$。

\begin{lemma}
\label{examples-lemma-nonfinite-h0}
非有限的 $H^0$。
\begin{enumerate}
\item 存在环 $R$ 和 $R$ 上的射影概形 $X$，使得
$H^0(X, \mathcal{O}_X)$ 不是有限 $R$-模。
\item 存在环 $R$ 和 $\mathbf{P}^1_R$ 上的有限型拟凝聚模
$\mathcal{F}$，使得 $H^0(\mathbf{P}^1_R, \mathcal{F})$
不是有限 $R$-模。
\end{enumerate}
\end{lemma}

\begin{proof}
见上述讨论。
\end{proof}









\section{射影性不在基底上局部}
\label{examples-section-non-descending-property-projective}

\noindent
在关于下降的章中，我们已经看到，许多态射性质在基底上局部，甚至在
fpqc 拓扑中也是如此。见《下降》的章节
\ref{descent-section-descending-properties-morphisms}、
\ref{descent-section-descending-properties-morphisms-fpqc} 和
\ref{descent-section-descending-properties-morphisms-fppf}。
对态射的射影性而言，情形并非如此。

\begin{lemma}
\label{examples-lemma-non-descending-property-projective}
性质
\begin{enumerate}
\item[] $\mathcal{P}(f) =$``$f$ 是射影的''，以及
\item[] $\mathcal{P}(f) =$``$f$ 是拟射影的''
\end{enumerate}
在基底上都不是 Zariski 局部的。因此，它们更不是在基底上 fpqc 局部的。
\end{lemma}

\begin{proof}
依照 Hironaka \cite[例 B.3.4.1]{H}，我们构造光滑复三维簇之间的真态射
$f:V_Y\to Y$；它在 Zariski 局部是射影的，但并非射影的。由于 $f$
是真而非射影的，它也不是拟射影的。

\medskip\noindent
令 $Y$ 为复数域 ${\mathbf C}$ 上的射影三维空间。取 $Y$ 中的光滑二次曲线
$C$ 和 $D$，使得闭子概形 $C\cap D$ 约化，且由两个复点 $P$ 和 $Q$
组成。（例如，取
$C=\{ [x,y,z,w]: xy=z^2, w=0\}$，$D=\{ [x,y,z,w]:
xy=w^2, z=0\}$，$P=[1,0,0,0]$，且
$Q=[0,1,0,0]$。）
在 $Y-Q$ 上，先沿曲线 $C$ 吹起，再沿曲线 $D$ 的严格变换吹起
（《除子》定义 \ref{divisors-definition-strict-transform}）。在 $Y-P$ 上，
先沿曲线 $D$ 吹起，再沿曲线 $C$ 的严格变换吹起。在 $Y-P-Q$ 上，
如此构造的两个簇典范同构，故可把它们沿 $Y-P-Q$ 粘合。所得
$V_Y$ 是 ${\mathbf C}$ 上的光滑真三维簇。态射 $f:V_Y\to Y$ 是真的，
并且在 Zariski 局部是射影的（因为它在 $Y-P$ 和 $Y-Q$ 上都是吹起）；
这由《除子》引理 \ref{divisors-lemma-blowing-up-projective} 得到。
我们将证明 $V_Y$ 在 ${\mathbf C}$ 上不是射影的。这将推出 $f$ 不是射影的。

\medskip\noindent
为此，令 $L$ 为 $C-P-Q$ 的一个复点在 $V_Y$ 中的逆像，令 $M$ 为
$D-P-Q$ 的一个复点的逆像。则 $L$ 和 $M$ 都同构于射影直线
${\mathbf P}^1_{{\mathbf C}}$。接着，在 $V_Y$ 中令 $E$ 为
$C\cup D\subset Y$ 在 $V_Y$ 中的逆像；于是
$E\rightarrow C\cup D$ 是真态射，在
$(C\cup D)-\{P,Q\}$ 上的纤维同构于 ${\mathbf P}^1$。点 $P$ 在 $E$
中的逆像是两条直线 $L_0$ 和 $M_0$ 的并，并且在 $E$ 上有循环的有理等价
$L\sim L_0+M_0$ 和 $M\sim M_0$（这里利用 $C$ 和 $D$ 同构于
${\mathbf P}^1$）。注意，这种不对称性来自我们吹起两条曲线的次序。
在 $Q$ 附近，发生相反的情形。因此，$Q$ 的逆像是两条直线
$L_0'$ 和 $M_0'$ 的并，并且在 $E$ 上有有理等价
$L\sim L_0'$ 和 $M\sim L_0'+M_0'$。结合这些等价关系，得到在 $E$
上 $L_0+M_0'\sim 0$，因而在 $V_Y$ 上也如此。若 $V_Y$ 在
${\mathbf C}$ 上射影，则它具有丰富线丛 $H$，而该线丛在 $V_Y$ 的所有曲线
上的次数都 $> 0$。特别地，$H$ 在 $L_0+M_0'$ 上的次数为正；但在线性丛的
次数在域上真概形的模有理等价的一维循环上良定义，这给出矛盾
（《Chow 同调》引理
\ref{chow-lemma-proper-pushforward-rational-equivalence}
和引理 \ref{chow-lemma-factors}）。
所以 $V_Y$ 在 ${\mathbf C}$ 上不是射影的。
\end{proof}

\noindent
换一种术语说，Hironaka 的三维簇 $V_Y$ 是 $Y$ 沿约化子概形
$C\cup D$ 的吹起 $Y'$ 的一个小解消；这里 $Y'$ 有两个结点奇点。
若把 $Z$ 定义为先沿 $C$ 吹起 $Y$，再沿 $D$ 的严格变换吹起所得的概形，
则 $Z$ 是光滑射影三维簇，而非射影三维簇 $V_Y$ 与 $Z$ 的差别是在
$Y-P$ 上做一次 ``flop''。


\section{射影概形的非有效下降数据}
\label{examples-section-non-effective-descent-projective}

\noindent
在关于下降的章中，我们已看到，相对于 fpqc 态射的概形下降数据对于若干类
态射是有效的。特别地，仿射态射以及更一般的拟仿射态射满足 fpqc 覆盖的下降
（《下降》引理 \ref{descent-lemma-quasi-affine}）。对射影态射而言并非如此。

\begin{lemma}
\label{examples-lemma-non-effective-descent-projective}
存在概形的 étale 覆盖 $X\to S$，以及相对于 $X\to S$ 的下降数据
$(V/X,\varphi)$，使得 $V\to X$ 是射影的，但该下降数据在概形范畴中
不是有效的。
\end{lemma}

\begin{proof}
我们仿照 Hironaka 的一个例子：一个不是概形的三维光滑分离复代数空间
\cite[例 B.3.4.2]{H}。

\medskip\noindent
考虑群 $G = \mathbf{Z}/2 = \{1, g\}$ 在复数域上的射影三维空间
$\mathbf{P}^3$ 上的作用：
$$
g[x,y,z,w] = [y,x,w,z].
$$
除 ${\mathbf P}^3$ 中两条不交直线
$L_1=\{ [x,x,z,z]\}$ 和 $L_2=\{ [x,-x,z,-z]\}$ 外，该作用是自由的。
令 $Y={\mathbf P}^3-(L_1\cup L_2)$。存在 ${\mathbf C}$ 上的光滑拟射影概形
$S=Y/G$，使得 $Y\to S$ 是 $G$-挠子（《群胚》定义
\ref{groupoids-definition-principal-homogeneous-space}）。
具体地，可以把 $S$ 定义为开子集 $Y$ 在 ${\mathbf P}^3$ 的态射
\begin{align*}
{\mathbf P}^3 & \to \text{Proj } {\mathbf C}[x,y,z,w]^G\\
   & = \text{Proj } {\mathbf C}[u_0,u_1,v_0,v_1,v_2]/(v_0v_1=v_2^2),
\end{align*}
下的像，其中 $u_0=x+y$、$u_1=z+w$、$v_0=(x-y)^2$、
$v_1=(z-w)^2$ 且 $v_2=(x-y)(z-w)$；环的分次规定为
$u_0,u_1$ 的次数为 1，而 $v_0,v_1,v_2$ 的次数为 2。

\medskip\noindent
令 $C=\{ [x,y,z,w]: xy=z^2, w=0\}$，并令
$D=\{ [x,y,z,w]: xy=w^2, z=0\}$。它们是 ${\mathbf P}^3$ 中的光滑二次
曲线，包含于 $G$-不变开子集 $Y$，且 $g(C)=D$。此外，$C\cap D$
由两点 $P:=[1,0,0,0]$ 和 $Q:=[0,1,0,0]$ 组成，这两个点在 $G$
的作用下彼此交换。

\medskip\noindent
令 $V_Y\to Y$ 为如下概形：在 $Y-P$ 上，先沿 $D$ 吹起，再沿 $C$
的严格变换吹起；在 $Y-Q$ 上，先沿 $C$ 吹起，再沿 $D$ 的严格变换吹起。
（这与引理 \ref{examples-lemma-non-descending-property-projective} 的证明中的构造相同，
不同之处是，这里的 $Y$ 表示 ${\mathbf P}^3$ 的开子集，而不是整个
${\mathbf P}^3$。）于是 $G$ 在 $Y$ 上的作用提升为 $G$ 在 $V_Y$ 上的作用，
该作用交换 $Y-P$ 和 $Y-Q$ 的逆像。$G$ 在 $V_Y$ 上的这个作用给出
$V_Y$ 上相对于 $G$-挠子 $Y\to S$ 的下降数据
$(V_Y/Y,\varphi_Y)$。如引理
\ref{examples-lemma-non-descending-property-projective} 的证明所示，态射
$V_Y\to Y$ 是真的，但不是射影的。

\medskip\noindent
令 $X$ 为开子集 $Y-P$ 和 $Y-Q$ 的不交并；于是有满 étale 态射
$X\to Y\to S$。令 $V$ 为 $V_Y\to Y$ 到 $X$ 的拉回；则态射
$V\to X$ 是射影的，因为 $V_Y\to Y$ 在开子集 $Y-P$ 和 $Y-Q$
的每一个上都是吹起。此外，下降数据 $(V_Y/Y,\varphi_Y)$ 拉回为相对于
étale 覆盖 $X\to S$ 的下降数据 $(V/X,\varphi)$。

\medskip\noindent
假设该下降数据在概形范畴中有效。也就是说，存在概形 $U\to S$，其拉回连同
下降数据就是态射 $V\to X$。于是 $U$ 将是 $V_Y$ 对其 $G$-作用的商。
$$
\xymatrix{
V \ar[r]\ar[d]&  X\ar[d] \\
V_Y \ar[r]\ar[d]& Y\ar[d] \\
U \ar[r]& S
}
$$

\medskip\noindent
令 $E$ 为 $C\cup D\subset Y$ 在 $V_Y$ 中的逆像；于是
$E\rightarrow C\cup D$ 是真态射，并且在 $(C\cup D)-\{P,Q\}$
上的纤维同构于 ${\mathbf P}^1$。$P$ 在 $E$ 中的逆像是两条直线
$L_0$ 和 $M_0$ 的并。因此，$Q=g(P)$ 在 $E$ 中的逆像是两条直线
$L_0'=g(M_0)$ 和 $M_0'=g(L_0)$ 的并。如引理
\ref{examples-lemma-non-descending-property-projective} 的证明所示，在 $E$ 上有
有理等价 $L_0+M_0'=L_0+g(L_0)\sim 0$。

\medskip\noindent
由闭子概形的下降，存在曲线 $L_1\subset U$（同构于 ${\mathbf P}^1$），
其在 $V_Y$ 中的逆像是 $L_0\cup g(L_0)$。（使用《下降》引理
\ref{descent-lemma-affine}，并注意闭浸入是仿射态射。）
令 $R$ 为 $L_1$ 的一个复点。由于假设 $U$ 是概形，可在局部环
% P12-ERR-0375
$O_{U,R}$ 中选择函数 $f$，它在 $R$ 处消失，但不在整条曲线 $L_1$
上消失。令 $D_{\text{loc}}$ 为 $\Spec O_{U,R}$ 中闭子集
$\{f = 0\}$ 的一个不可约分支；则 $D_{\text{loc}}$ 的余维为 1。
$D_{\text{loc}}$ 在 $U$ 中的闭包是 $U$ 中的不可约除子 $D_U$，它包含点
$R$，但不包含整条曲线 $L_1$。$D_U$ 在 $V_Y$ 中的逆像是有效除子 $D$，
它与 $L_0\cup g(L_0)$ 相交，但不包含曲线 $L_0$ 或 $g(L_0)$ 中的任一条。

\medskip\noindent
由于复三维簇 $V_Y$ 光滑，
% P12-ERR-0376
$O(D)$ 是 $V_Y$ 上的线丛。这里我们用到了正则局部环是因子环，换言之是
UFD；见《代数进阶》引理 \ref{more-algebra-lemma-regular-local-UFD}。
$O(D)$ 到真曲面 $E\subset V_Y$ 的限制是线丛；根据我们关于 $D$ 的信息，
它在一维循环 $L_0+g(L_0)$ 上的次数为正。由于在 $E$ 上
$L_0+g(L_0)\sim 0$，这与线丛的次数在域上真概形的模有理等价的一维循环上
良定义相矛盾（《Chow 同调》引理
\ref{chow-lemma-proper-pushforward-rational-equivalence}
和引理 \ref{chow-lemma-factors}）。因此，下降数据 $(V/X,\varphi)$
事实上不是有效的；也就是说，$U$ 不作为概形存在。
\end{proof}

\noindent
在这个例中，该下降数据在代数空间范畴中{\it 是}有效的。更确切地说，
$U$ 作为 ${\mathbf C}$ 上的三维光滑分离代数空间存在；例如，这由
《代数空间》引理 \ref{spaces-lemma-quotient} 得到。Hironaka 的三维空间
$U$ 是光滑拟射影三维簇 $S$ 沿不可约结点曲线 $(C\cup D)/G$ 的吹起
$S'$ 的一个小解消；三维簇 $S'$ 有一个结点奇点。$S'$ 的另一个小解消
（与 $U$ 相差一次 ``flop''）仍是一个不是概形的代数空间。






\section{全空间不是概形的曲线族}
\label{examples-section-family-of-curves}

\noindent
在《Quot》章节 \ref{quot-section-curves} 中，我们把概形 $S$ 上的曲线族
定义为从代数空间 $X$ 到 $S$ 的相对维数 $\leq 1$ 的真、平坦、有限表现态射。
若 $S$ 是完备 Noether 局部环的谱，则 $X$ 是概形；见《空间态射进阶》引理
\ref{spaces-more-morphisms-lemma-projective-over-complete}。
本节说明，这在一般情形下并不成立。

\medskip\noindent
设 $k$ 为域。我们从真平坦态射
$$
Y \longrightarrow \mathbf{A}^1_k
$$
以及位于 $0 \in \mathbf{A}^1_k(k)$ 上方的点 $y \in Y(k)$ 出发，
并要求它们满足如下性质：
\begin{enumerate}
\item 纤维 $Y_0$ 是 $k$ 上的光滑几何不可约曲线；
\item 对任意支配 $\mathbf{A}^1_k$ 的真闭子概形 $T \subset Y$，交
$T \cap Y_0$ 都包含至少一个不同于 $y$ 的点。
\end{enumerate}
给定这样的曲面，我们如下构造所需例子：
$$
\xymatrix{
Y \ar[rd] & Z \ar[d] \ar[l] \ar[r] & X \ar[ld] \\
& \mathbf{A}^1_k
}
$$
这里 $Z \to Y$ 是 $Y$ 在 $y$ 处的吹起。令 $E \subset Z$ 为例外除子，
并令 $C \subset Z$ 为 $Y_0$ 的严格变换。概形论意义下有
$Z_0 = E \cup C$（为看出这一点，可用 $Y$ 在 $y$ 处光滑以及
$Y \to \mathbf{A}^1_k$ 在 $y$ 处光滑）。由 Artin 的结果
（\cite{ArtinII}；使用《半稳定约化》引理
\ref{models-lemma-properties-form} 说明 $C$ 的法丛是负的），
可在 $Z$ 中吹下曲线 $C$，得到图中的代数空间 $X$。令
$x \in X(k)$ 为 $C$ 的像。

\medskip\noindent
我们断言 $X$ 不是概形。事实上，若它是概形，则 $x$ 有仿射开邻域
$U \subset X$。置 $T = X \setminus U$。则 $T$ 支配
$\mathbf{A}^1_k$（因为 $X \to \mathbf{A}^1_k$ 的纤维是真且维数为 $1$ 的，而
$U \to \mathbf{A}^1_k$ 的纤维是仿射的，因而二者不同）。令
$T' \subset Z$ 为同构地映到 $T$ 的闭子概形（因为 $x \not \in T$）。
于是
% P12-ERR-0377
$T'$ 在 $X$ 中的像与上面的条件 (2) 矛盾（因为 $T' \cap Z_0$
包含于吹起 $Z \to Y$ 的例外除子 $E$）。

\medskip\noindent
为结束讨论，还需构造这样的 $Y$。我们假设 $k$ 的特征不是 $3$。
写成 $\mathbf{A}^1_k = \Spec(k[t])$，并在 $\mathbf{P}^2_{k[t]}$ 中取
$$
Y \quad : \quad T_0^3 + T_1^3 + T_2^3 - tT_0T_1T_2 = 0
$$
它在 $t = 0$ 处的纤维是光滑射影亏格为 $1$ 的曲线。在仿射块
% P12-ERR-0378
$V_+(T_0)$ 上，得到仿射方程
$$
1 + x^3 + y^3 - txy = 0
$$
它定义 $k$ 上的光滑曲面。由对称性，其他仿射块上也成立，故 $Y$ 是光滑
曲面。最后，从该仿射方程还可看出分式域是 $k(x, y)$，因而 $Y$ 是有理曲面。
有理曲面的 Picard 群有限生成
% P12-ERR-0379
（在此插入将来的参考文献）。因此，为选出具有性质 (2) 的
$y \in Y_0(k)$，只需选取 $y$ 使得
\begin{equation}
\label{examples-equation-three}
\mathcal{O}_{Y_0}(ny) \not \in \Im(\Pic(Y) \to \Pic(Y_0))
\text{ 对所有 }n > 0
\end{equation}
事实上，与 (2) 矛盾的 $T$ 的所有 $1$-维不可约分支之和将给出一个有效
Cartier 除子，
% P12-ERR-0380
其交 $Y_0$ 为某个 $n \geq 1$ 的除子 $ny$，由此可知
$\mathcal{O}_{Y_0}(ny)$ 位于限制映射的像中。注意，因为 $Y_0$ 的亏格
$\geq 1$，映射
$$
Y_0(k) \to \Pic(Y_0),\quad y \mapsto \mathcal{O}_{Y_0}(y)
$$
是单射。现在，若 $k$ 是不可数代数闭域，则由 $\Pic(Y)$ 的可数性和刚才的
观察，可以找到满足 (\ref{examples-equation-three})、从而满足 (2) 的点
$y \in Y_0(k)$。

\begin{lemma}
\label{examples-lemma-family-of-curves-not-scheme}
存在域 $k$ 和曲线族 $X \to \mathbf{A}^1_k$，使得 $X$ 不是概形。
\end{lemma}

\begin{proof}
见上述讨论。
\end{proof}






\section{导出基变换}
\label{examples-section-derived-base-change}

\noindent
设 $R \to R'$ 为环映射。在《代数进阶》章节
\ref{more-algebra-section-derived-base-change} 中，我们构造了导出基变换函子
$- \otimes_R^\mathbf{L} R' : D(R) \to D(R')$。接着，设
$R \to A$ 为第二个环映射。图示为
$$
\xymatrix{
	A \ar[r] & A \otimes_R R' \ar@{=}[r] & A' \\
R \ar[u] \ar[r] & R' \ar[u] \ar[ur]
}
$$
给定 $A$-模 $M$，张量积 $M \otimes_R R'$ 是
$A \otimes_R R'$-模，即 $A'$-模。对环映射 $A \to A'$，有导出函子
$$
- \otimes_A^\mathbf{L} A' : D(A) \longrightarrow D(A')
$$
但一般而言，该函子与 $- \otimes_R^\mathbf{L} R'$ 不一致。更确切地说，
对 $K \in D(A)$，在《代数进阶》方程
(\ref{more-algebra-equation-comparison-map}) 中构造的 $D(R')$ 中的典范映射
$$
K \otimes_R^{\mathbf{L}} R' \longrightarrow
K \otimes_A^{\mathbf{L}} A'
$$
一般不是同构。因此可以问，是否存在一个“导出基变换函子”
$T : D(A) \to D(A')$，即一个使得 $T(K)$ 在 $D(R')$ 中映到
$K \otimes_R^\mathbf{L} R'$ 的函子。本节将说明，一般而言它不存在。

\medskip\noindent
设 $k$ 为域。置 $R = k[x, y]$、$R' = R/(xy)$ 且 $A = R/(x^2)$。
$D(R')$ 中的对象 $A \otimes_R^\mathbf{L} R'$ 由
$$
x^2 : R' \longrightarrow R'
$$
表示，并且 $H^0(A \otimes_R^\mathbf{L} R') = A \otimes_R R'$。
我们断言，不存在 $D(A \otimes_R R')$ 的对象 $E$，使它在 $D(R')$
中映到 $A \otimes_R^\mathbf{L} R'$。事实上，对这样的 $E$，模
$H^0(E)$ 将是自由的，因而 $E$ 将分解成
$H^0(E)[0] \oplus H^{-1}(E)[1]$。但容易看出，在 $D(R')$ 中，
$A \otimes_R^\mathbf{L} R'$ 不同构于其各上同调群的直和。

\begin{lemma}
\label{examples-lemma-no-derived-base-change}
设 $R \to R'$ 和 $R \to A$ 为环映射。一般而言，不存在三角范畴的函子
$T : D(A) \to D(A \otimes_R R')$，使得 $A$-模 $M$ 给出
$D(A \otimes_R R')$ 中的对象 $T(M)$，且该对象在映射
$D(A \otimes_R R') \to D(R')$ 下映到
$M \otimes_R^\mathbf{L} R'$。
\end{lemma}

\begin{proof}
见上述讨论。
\end{proof}




\section{一个有趣的紧对象}
\label{examples-section-interesting-compact}


\noindent
设 $R$ 为环，并设 $(A, \text{d})$ 为微分分次 $R$-代数。若 $A = R$，
则我们知道 $D(A, \text{d}) = D(R)$ 的每个紧对象都由有限投射模的有限复形
表示。换言之，紧对象是完美的；见《代数进阶》命题
\ref{more-algebra-proposition-perfect-is-compact}。用微分分次模的语言，
相应问题是：“$D(A, \text{d})$ 的每个紧对象是否都可由一个有限且分次投射的
微分分次 $A$-模 $P$ 表示？”

\medskip\noindent
对一般的微分分次代数，这并不成立。事实上，设 $k$ 为特征 $2$ 的域
（这样便不必顾虑符号）。令 $A = k[x, y]/(y^2)$，其中
\begin{enumerate}
\item $x$ 的次数为 $0$；
\item $y$ 的次数为 $-1$；
\item $\text{d}(x) = 0$；并且
\item $\text{d}(y) = x^2 + x$。
\end{enumerate}
于是 $x : A \to A$ 是 $K(A, \text{d})$ 中的投影算子。因此
$$
A = \Ker(x) \oplus \Im(1 - x)
$$
在 $K(A, \text{d})$ 中成立；见《微分分次代数》引理
\ref{dga-lemma-homotopy-direct-sums} 和《导出范畴》引理
\ref{derived-lemma-projectors-have-images-triangulated}。
显然，$A$ 是 $D(A, \text{d})$ 的紧对象。由《导出范畴》引理
\ref{derived-lemma-compact-objects-subcategory}，
$\Ker(x)$ 也是 $D(A, \text{d})$ 的紧对象。

\medskip\noindent
接着，假设 $M$ 是表示 $\Ker(x)$ 的微分分次（右）$A$-模，并假设
$M$ 作为分次 $A$-模是有限且投射的。由于
$k[x, y]/(y^2)$ 上的每个有限分次投射模都是分次自由的，可知 $M$ 作为分次
$k[x, y]/(y^2)$-模是有限自由的（即忘掉微分以后）。置
$N = M/M(x^2 + x)$。考虑正合列
$$
0 \to M \xrightarrow{x^2 + x} M \to N \to 0
$$
由于 $x^2 + x$ 的次数为 $0$，位于 $A$ 的中心，并且
$\text{d}(x^2 + x) = 0$，可知这是微分分次 $A$-模的短正合列。此外，
由 $\text{d}(y) = x^2 + x$ 可知 $N$ 上的微分是线性的。映射
$$
H^{-1}(N) \to H^0(M)
\quad\text{且}\quad
H^0(M) \to H^0(N)
$$
是同构，因为 $M \cong \Ker(x)$ 在 $D(A, \text{d})$ 中成立，所以
$H^*(M) = H^0(M) = k$。计算边界映射可知，作为分次模有
$H^*(N) = k[x, y]/(x, y^2)$；我们略去细节。因为 $N$ 是自由
$k[x, y]/(y^2, x^2 + x)$-模，存在与微分相容的解消
$$
\ldots \to N[2] \xrightarrow{y} N[1] \xrightarrow{y} N \to N/Ny \to 0
$$
由于 $N$ 有界且 $H^0(N) = k[x,y]/(x, y^2)$，由《同调》引理
\ref{homology-lemma-first-quadrant-ss} 可知 $H^0(N/Ny) = k[x]/(x)$。
但 $N/Ny$ 是自由 $k[x]/(x^2 + x) = k \times k$-模的有限复形，
故其上同调的维数必为偶数，矛盾。

\begin{lemma}
\label{examples-lemma-no-good-representatif-compact-object}
存在微分分次代数 $(A, \text{d})$ 和 $D(A, \text{d})$ 的紧对象 $E$，
使得 $E$ 不能由有限且分次投射的微分分次 $A$-模表示。
\end{lemma}

\begin{proof}
见上述讨论。
\end{proof}








\section{两个微分分次范畴}
\label{examples-section-nongraded-differential-graded}

\noindent
本节构造两个满足《微分分次代数》情形 \ref{dga-situation-ABC} 中公理
(A)、(B) 和 (C) 的微分分次范畴，而它们的对象并不带
$\mathbf{Z}$-分次。

\medskip\noindent
{\bf 例 I。}设 $X$ 为拓扑空间。以
% P12-ERR-0381
$\underline{\mathbf{Z}}$ 表示取值为 $\mathbf{Z}$ 的常值层。设 $A$ 为
$\underline{\mathbf{Z}}$-挠子。在此情形下，我们称 Abel 群层
$\mathcal{F}$ 是{\it $A$-分次的}，若对给定的局部截面
$a \in A(U)$，存在投影算子
$p_a : \mathcal{F}|_U \to \mathcal{F}|_U$，使得每当存在局部同构
$\underline{\mathbf{Z}}|_U \to A|_U$ 时，都有
$\mathcal{F}|_U = \bigoplus_{n \in \mathbf{Z}} p_n(\mathcal{F})$。
换一种说法，Abel 层 $\mathcal{F}$ 在 $X$ 上局部具有
$\mathbf{Z}$-分次，但在交叠处，不同分次选择相差由挠子 $A$ 的转移函数
给出的次数平移。
如果 $\mathcal{F}$ 是 $A$-分次 Abel 层，而
$\text{d} : \mathcal{F} \to \mathcal{F}$ 是满足
$\text{d}^2 = 0$ 的微分，并且对 $A$ 的每个局部截面 $a$ 都有
$p_{a + 1} \circ \text{d} = \text{d} \circ p_a$，则称偶对
$(\mathcal{F}, \text{d})$ 是{\it Abel 层的 $A$-分次复形}。
换言之，$\text{d}(p_a(\mathcal{F}))$ 包含于
$p_{a + 1}(\mathcal{F})$。

\medskip\noindent
接着，考虑范畴 $\mathcal{A}$，其中
\begin{enumerate}
\item 对象是 Abel 层的 $A$-分次复形；并且
\item 对对象 $(\mathcal{F}, \text{d})$、$(\mathcal{G}, \text{d})$，置
$$
\Hom_\mathcal{A}((\mathcal{F}, \text{d}), (\mathcal{G}, \text{d})) =
\bigoplus \Hom^n(\mathcal{F}, \mathcal{G})
$$
其中 $\Hom^n(\mathcal{F}, \mathcal{G})$ 是满足下述条件的 Abel 层映射
$f$ 所成的群：对 $A$ 的所有局部截面 $a$，都有
$f(p_a(\mathcal{F})) \subset p_{a + n}(\mathcal{G})$。微分取为
$\text{d}(f) = \text{d} \circ f - (-1)^n f \circ \text{d}$；见
《微分分次代数》例 \ref{dga-example-category-complexes}。
\end{enumerate}
我们略去验证：这确实是满足 (A)、(B) 和 (C) 的微分分次范畴。所有性质都可
在 $X$ 上局部验证；在那里，我们恰好恢复 Abel 层复形所成的微分分次范畴。
因而得到三角范畴 $K(\mathcal{A})$。

\medskip\noindent
$X$ 的扭曲导出范畴。注意，给定 $\mathcal{A}$ 的对象
$(\mathcal{F}, \text{d})$，存在良定义的 $A$-分次上同调层
$H(\mathcal{F}, \text{d})$。因此，$K(\mathcal{A})$ 中的拟同构之含义是
清楚的。将拟同构取逆，就得到{\it $A$-分次层复形的导出范畴
$D(\mathcal{A})$}。若 $A$ 是平凡挠子，则 $D(\mathcal{A})$ 等于
$D(X)$；但对
% P12-ERR-0382
非零挠子，会得到 $X$ 的一种{\it 扭曲}导出范畴。

\medskip\noindent
{\bf 例 II。}设 $C$ 为完美域 $k$ 上的光滑曲线，且其特征为 $2$。
则 $\Omega_{C/k}$ 带有典范平方根。也就是说，可以写成
$\Omega_{C/k} = \mathcal{L}^{\otimes 2}$，使得对 $C$ 上每个局部函数
$f$，截面 $\text{d}(f)$ 都等于某个 $\mathcal{L}$ 的局部截面 $s$
的 $s^{\otimes 2}$。其“理由”是
$$
\text{d}(a_0 + a_1t + \ldots +a_dt^d) =
(\sum\nolimits_{i\text{ 为奇数}} a_i^{1/2} t^{(i - 1)/2})^2\text{d}t
$$
% P12-ERR-0383
（在此插入将来的参考文献）。特别地，这确定了典范联络
$$
\nabla_{can} :
\Omega_{C/k}
\longrightarrow
\Omega_{C/k} \otimes_{\mathcal{O}_C} \Omega_{C/k}
$$
其 $2$-曲率为零（即唯一使平方的导数为零的联络）。注意，带联络的向量丛范畴
是张量范畴，故对所有 $n \in \mathbf{Z}$，我们还在可逆层
$\Omega_{C/k}^{\otimes n}$ 上得到典范联络 $\nabla_{can}$。

\medskip\noindent
令 $\mathcal{A}$ 为如下范畴：
\begin{enumerate}
\item 对象是偶对 $(\mathcal{F}, \nabla)$，其中 $\mathcal{F}$ 是带有联络
$$
\nabla :
\mathcal{F}
\longrightarrow
\mathcal{F} \otimes_{\mathcal{O}_C} \Omega_{C/k}
$$
的有限局部自由层，并且该联络的 $2$-曲率为零；并且
\item $(\mathcal{F}, \nabla_\mathcal{F})$ 和
$(\mathcal{G}, \nabla_\mathcal{G})$ 之间的态射由
$$
\Hom_\mathcal{A}((\mathcal{F}, \nabla_\mathcal{F}),
(\mathcal{G}, \nabla_\mathcal{G})) =
\bigoplus \Hom_{\mathcal{O}_C}(\mathcal{F},
\mathcal{G} \otimes_{\mathcal{O}_C} \Omega_{C/k}^{\otimes n})
$$
给出。对次数为 $n$ 的元素
$f : \mathcal{F} \to \mathcal{G} \otimes \Omega_{C/k}^{\otimes n}$，置
$$
\text{d}(f) =
\nabla_{\mathcal{G} \otimes \Omega_{C/k}^{\otimes n}} \circ f +
f \circ \nabla_\mathcal{F}
$$
其中采用适当的等同。
\end{enumerate}
我们略去验证：这构成具有性质 (A)、(B)、(C) 的微分分次范畴。因此得到三角
同伦范畴 $K(\mathcal{A})$。

\medskip\noindent
若 $C = \mathbf{P}^1_k$，则 $K(\mathcal{A})$ 是零范畴。然而，若 $C$
是亏格 $> 1$ 的光滑真曲线，则 $K(\mathcal{A})$ 非零。事实上，假设
$\mathcal{N}$ 是次数满足 $0 \leq d < g - 1$ 且具有非零截面
$\sigma$ 的可逆层。置
$(\mathcal{F}, \nabla_\mathcal{F}) = (\mathcal{O}_C, \text{d})$
以及
$(\mathcal{G}, \nabla_\mathcal{G}) = (\mathcal{N}^{\otimes 2}, \nabla_{can})$。
可见
% P12-ERR-0384
$$
\Hom_\mathcal{A}^n((\mathcal{F}, \nabla_\mathcal{F}),
(\mathcal{G}, \nabla_\mathcal{G})) =
\left\{
\begin{matrix}
0 & \text{若} & n < 0 \\
\Gamma(C, \mathcal{N}^{\otimes 2}) & \text{若} & n = 0 \\
\Gamma(C, \mathcal{N}^{\otimes 2} \otimes \Omega_{C/k}) & \text{若} & n = 1
\end{matrix}
\right.
$$
% P12-ERR-0385
第一个 $0$ 是因为，由条件 $d < g - 1$，
$\mathcal{N}^{\otimes 2} \otimes \Omega_{C/k}^{\otimes -1}$ 的次数为负。
现在，截面 $\sigma^{\otimes 2}$ 的导数等于零，因而同态群
$$
\Hom_{K(\mathcal{A})}((\mathcal{F}, \nabla_\mathcal{F}),
(\mathcal{G}, \nabla_\mathcal{G}))
$$
非零。







\section{真代数空间的叠不是代数叠}
\label{examples-section-proper-spaces-not-algebraic}

\noindent
在《Quot》章节 \ref{quot-section-stack-of-spaces} 中，我们引入并研究了群胚叠
$$
p'_{fp, flat, proper} :
\Spacesstack'_{fp, flat, proper}
\longrightarrow
\Sch_{fppf}
$$
它在概形 $S$ 上的截面范畴是 $S$ 上平坦、真、有限表现的代数空间所成的
范畴。我们证明了它满足许多 Artin 公理。本节将
% P12-ERR-0386
说明这个叠为什么不是代数叠：一般而言，形式有效性不成立。

\medskip\noindent
典范例利用如下事实：$g$ 维 Abel 簇的泛形变空间有 $g^2$ 个形式参数，
而任何有效形式形变都可定义在维数 $\leq g(g + 1)/2$ 的完备局部环上。
我们的例子将通过写下 K3 曲面的适当非有效形变来构造。我们只勾勒论证，
不提供全部细节。

\medskip\noindent
令 $k = \mathbf{C}$ 为复数域。设 $X \subset \mathbf{P}^3_k$ 为 $k$ 上的
光滑 $4$ 次曲面。我们有
$\omega_X \cong \Omega^2_{X/k} \cong \mathcal{O}_X$。最后，还有
$\dim_k H^0(X, T_{X/k}) = 0$、
$\dim_k H^1(X, T_{X/k}) = 20$，以及
$\dim_k H^2(X, T_{X/k}) = 0$。
由于 $X$ 在 $k$ 上光滑，$L_{X/k} = \Omega_{X/k}$；又因为
$\Ext^i_{\mathcal{O}_X}(\Omega_{X/k}, \mathcal{O}_X) =
H^i(X, T_{X/k})$，并利用《余切复形》引理
\ref{cotangent-lemma-find-obstruction-ringed-topoi}，可知 $X$ 在
$$
k[[x_1, \ldots, x_{20}]]
$$
上存在泛形变。假设该泛形变是有效的（定义见《Artin 公理》章节
\ref{artin-section-formal-objects}）。那么将得到平坦真态射
$$
f : Y \longrightarrow \Spec(k[[x_1, \ldots, x_{20}]])
$$
其中 $Y$ 是恢复该泛形变的代数空间。下面说明这是不可能的。因为 $Y$ 分离，
可找到仿射开子概形 $V \subset Y$。由于 $Y$ 的特殊纤维 $X$ 光滑，可知
$f$ 光滑。因此，$Y$ 在正则空间上光滑，故为正则的；从而 $V$ 在 $Y$
中的补集 $D$ 是有效 Cartier 除子。于是 $\mathcal{O}_Y(D)$ 是
$\Pic(Y)$ 的非平凡元素（为证明这一点，可说明域上真光滑代数空间中非空仿射
开集的补集总在 Picard 群中给出
% P12-ERR-0387
一个非平凡元素，再把它应用于 $f$ 的泛纤维）。最后，为得到矛盾，我们证明
$\Pic(Y) = 0$。事实上，映射 $\Pic(Y) \to \Pic(X)$ 是单射，因为
$H^1(X, \mathcal{O}_X) = 0$（所以 $\mathcal{O}_X$ 到
% P12-ERR-0388
$Y \times \Spec(k[[x_i]]/\mathfrak m^n)$ 的所有形变都是平凡的），
并且由 Grothendieck 存在定理（它断言：在各加厚上给出相同层的凝聚模彼此
同构）。若 $X$ 足够一般，则 $\Pic(X) = \mathbf{Z}$，由
$\mathcal{O}_X(1)$ 生成。因此只须说明，对 $n > 0$，
$\mathcal{O}_X(n)$ 不形变到一阶邻域\footnote{只要映射
$c_1 : \Pic(X) \to H^1(X, \Omega_{X/k})$ 是单射，该论证就成立；
对任意特征为零的域 $k$ 以及 $\mathbf{P}^3_k$ 中任意光滑 $4$ 次超曲面
$X$，这一点都成立。}。考虑杯积
$$
H^1(X, \Omega_{X/k}) \times H^1(X, T_{X/k})
\longrightarrow H^2(X, \mathcal{O}_X)
$$
由凝聚对偶，这是非退化配对。计算表明，Hodge 上同调中的 Chern 类
$c_1(\mathcal{O}_X(n)) \in H^1(X, \Omega_{X/k})$ 非零。
因此存在一个一阶形变，它与 $c_1(\mathcal{O}_X(n))$ 的杯积非零。
最后可说明，该杯积就是提升的障碍类。

\begin{lemma}
\label{examples-lemma-proper-spaces-not-algebraic}
群胚叠
$$
p'_{fp, flat, proper} :
\Spacesstack'_{fp, flat, proper}
\longrightarrow
\Sch_{fppf}
$$
在概形 $S$ 上的截面范畴是 $S$ 上平坦、真、有限表现的代数空间所成的范畴
（见《Quot》章节 \ref{quot-section-stack-of-spaces}）；它不是代数叠。
\end{lemma}

\begin{proof}
若它是代数叠，则每个形式对象都将是有效的；见《Artin 公理》引理
\ref{artin-lemma-effective}。上述讨论
% P12-ERR-0389
表明，在基变换到 $\Spec(\mathbf{C})$ 后并非如此。结论得证。
\end{proof}






\section{非代数 Hom-叠的一个例}
\label{examples-section-non-algebraic-hom-stack}

\noindent
设 $\mathcal{Y}, \mathcal{Z}$ 为概形 $S$ 上的代数叠。{\it Hom-叠}
$\underline{\Mor}_S(\mathcal{Y}, \mathcal{Z})$ 是 $S$ 上的群胚叠，
其在概形 $T$ 上的截面范畴为
$$
\Mor_T(\mathcal{Y} \times_S T, \mathcal{Z} \times_S T)
$$
这里的对象是 $1$-态射，而态射是 $2$-态射。我们略去证明它确实是
$(\Sch/S)_{fppf}$ 上的群胚叠
% P12-ERR-0390
（在此插入将来的参考文献）。当然，一般而言，Hom-叠并非代数叠。
本节给出一个反例，其中 $\mathcal{Y}$ 可由 $S$ 上的真平坦概形表示，
而 $\mathcal{Z}$ 在 $S$ 上光滑且真。

\medskip\noindent
设 $k$ 为代数闭域，但不是有限域的代数闭包。令
$S = \Spec(k[[t]])$ 且 $S_n = \Spec(k[t]/(t^n)) \subset S$。
设 $f : X \to S$ 为满足下列条件的映射：
\begin{enumerate}
\item $f$ 射影且平坦，并且它的纤维是几何连通曲线；
\item
% P12-ERR-0391
纤维 $X_0 = X \times_S S_0$ 是结点曲线，其光滑不可约分支的对偶图含有
一个由有理曲线组成的圈；
\item $X$ 是正则概形。
\end{enumerate}
例如，为构造这样的曲面 $X$，可在 $\mathbf{P}^2_{k[[t]]}$ 中取
$$
X\quad :\quad T_0T_1T_2 - t(T_0^3 + T_1^3 + T_2^3) = 0
$$
令 $A_0$ 为 $k$ 上的非零 Abel 簇，例如一条椭圆曲线。令
$A = A_0 \times_{\Spec(k)} S$ 为与 $A_0$ 对应的 $S$ 上常值 Abel
概形\footnote{基底 $S$ 上的 Abel 概形是 $S$ 上平坦、真、有限表现的群概形，
且其所有纤维都是 Abel 簇。}。我们将证明叠
% P12-ERR-0392
$\mathcal{X} = \underline{\Mor}_S(X, [S/A]))$ 不是代数叠。

\medskip\noindent
回忆，$[S/A]$ 一方面是 $A$ 在 $S$ 上平凡作用的商叠，另一方面等于分类
fppf $A$-挠子的叠；见《叠的例》命题
\ref{examples-stacks-proposition-equal-quotient-stacks}。
注意，$[S/A] = [\Spec(k)/A_0] \times_{\Spec(k)} S$。这使我们可以如下描述
概形 $T$ 上的纤维范畴：
\begin{align*}
\mathcal{X}_T
& =
\underline{\Mor}_S(X, [S/A])_T \\
& =
\Mor_T(X \times_S T, [S/A] \times_S T) \\
& =
\Mor_S(X \times_S T, [S/A]) \\
& =
\Mor_{\Spec(k)}(X \times_S T, [\Spec(k)/A_0])
\end{align*}
这里 $T$ 是任意 $S$-概形。换言之，群胚 $\mathcal{X}_T$ 是
$X \times_S T$ 上 fppf $A_0$-挠子所成的群胚。在讨论
$\mathcal{X}$ 为什么不是代数叠之前，需要几个引理。

\begin{lemma}
\label{examples-lemma-torsors-over-two-dimensional-regular}
设 $W$ 为二维正则整 Noether 概形，分式域为 $K$。设
$G \to W$ 为 Abel 概形。则映射
$H^1_{fppf}(W, G) \to H^1_{fppf}(\Spec(K), G)$ 是单射。
\end{lemma}

\begin{proof}[证明概要]
设 $P \to W$ 为在泛点处平凡的 fppf $G$-挠子。于是有 $W$ 上的态射
$\Spec(K) \to P$，并可取其概形论像 $Z \subset P$。由于
$P \to W$ 是真的（因为它是 $W$ 上真群代数空间的挠子），可知
$Z \to W$ 是真双有理态射。由《域上的空间》引理
\ref{spaces-over-fields-lemma-finite-in-codim-1}，态射 $Z \to W$
在除 $W$ 的有限多个闭点以外之处都是有限的。由
% P12-ERR-0393
（在此插入关于通过曲面二次变换的吹起消解态射不定性的将来参考文献），
$Z \to W$ 的几何纤维的不可约分支都是有理曲线。由《空间中的群胚进阶》引理
\ref{spaces-more-groupoids-lemma-no-nonconstant-morphism-from-P1-to-group}，
不存在从有理曲线到群概形或其挠子的非常值态射。因此 $Z \to W$ 是有限的，
从而 $Z$ 是概形，并且由《态射》引理
\ref{morphisms-lemma-finite-birational-over-normal}，$Z \to W$ 是同构。
换言之，挠子 $P$ 是平凡的。
\end{proof}

\begin{remark}
\label{examples-remark-torsors-over-regular}
可以去掉引理 \ref{examples-lemma-torsors-over-two-dimensional-regular} 中
$W$ 的 $2$-维条件。也就是说，若 $W$ 是 Noether、正则、整的，则：
(1) 对每个 $G$-挠子 $P$，$P(W)\to P(K)$ 是双射；并且
(2) $H^1(W,G)\to H^1(K,G)$ 是单射。证明概要如下。
首先，(2) 显然由 (1) 得到。对 (1)，单射性显然；满射性则像当前证明一样，
归结为证明有理截面的概形论闭包是截面。这个问题在 $W$ 上是 fppf 局部的，
所以经过 \'etale 基变换后，可假设 $P=G$（注意 \'etale 覆盖保持正则性）。
然后注意，$G=\underline{\mathrm{Pic}}^0_{G'/W}$，其中 $G'$ 是对偶 Abel
概形。结论来自如下事实：$G'$ 是正则概形，故泛定义的可逆层可以延拓。
事实上，可把“$W$ 正则”弱化为“光滑 $W$-概形局部为因子概形”。
Laurent Moret-Bailly 从 Raynaud 那里学到了这个 Picard 技巧；
Raynaud 的著作中某处应当有相应参考文献。
\end{remark}

\begin{lemma}
\label{examples-lemma-torsors-over-field-torsion}
设 $G$ 为域 $K$ 上的光滑交换群代数空间。则
$H^1_{fppf}(\Spec(K), G)$ 是挠群。
\end{lemma}

\begin{proof}
作为 $G$ 的形式，$\Spec(K)$ 上的每个 $G$-挠子 $P$ 都在 $K$ 上光滑。
因此，$P$ 在某个有限可分扩张 $L/K$ 上有点。设 $[L : K] = n$。
以 $[n](P)$ 表示如下 $G$-挠子：其类是 $P$ 的类在
$H^1_{fppf}(\Spec(K), G)$ 中的 $n$ 倍。存在 $K$ 上代数空间的典范态射
$$
P \times_{\Spec(K)} \ldots \times_{\Spec(K)} P \to [n](P)
$$
由于 $G$ 交换，该态射是对称的，故它通过商
$$
(P \times_{\Spec(K)} \ldots \times_{\Spec(K)} P)/S_n
$$
分解。另一方面，态射 $\Spec(L) \to P$ 定义了态射
$$
(\Spec(L) \times_{\Spec(K)} \ldots \times_{\Spec(K)} \Spec(L))/S_n
\longrightarrow (P \times_{\Spec(K)} \ldots \times_{\Spec(K)} P)/S_n
$$
读者可以验证，左边的概形有一个 $K$-有理点。因此，$[n](P)$ 是平凡挠子。
\end{proof}

\noindent
为证明 $\mathcal{X} = \underline{\Mor}_S(X, [S/A])$ 不是代数叠，
由《Artin 公理》引理 \ref{artin-lemma-effective}，只须证明如下结论。

\begin{lemma}
\label{examples-lemma-not-essentially-surjective}
典范映射 $\mathcal{X}(S) \to \lim \mathcal{X}(S_n)$ 不是本质满的。
\end{lemma}

\begin{proof}[证明概要]
展开定义，只须验证 $H^1(X, A_0) \to \lim H^1(X_n, A_0)$ 不是满射。
由于 $X$ 正则且射影，由引理 \ref{examples-lemma-torsors-over-field-torsion} 和
\ref{examples-lemma-torsors-over-two-dimensional-regular}，$X$ 上的每个
$A_0$-挠子都是挠的。特别地，群 $H^1(X, A_0)$ 是挠群。因此只须说明：
(a) 群 $H^1(X_0, A_0)$ 不是挠群；以及
(b) 对所有 $n$，映射
$H^1(X_{n + 1}, A_0) \to H^1(X_n, A_0)$ 都是满射。

\medskip\noindent
关于 (a)。在结点曲线 $X_0$ 上，通过在 $X_0$ 的每个分支上取平凡
$A_0$-挠子，并以 $A_0$ 的非挠点作为各结点上方的同构来粘合，可以构造
非挠 $A_0$-挠子 $P_0$。更确切地说，令 $x \in X_0$ 为由有理曲线组成的
一个圈中出现的结点。令 $X'_0 \to X_0$ 为 $X_0$ 在
$X_0 \setminus \{x\}$ 中的正规化。令 $x', x'' \in X'_0$ 为映到
% P12-ERR-0394
$x_0$ 的两个点。然后取 $A_0 \times_{\Spec(k)} X'_0$，并利用非挠点
$a_0 \in A_0(k)$ 给出的平移 $A_0 \to A_0$，把
$A_0 \times {x'}$ 与 $A_0 \times \{x''\}$ 等同（这样的非挠点存在，
因为 $k$ 代数闭但不是有限域的代数闭包；这实际上并非容易证明）。
可以说明该粘合是代数空间（事实上，可以说明它是概形），并且它是 $X_0$
上的
% P12-ERR-0395
非挠 $A_0$-挠子。它之所以非挠，是因为若 $[n](P_0)$ 有截面，则该截面
产生态射 $s : X'_0 \to A_0$，使得在 $A_0(k)$ 的群律下有
$[n](a_0) = s(x') - s(x'')$。然而，由于圈的不可约分支都是有理曲线，
% P12-ERR-0396
截面 $s$ 在其上为常值（《空间中的群胚进阶》引理
\ref{spaces-more-groupoids-lemma-no-nonconstant-morphism-from-P1-to-group}）。
故 $s(x') = s(x'')$，得到矛盾。

\medskip\noindent
关于 (b)。形变理论表明，把 $A_0$-挠子 $P_n \to X_n$ 形变为
$A_0$-挠子 $P_{n + 1} \to X_{n + 1}$ 的障碍位于
$H^2(X_0, \omega)$ 中，其中 $\omega$ 是 $X_0$ 上的某个向量丛。
后一个群因为 $X_0$ 是曲线而消失，这就证明了断言。
\end{proof}

\begin{proposition}
\label{examples-proposition-nonalghomstack}
叠 $\mathcal{X} = \underline{\Mor}_S(X, [S/A])$ 不是代数叠。
\end{proposition}

\begin{proof}
见上述讨论。
\end{proof}

\begin{remark}
\label{examples-remark-contradict-aoki}
命题 \ref{examples-proposition-nonalghomstack} 与 \cite[定理 1.1]{AokiHomStacks}
矛盾。问题在于 $\underline{\Mor}_S(X, [S/A])$ 的形式对象不是有效的。
勘误 \cite{AokiHomStacksErr} 也提到了同一问题；这是对
\cite{AokiHomStacks} 的勘误。
遗憾的是，该勘误接着断言：若 $\mathcal{Z}$ 分离，则
$\underline{\Mor}_S(\mathcal{Y}, \mathcal{Z})$ 是代数叠；
这也与命题 \ref{examples-proposition-nonalghomstack} 矛盾，因为 $[S/A]$ 是分离的。
\end{remark}









\section{不满足强形式有效性的代数叠}
\label{examples-section-non-formal-effectiveness}

\noindent
这是 \cite[例 4.12]{Bhatt-Algebraize}。设 $k$ 为代数闭域，并设 $A$ 为
$k$ 上的 Abel 簇。假设 $A(k)$ 不是挠群（若 $k$ 不是有限域的代数闭包，
这总成立）。令 $\mathcal{X} = [\Spec(k)/A]$。我们断言，存在理想
$I \subset k[x, y]$，使得
$$
\mathcal{X}_{\Spec(k[x, y]^\wedge)}
\longrightarrow
\lim \mathcal{X}_{\Spec(k[x, y]/I^n)}
$$
不是本质满的。事实上，令 $I$ 为由 $xy(x + y - 1)$ 生成的理想。
于是 $X_0 = V(I)$ 由三份 $\mathbf{A}^1_k$ 沿三个点粘成三角形而得。
因此，可以在 $X_0$ 的各不可约分支上取平凡挠子，并使用非挠点给出的平移
进行粘合，从而构造 $X_0$ 上的无限阶 $A$-挠子 $P_0$。与引理
\ref{examples-lemma-not-essentially-surjective} 的证明完全相同，可以把 $P_0$
提升为 $X_n = \Spec(k[x, y]/I^n)$ 上的挠子 $P_n$。
由于 $k[x, y]^\wedge$ 是二维正则整环，由引理
\ref{examples-lemma-torsors-over-two-dimensional-regular} 和
\ref{examples-lemma-torsors-over-field-torsion} 可知，
$\Spec(k[x, y]^\wedge)$ 上任意 $A$-挠子 $P$ 都是挠的。
所以这族挠子不在所显示函子的像中。

\begin{lemma}
\label{examples-lemma-non-formal-effectiveness}
设 $k$ 为不是有限域
% P12-ERR-0397
闭包的代数闭域。设 $A$ 为 $k$ 上的 Abel 簇，并令
$\mathcal{X} = [\Spec(k)/A]$。存在具有满转移映射的 $k$-代数逆系
$R_n$，其转移映射的核局部幂零；并存在位于系统
% P12-ERR-0398
$(\Spec(R_n))$ 上方的 $\mathcal{X}$ 的系统 $(\xi_n)$，使得该系统在
《Artin 公理》注 \ref{artin-remark-strong-effectiveness} 的意义下不是
有效的。
\end{lemma}

\begin{proof}
见上述讨论。
\end{proof}








\section{Grothendieck 存在定理的一个反例}
\label{examples-section-Grothendieck-existence}

\noindent
设 $k$ 为域，并令 $A = k[[t]]$。令 $X$ 为如下粘合：把
$U = \Spec(A[x])$ 和 $V = \Spec(A[y])$ 通过等同
$$
U \setminus \{0_U\} \longrightarrow V \setminus \{0_V\}
$$
粘合，其中该等同把 $x$ 送到 $y$，而
% P12-ERR-0399
$0_U \in U$ 和 $O_V \in V$ 是分别对应于极大理想 $(x, t)$ 和
$(y, t)$ 的点。置 $A_n = A/(t^n)$，并置
$X_n = X \times_{\Spec(A)} \Spec(A_n)$。令 $\mathcal{F}_n$ 为
$X_n$ 上的凝聚层；它对应于 $A_n[x]$-模
$A_n[x]/(x) \cong A_n$ 和 $A_n[y]$-模 $0$，二者按显然方式粘合。
令 $\mathcal{I} \subset \mathcal{O}_X$ 为
% P12-ERR-0400
由 $t$ 生成的理想层。则 $(\mathcal{F}_n)$ 是《概形的上同调》章节
\ref{coherent-section-existence-proper-support} 中定义的范畴
$\textit{Coh}_{\text{支撑在 }A\text{ 上为固有}}(X, \mathcal{I})$ 的对象。
另一方面，该对象不在《概形的上同调》方程
(\ref{coherent-equation-completion-functor-proper-over-A}) 中函子的像内。
事实上，若它
% P12-ERR-0401
在那里，则将存在有限 $A[x]$-模 $M$、有限 $A[y]$-模 $N$ 和同构
$M[1/t] \cong N[1/t]$，使得对所有 $n$ 都有
$M/t^nM \cong A_n[x]/(x)$ 且 $N/t^nN = 0$。容易看出这是不可能的。

\begin{lemma}
\label{examples-lemma-counter-Grothendieck-existence}
凝聚层代数化的反例。
\begin{enumerate}
\item 若删去 $X \to \Spec(A)$ 分离的假设，则《概形的上同调》定理
\ref{coherent-theorem-grothendieck-existence} 中陈述的 Grothendieck
存在定理为假。
\item 若删去 $X \to S$ 分离的假设，则《Quot》定理
\ref{quot-theorem-coherent-algebraic-general} 和
\ref{quot-theorem-coherent-algebraic} 中的凝聚层叠
$\Cohstack_{X/B}$ 一般不是代数叠。
\item 若删去 $X \to B$ 分离的假设，则《Quot》命题
\ref{quot-proposition-quot} 中的函子
$\Quotfunctor_{\mathcal{F}/X/B}$ 一般不是代数空间。
\end{enumerate}
\end{lemma}

\begin{proof}
第 (1) 部分已在上面说明。这表明 $\textit{Coh}_{X/A}$ 不满足
《Artin 公理》章节 \ref{artin-section-axioms} 的公理 [4]。
因此，由《Artin 公理》引理 \ref{artin-lemma-effective}，它不可能是代数叠。
这就说明了 (2)。为看出 (3)，注意对所有 $n$ 都有相容的满射
$\mathcal{O}_{X_n} \to \mathcal{F}_n$。因此，
$\Quotfunctor_{\mathcal{O}_X/X/A}$ 不满足公理 [4]，并可像前面一样看出
(3) 成立。
\end{proof}






\section{仿射形式代数空间}
\label{examples-section-affine-formal-algebraic-space}

\noindent
设 $K$ 为域，并设 $(V_i)_{i \in I}$ 为 $K$ 上非零向量空间的有向逆系，
其转移映射为满射且 $\lim V_i = 0$；见章节 \ref{examples-section-zero-limit}。
令 $R_i = K \oplus V_i$ 为 $K$-代数，其中 $V_i$ 是平方为零的理想。
于是 $R_i$ 构成 $K$-代数的逆系，其转移映射为满射，核为幂零理想，
并且 $\lim R_i = K$。仿射形式代数空间
$X = \colim \Spec(R_i)$ 是一个并非 McQuillan 的仿射形式代数空间之例。

\begin{lemma}
\label{examples-lemma-affine-not-mcquillan}
存在一个并非 McQuillan 的仿射形式代数空间。
\end{lemma}

\begin{proof}
见上述讨论。
\end{proof}

\noindent
令 $0 \to W_i \to V_i \to K \to 0$ 为章节 \ref{examples-section-zero-limit}
中的正合列系统。令
$A_i = K[V_i]/(ww'; w, w' \in W_i)$。
于是存在相容的满射系统 $A_i \to K[t]$，其核幂零；而转移映射
$A_i \to A_j$ 也是核幂零的满射。回忆，$V_i$ 在 $K$ 上自由，其基由
$s \in S_i$ 给出。于是，若 $K$ 的特征为零，$A_i$ 的次数 $d$ 部分在
$K$ 上自由，其基由 $s^d$（$s \in S_i$）给出，而每个基元都映到 $t^d$。
故 $A_i$ 的次数 $d$ 部分所成的逆系同构于向量空间 $V_i$ 所成的逆系。
由于 $\lim V_i = 0$，可知 $\lim A_i = K$，至少当 $K$ 的特征为零时如此。
这给出一个仿射形式代数空间的例，其“正则函数”不能分离点。

\begin{lemma}
\label{examples-lemma-affine-formal-functions-do-not-separate-points}
存在仿射形式代数空间 $X$，其正则函数在如下意义下不能分离点：
若按《形式空间》定义
\ref{formal-spaces-definition-affine-formal-algebraic-space}
写成 $X = \colim X_\lambda$，则
$\lim \Gamma(X_\lambda, \mathcal{O}_{X_\lambda})$ 是域，但
$X_{red}$ 有无穷多个点。
\end{lemma}

\begin{proof}
见上述讨论。
\end{proof}

\noindent
设 $K$、$I$ 和 $(V_i)$ 如上。考虑系统
$$
\Phi = (\Lambda, J_i \subset \Lambda, (M_i) \to (V_i))
$$
其中 $\Lambda$ 是增广 $K$-代数；对 $i \in I$，
$J_i \subset \Lambda$ 是平方为零的理想；$(M_i) \to (V_i)$ 是
$K$-向量空间逆系的映射，使得对每个 $i$，$M_i \to V_i$ 都是满射；
$M_i$ 带有 $\Lambda$-模结构；转移映射 $M_i \to M_j$（$i > j$）
为 $\Lambda$-线性的；并且当 $i > j$ 时，
$J_j M_i \subset \Ker(M_i \to M_j)$。
断言：存在上述系统，使得对所有 $i > j$ 都有 $M_j = M_i/J_j M_i$。

\medskip\noindent
若断言成立，就得到仿射形式代数空间的可表态射
$$
\colim_{i \in I} \Spec(\Lambda/J_i \oplus M_i)
\longrightarrow
\text{Spf}(\lim \Lambda/J_i)
$$
其源不是 McQuillan 的，而靶是 McQuillan 的。这里
$\Lambda/J_i \oplus M_i$ 带有通常的 $\Lambda/J_i$-代数结构，
其中 $M_i$ 是平方为零的理想。可表性恰好转化为条件：
当 $i > j$ 时，$M_i/J_jM_i = M_j$。该态射的源不是 McQuillan 的，
因为投影
% P12-ERR-0402
$\lim_{i \in I} M_i \to M_i$ 不是满射。这是因为映射
$\lim V_i \to V_i$ 不是满射，并且有满射 $M_i \to V_i$。
略去一些细节。

\medskip\noindent
断言的证明。首先注意，至少存在一个系统，即
$$
\Phi_0 = (K, J_i = (0), (V_i) \xrightarrow{\text{id}} (V_i))
$$
给定系统 $\Phi$，我们将证明存在系统态射 $\Phi \to \Phi'$（系统态射按
显然方式定义），使得 $\Ker(M_i/J_j M_i \to M_j)$ 在
$M'_i/J'_j M'_i$ 中的像为零。做到这一点以后，可以使用通常技巧：
对 $n \geq 1$ 归纳地置 $\Phi_n = (\Phi_{n - 1})'$，再取
$\Phi = \colim \Phi_n$，得到具有所需性质的系统。略去细节。

\medskip\noindent
给定 $\Phi$，构造 $\Phi'$。考虑三元组
$u = (i, j, \xi)$ 所成的集合 $U$，其中 $i > j$ 且
$\xi \in \Ker(M_i \to M_j)$。以 $s, t : U \to I$ 表示映射
$s(i, j, \xi) = i$ 和 $t(i, j, \xi) = j$。令
$\xi_u \in M_{s(u)}$ 为 $u$ 的第三个分量。取
$$
\Lambda' = \Lambda[x_u; u \in U]/(x_u x_{u'}; u, u' \in U)
$$
其中增广 $\Lambda' \to K$ 由 $\Lambda$ 的增广给出，并把 $x_u$
送到零。取
% P12-ERR-0403
$J'_k = J_k \Lambda' + (x_{u,\ t(u) \geq k})$。置
$$
M'_i = M_i \oplus \bigoplus\nolimits_{s(u) \geq i} K\epsilon_{i, u}
$$
作为 $i > j$ 时的转移映射 $M'_i \to M'_j$，我们使用给定映射
$M_i \to M_j$，并把 $\epsilon_{i, u}$ 送到 $\epsilon_{j, u}$。
映射 $M'_i \to V_i$ 诱导给定映射 $M_i \to V_i$，并把
$\epsilon_{i, u}$ 送到零。最后，令 $\Lambda'$ 如下作用在 $M'_i$ 上：
对 $\lambda \in \Lambda$，在 $M_i$ 上按 $\Lambda$-模结构作用，并通过增广
$\Lambda \to K$ 作用在 $\epsilon_{i, u}$ 上。对所有 $i$，元素 $x_u$
在 $M_i$ 上作用为 $0$。最后定义
$$
x_u \epsilon_{i, u} = \text{元素 }\xi_u\text{ 映至 }M_i
$$
并令所有其他乘积 $x_{u'} \epsilon_{i, u}$ 等于零。所显示公式有意义，
因为 $s(u) \geq i$ 且 $\xi_u \in M_{s(u)}$。主要需要
% P12-ERR-0404
检验的是：当 $i > j$ 时，$J'_j M'_i \subset M'_i$ 在 $M'_j$ 中映为零；
以及
% P12-ERR-0405
$\Ker(M_i \to M_j)$ 在 $M'_i/J_j M'_i$ 中映为零。后一个事实的理由是，
对任意 $\xi \in \Ker(M_i \to M_j)$，都有
$\xi = x_{(i, j, \xi)} \epsilon_{i, (i, j, \xi)} \in J'_j M'_i$。
我们略去细节。

\begin{lemma}
\label{examples-lemma-representable-morphism-affine-formal-not-mcquillan-top}
存在仿射形式代数空间的可表态射 $f : X \to Y$，使得 $Y$ 是
McQuillan 的，而 $X$ 不是 McQuillan 的。
\end{lemma}

\begin{proof}
见上述讨论。
\end{proof}









\section{平坦映射不一定是有限表现平坦映射的滤过余极限}
\label{examples-section-flat-not-colimit-flat-finitely-presented}

\noindent
本节的目标是给出一个平坦环映射的例，它不是平坦有限表现环映射的滤过
余极限。\cite{gabber-nonexcellent} 证明了：若 $A$ 是维数为 $1$、剩余特征
为零的非优局部环，则到其完备化的（平坦）环映射
$A \to A^\wedge$ 不是有限型平坦环映射的滤过余极限。本节例子的源将是优环。
我们鼓励读者提交其他例子；若你有例子，请发送电子邮件到
\href{mailto:stacks.project@gmail.com}{stacks.project@gmail.com}。

\medskip\noindent
为进行构造，固定素数 $p$，并令
$A = \mathbf{F}_p[x_1, \ldots, x_n]$。选取 $A$ 的绝对整闭包 $A^+$；
也就是说，$A^+$ 是 $A$ 在其分式域的某个代数闭包中的整闭包。
\cite[\S 6.7]{HHBigCM} 证明了 $A \to A^+$ 平坦。

\medskip\noindent
我们断言，当 $n \geq 3$ 时，$A$-代数 $A^+$ 不是有限表现平坦
$A$-代数的滤过余极限。

\medskip\noindent
我们勾勒 $n = 3$ 情形的论证，把推广到更高 $n$ 的工作留给读者。
只须对映射 $R \to R^+$ 证明类似陈述，其中 $R$ 是 $A$ 在原点处的严格
Hensel 化，而 $R^+$ 是它的绝对整闭包。注意，$R$ 是 Hensel 正则局部环，
其剩余域 $k$ 是 $\mathbf{F}_p$ 的代数闭包。

\medskip\noindent
选取 $k$ 上的普通 Abel 曲面 $X$ 和 $X$ 上的甚丰富线丛 $L$。
截面环 $\Gamma_*(X, L) = \bigoplus_n H^0(X,L^n)$ 是 $X$ 关于 $L$
的仿射锥的坐标环。当 $L$ 足够正时，它是正规环。令 $S$ 表示
$\Gamma_*(X, L)$ 在
% P12-ERR-0406
锥顶点处的 Hensel 化。则 $S$ 是维数为 $3$ 的 Hensel Noether 正规整环。
取有限型 $k$-代数 $\Gamma_*(X, L)$ 的一个 Noether 正规化并作 Hensel 化，
得到有限单射 $R \to S$。由于 $R^+$ 是 $R$ 的绝对整闭包，还可固定嵌入
$S \to R^+$。因此，$R^+$ 也是 $S$ 的绝对整闭包。为说明 $R^+$ 不是
平坦 $R$-代数的滤过余极限，只须说明：
\begin{enumerate}
\item 若 $S \to P$ 为环映射，且 $P$ 在 $R$ 上平坦、有限型，则存在环映射
$S \to T$，其中 $T$ 在 $R$ 上有限平坦。
\item 对任意有限环映射 $S \to T$，环 $T$ 都不是 $R$-平坦的。
\end{enumerate}
事实上，因为 $S$ 在 $R$ 上有限表现，若可以把
$R^+ = \colim_i P_i$ 写成有限表现平坦 $R$-代数 $P_i$ 的滤过余极限，
则对 $i \gg 0$，$S \to R^+$ 将分解为 $S \to P_i \to R^+$，
这与上述两个断言矛盾。断言 (1) 来自 $R$ 是 Hensel 环这一事实和一个切片
论证；见《态射进阶》引理
\ref{more-morphisms-lemma-qf-fp-flat-neighbourhood-dominates-fppf}。
第 (2) 部分已在 \cite{BhattSmallCMMod} 中证明；为方便读者，我们回忆该论证。

\medskip\noindent
令 $U \subset \Spec(S)$ 为穿孔谱，于是有自然映射
$X \leftarrow U \subset \Spec(S)$。第一个映射给出等同
$H^1(U, \mathcal{O}_U) \simeq H^1(X, \mathcal{O}_X)$。
通过完美化的 Witt 向量并使用 Artin--Schreier 序列\footnote{这里用到：
$S$ 是特征 $p$ 的严格 Hensel 局部环，因而
$S \to S$，$f \mapsto f^p - f$ 是满射。又因为 $S$ 是正规整环，
$\Gamma(U, \mathcal{O}_U) = S$。所以
$H^1_\etale(U, \mathbf{Z}/p)$ 是由 $f \mapsto f^p - f$ 诱导的映射
$H^1(U, \mathcal{O}_U) \to H^1(U, \mathcal{O}_U)$ 的核。}，
得到等同
$H^1_\etale(U, \mathbf{Z}_p) \simeq H^1_\etale(X, \mathbf{Z}_p)$。
特别地，这个群是秩为 $2$ 的有限自由 $\mathbf{Z}_p$-模
（因为 $X$ 是普通的）。为得到矛盾，假设存在上述 (2) 中的
$R$-平坦 $T$。令 $V \subset \Spec(T)$ 表示 $U$ 的逆像，并以
$f : V \to U$ 表示诱导的有限满射。由于 $U$ 正规，$U_\etale$ 上存在迹映射
$f_*\mathbf{Z}_p \to \mathbf{Z}_p$，它与拉回
$\mathbf{Z}_p \to f_*\mathbf{Z}_p$ 的复合是乘以
$d = \deg(f)$。传到上同调，并利用
$H^1_\etale(U, \mathbf{Z}_p)$ 无挠，可知
$H^1_\etale(V, \mathbf{Z}_p)$ 非零。由于不存在 $R^1\lim$ 干扰，有
$H^1_\etale(V, \mathbf{Z}_p) \simeq \lim H^1_\etale(V, \mathbf{Z}/p^n)$，
故群 $H^1(V_\etale,\mathbf{Z}/p)$ 必非零。由于 $T$ 是 $R$-平坦的，
有 $\Gamma(V, \mathcal{O}_V) = T$，这是严格 Hensel 的；Artin--Schreier
序列表明 $H^1(V, \mathcal{O}_V) \neq 0$。这等价于
$H^2_\mathfrak m(T) \neq 0$，其中 $\mathfrak m \subset R$ 是极大理想。
于是得到矛盾，因为 $T$ 作为 $R$-模有限平坦（即有限自由），而
$H^2_\mathfrak m(R) = 0$。这个矛盾证明了 (2)。

\begin{lemma}
\label{examples-lemma-weird-flat-map}
存在交换环 $A$ 和平坦 $A$-代数 $B$，使得该代数不能写成有限表现平坦
$A$-代数的滤过余极限。事实上，可以选取 $A$ 为有限型
$\mathbf{F}_p$-代数，也可以选为剩余域特征为 $0$ 的 $1$-维 Noether
局部环。
\end{lemma}

\begin{proof}
见上述讨论。
\end{proof}

\section{模范畴对挠模取商}
\label{examples-section-serre-quotient-modulo-torsion-modules}

\noindent
挠群范畴是全体 Abel 群范畴的 Serre 子范畴
（《同调》定义 \ref{homology-definition-serre-subcategory}）。更一般地，对任意环
$A$，挠 $A$-模范畴是全体 $A$-模范畴的 Serre 子范畴；见《代数进阶》第
\ref{more-algebra-section-abelian-categories-modules} 节。若 $A$ 是整环，则其商范畴
（《同调》引理 \ref{homology-lemma-serre-subcategory-is-kernel}）等价于分式域上的
向量空间范畴。这可由下面更一般的命题推出。

\begin{proposition}
\label{examples-proposition-localization-and-serre-quotients}
设 $A$ 为环，$S$ 为 $A$ 的乘法子集。以 $\text{Mod}_A$ 表示 $A$-模范畴，
以 $\mathcal{T}$ 表示其如下 Serre 子范畴：其中每个模的任一元素都被 $S$ 的某个
元素零化。则存在典范等价
$\text{Mod}_A/\mathcal{T} \rightarrow \text{Mod}_{S^{-1}A}$。
\end{proposition}

\begin{proof}
由 $M \mapsto M \otimes_A S^{-1}A$ 给出的函子
$\text{Mod}_A \to \text{Mod}_{S^{-1}A}$ 是正合的（由《代数》命题
\ref{algebra-proposition-localization-exact}），并把 $\mathcal{T}$ 中的模映为零。
因此，由《同调》引理 \ref{homology-lemma-serre-subcategory-is-kernel} 给出的泛性质，
该函子下降为函子
$\text{Mod}_A/\mathcal{T} \to \text{Mod}_{S^{-1}A}$。

\medskip\noindent
反过来，任意满足 $M \otimes_A S^{-1}A = 0$ 的 $A$-模 $M$ 都是
$\mathcal{T}$ 的对象，因为
$M \otimes_A S^{-1}A \cong S^{-1} M$
（《代数》引理 \ref{algebra-lemma-tensor-localization}）。因而《同调》引理
\ref{homology-lemma-quotient-by-kernel-exact-functor} 表明函子
$\text{Mod}_A/\mathcal{T} \to \text{Mod}_{S^{-1}A}$ 是忠实的。

\medskip\noindent
此外，这个嵌入还是本质满的：$S^{-1}A$-模 $N$ 的一个原像是 $N_A$，即把
$N$ 看作 $A$-模；这是因为把 $x \otimes a/s$ 映为 $(a/s) \cdot x$ 的典范映射
$N_A \otimes_A S^{-1}A \to N$ 是 $S^{-1}A$-模的同构。
\end{proof}

\begin{proposition}
\label{examples-proposition-quotient-by-torsion-modules}
设 $A$ 为环。以 $Q(A)$ 表示其全商环（如《代数》例
\ref{algebra-example-localize-at-prime}）。以 $\text{Mod}_A$ 表示 $A$-模范畴，
以 $\mathcal{T}$ 表示其挠模 Serre 子范畴，并以 $\text{Mod}_{Q(A)}$ 表示
$Q(A)$-模范畴。则存在典范等价
$\text{Mod}_A/\mathcal{T} \rightarrow \text{Mod}_{Q(A)}$。
\end{proposition}

\begin{proof}
将命题 \ref{examples-proposition-localization-and-serre-quotients} 应用于乘法子集
$S = \{f \in A \mid f \text{ 为非零因子于 }A\}$ 即立即得到结论，
因为一个模是挠模，当且仅当它的每个元素分别被 $S$ 的某个元素零化。
\end{proof}

\begin{proposition}
\label{examples-proposition-quotient-by-finitely-generated-torsion-modules}
设 $A$ 为 Noether 整环，$K$ 为其分式域。以 $\text{Mod}_A^{fg}$ 表示有限生成
$A$-模范畴，以 $\mathcal{T}^{fg}$ 表示其有限生成挠模 Serre 子范畴。则
$\text{Mod}_A^{fg}/\mathcal{T}^{fg}$ 典范等价于有限维 $K$-向量空间范畴。
\end{proposition}

\begin{proof}
命题 \ref{examples-proposition-quotient-by-torsion-modules} 给出的等价沿嵌入
$\text{Mod}_A^{fg}/\mathcal{T}^{fg} \to \text{Mod}_A/\mathcal{T}$ 限制为等价
$\text{Mod}_A^{fg}/\mathcal{T}^{fg} \to \text{Vect}_K^{fd}$。
Noether 假设保证 $\text{Mod}_A^{fg}$ 是 Abel 范畴（见《代数进阶》第
\ref{more-algebra-section-abelian-categories-modules} 节），并保证典范函子
$\text{Mod}_A^{fg}/\mathcal{T}^{fg} \to \text{Mod}_A/\mathcal{T}$ 是满的
（否则，有限生成模的挠子模可能不是 $\mathcal{T}^{fg}$ 的对象）。
\end{proof}

\begin{proposition}
\label{examples-proposition-quotient-abelian-groups-by-torsion-groups}
Abel 群范畴对其挠群 Serre 子范畴的商范畴是 $\mathbf{Q}$-向量空间范畴。
\end{proposition}

\begin{proof}
该断言直接由命题 \ref{examples-proposition-quotient-by-torsion-modules} 得出。
\end{proof}




\section{不同的余极限拓扑}
\label{examples-section-colimit-topology}

\noindent
本例见 \cite[例 1.2，第 553 页]{TSH}。令
$G_n = \mathbf{Q} \times \mathbf{R}^n$，$n \geq 1$；把它视为加法拓扑群，并赋予
通常的（欧氏）拓扑。考虑闭嵌入 $G_n \to G_{n + 1}$，它把
$(x_0, \ldots, x_n)$ 映为 $(x_0, \ldots, x_n, 0)$。我们断言，给
$G = \colim G_n$ 赋予由下式定义的拓扑：
$$
U \subset G\text{ 为开集} \Leftrightarrow G_n \cap U\text{ 为开集 }\forall n
$$
则所得空间不是拓扑群。

\medskip\noindent
为说明这一点，考虑集合
$$
U = \{(x_0, x_1, x_2, \ldots)\text{，满足 }
|x_j| < |\cos(jx_0)| \text{ 对 } j > 0\}
$$
% P12-ERR-0407
利用 $\pi$ 不是有理数，容易从 $jx_0$ 绝不会是 $\pi/2$ 的整数倍推出
$U \cap G_n$ 是开集。由于 $0 \in U$，若上述拓扑使 $G$ 成为拓扑群，则存在
$0$ 的开邻域 $V \subset G$，使得 $V + V \subset U$。于是，对每个
$j \geq 0$，都存在 $\epsilon_j > 0$，使得当 $|x_j| < \epsilon_j$ 时有
$(0, \ldots, 0, x_j, 0, \ldots) \in V$。由于 $V + V \subset U$，当
$|x_0| < \epsilon_0$ 且 $|x_j| < \epsilon_j$ 时，必有
$$
(x_0, 0, \ldots, 0, x_j, 0, \ldots) \in U
$$
然而，取充分大的 $j$ 使 $j \epsilon_0 > \pi/2$，就可以选取
$x_0 \in \mathbf{Q}$，使得 $|\cos(jx_0)|$ 小于 $\epsilon_j$，进而存在
$x_j$ 满足 $|\cos(jx_0)| < |x_j| < \epsilon_j$。这个矛盾证明了该断言。

\begin{lemma}
\label{examples-lemma-colimit-topology}
存在（Abel）拓扑群组成的系统 $G_1 \to G_2 \to G_3 \to \ldots$，使得在
拓扑空间范畴中取的 $\colim G_n$ 不同于在拓扑群范畴中取的 $\colim G_n$。
\end{lemma}

\begin{proof}
见上述讨论。
\end{proof}







\section{普遍下沉但不是 V 覆盖}
\label{examples-section-universally-submersive-not-V}

\noindent
设 $A$ 为赋值环。设 $\mathfrak p \subset A$ 为既非极小素理想也非极大理想的
素理想。（一个便于记忆的情形是 $A$ 恰有三个素理想，而 $\mathfrak p$ 是
``中间''那个。）考虑仿射概形的态射
$$
\Spec(A_\mathfrak p) \amalg \Spec(A/\mathfrak p) \longrightarrow \Spec(A)
$$
我们断言它是普遍下沉的。为证明这一点，设
$\Spec(B) \to \Spec(A)$ 是由环映射 $A \to B$ 给出的仿射概形态射。于是必须证明
$$
\Spec(B_\mathfrak p) \amalg \Spec(B/\mathfrak pB) \to \Spec(B)
$$
是下沉的。首先，它是满射。其次，设 $T \subset \Spec(B)$ 为子集，并且
$T_1 = \Spec(B_\mathfrak p) \cap T$ 与
$T_2 = \Spec(B/\mathfrak p B) \cap T$ 都是闭集。由于 $T_1$ 和 $T_2$
都是 $B$-代数的谱，可知 $T$ 是某个 $B$-代数的谱的像。因此，要证明 $T$
为闭集，只须证明 $T$ 在专化下稳定；见《代数》引理
\ref{algebra-lemma-image-stable-specialization-closed}。为此，设
$p \leadsto q$ 是 $\Spec(B)$ 中点的一个专化，且 $p \in T$。设 $A'$ 为赋值环，
并设 $\Spec(A') \to \Spec(B)$ 为一个态射，使一般点 $\eta$（属于
$\Spec(A')$）映到 $p$、闭点 $s$（属于 $\Spec(A')$）映到 $q$；见《概形》引理
\ref{schemes-lemma-points-specialize}。注意，复合
$\gamma : \Spec(A') \to \Spec(A)$ 的像恰为点 $\xi \in \Spec(A)$ 之集，
这些点满足 $\gamma(\eta) \leadsto \xi \leadsto \gamma(s)$（略去细节）。
若 $\mathfrak p \not \in \Im(\gamma)$，则或者
$p, q \in \Spec(B_\mathfrak p)$，或者
$p, q \in \Spec(B/\mathfrak pB)$。此时，由 $T_1$（相应地，$T_2$）为闭集，
可得 $q \in T_1$（相应地，$q \in T_2$），因而 $q \in T$。
最后，设 $\mathfrak p \in \Im(\gamma)$，比如
$\mathfrak p = \gamma(r)$。则有专化 $p \leadsto r$ 和 $r \leadsto q$。
在这种情形下，$p, r \in \Spec(B_\mathfrak p)$，并且
$r, q \in \Spec(B/\mathfrak pB)$。于是，
% P12-ERR-0408
先得 $r \in T_1 \subset T$；再由 $r$ 映到 $\mathfrak p$ 得
$r \in T_2$；最后得到所需的 $q \in T_2 \subset T$。

\medskip\noindent
另一方面，我们断言单元素族
$$
\{\Spec(A_\mathfrak p) \amalg \Spec(A/\mathfrak p) \longrightarrow \Spec(A)\}
$$
不是 V 覆盖。见《拓扑》定义 \ref{topologies-definition-V-covering}。事实上，
% P12-ERR-0409
若它是 V 覆盖，就会有赋值环扩张 $A \subset B$，使得
$\Spec(B) \to \Spec(A)$ 分解经过
$\Spec(A_\mathfrak p) \amalg \Spec(A/\mathfrak p)$。
% P12-ERR-0410
这将推出 $\Spec(A')$ 不连通，荒谬。

\begin{lemma}
\label{examples-lemma-universally-submersive-not-V}
存在仿射概形的态射 $X \to Y$，它是普遍下沉的，但 $\{X \to Y\}$
不是 V 覆盖。
\end{lemma}

\begin{proof}
见上述讨论。
\end{proof}









\section{整数谱不是拟紧的}
\label{examples-section-canonical}

\noindent
当然，本节标题并不是把整数谱看作拓扑空间而言的，因为任意谱作为拓扑空间都是
拟紧的（《代数》引理 \ref{algebra-lemma-quasi-compact}）。不，这里所指的是概形范畴
上的典范拓扑中的整数谱，以及一个位点中拟紧对象的定义
（《位点》定义 \ref{sites-definition-quasi-compact}）。

\medskip\noindent
设 $U$ 为素数集 $P$ 上的一个非主超滤子。对一个子集 $T \subset P$，以
$T^c = P \setminus T$ 表示其补集。对 $A \in U$，令
$S_A \subset \mathbf{Z}$ 为由 $p \in A$ 生成的乘法子集。置
$$
\mathbf{Z}_A = S_A^{-1}\mathbf{Z}
$$
若把 $P$ 看作 $\Spec(\mathbf{Z})$ 的闭点集，则注意到
$\Spec(\mathbf{Z}_A) = \{(0)\} \cup A^c \subset \Spec(\mathbf{Z})$。
若 $A, B \in U$，则 $A \cap B \in U$ 且 $A \cup B \in U$，并且有
$$
\mathbf{Z}_{A \cap B} =
\mathbf{Z}_A \times_{\mathbf{Z}_{A \cup B}} \mathbf{Z}_B
$$
（环的纤维积）。特别地，对任意整数 $n$ 和元素
$A_1, \ldots, A_n \in U$，态射
$$
\Spec(\mathbf{Z}_{A_1}) \amalg \ldots \amalg \Spec(\mathbf{Z}_{A_n})
\longrightarrow \Spec(\mathbf{Z})
$$
% P12-ERR-0411
都分解经过某个 $p$ 的 $\Spec(\mathbf{Z}[1/p])$
（事实上，任取 $p \in A_1 \cap \ldots \cap A_n$ 即可）。因此，平坦态射族
$\{\Spec(\mathbf{Z}_A) \to \Spec(\mathbf{Z})\}_{A \in U}$ 联合满射，
但其任何有限子族都不是联合满射。

\medskip\noindent
对一个 $\mathbf{Z}$-模 $M$，置
$$
M_A = S_A^{-1}M = M \otimes_{\mathbf{Z}} \mathbf{Z}_A
$$
断言 I：对每个 $\mathbf{Z}$-模 $M$，都有
$$
M =
\text{等化子}\left(
\xymatrix{
\prod\nolimits_{A \in U} M_A \ar@<1ex>[r] \ar@<-1ex>[r] &
\prod\nolimits_{A, B \in U} M_{A \cup B}
}
\right)
$$
首先，设 $M$ 无挠。则对所有 $A \in U$，都有 $M_A \subset M_P$。
于是只须证明
$$
M = \bigcap\nolimits_{A \in U} M_A\text{，交取于 }M_P = M \otimes \mathbf{Q}
$$
事实上，因为 $U$ 非主，对任意素数 $p$ 都有 $\{p\}^c \in U$。此外，
$M_{\{p\}^c} = M_{(p)}$ 等于在素理想 $(p)$ 处的局部化。
% P12-ERR-0412
所以，上式是显然的，因为已经有 $M_{(2)} \cap M_{(3)} = M$。
其次，设 $M$ 为挠模。则有
$$
M = \bigoplus\nolimits_{p \in P} M[p^\infty]
$$
相应地有
$$
M_A = \bigoplus\nolimits_{p \not \in A} M[p^\infty]
$$
因为我们正对 $A$ 中的素数作局部化。设
$(x_A) \in \prod M_A$ 属于该等化子。记
$x_p = x_{\{p\}^c} \in M[p^{\infty}]$。等化子性质给出
$$
x_A = (x_p)_{p \not \in A}
$$
特别地，它说明除有限多个 $p \not \in A$ 外，$x_p$ 都为零。为完成挠模情形的
证明，只须说明除有限多个素数 $p$ 外，$x_p$ 都为零。若非如此，把
$\{p \in P \mid x_p \not = 0\} = T \amalg T'$ 写成两个无限集合的不交并。
由于 $U$ 是超滤子，或者 $T \not \in U$，或者 $T' \not \in U$
（事实上，若 $T, T'$ 都属于 $U$，则 $U$ 包含
$T \cap T' = \emptyset$，这是不允许的）。不妨设 $T \not \in U$。于是
$T = A^c$，这与上面提到的有限性矛盾。
最后，设 $M$ 为一般的模。考察短正合列
$$
0 \to M_{tors} \to M \to M/M_{tors} \to 0
$$
以及下面的大交换图
$$
\xymatrix{
M_{tors} \ar[r] \ar[d] &
\prod\nolimits_{A \in U} M_{tors, A} \ar@<1ex>[r] \ar@<-1ex>[r] \ar[d] &
\prod\nolimits_{A, B \in U} M_{tors, A \cup B} \ar[d] \\
M \ar[r] \ar[d] &
\prod\nolimits_{A \in U} M_A \ar@<1ex>[r] \ar@<-1ex>[r] \ar[d] &
\prod\nolimits_{A, B \in U} M_{A \cup B} \ar[d] \\
M/M_{tors} \ar[r] &
\prod\nolimits_{A \in U} (M/M_{tors})_A \ar@<1ex>[r] \ar@<-1ex>[r] &
\prod\nolimits_{A, B \in U} (M/M_{tors})_{A \cup B} \\
}
$$
% P12-ERR-0413
利用各列的正合性，以及对挠模 $M_{tors}$ 和无挠模 $M/M_{tors}$ 已得的结论
作图追逐，就为 $M$ 证明了断言 I。这给出了下述引理中现象的一个例子。

\begin{lemma}
\label{examples-lemma-non-fpqc-descent}
存在环 $A$ 和平坦环映射的无限族 $\{A \to A_i\}_{i \in I}$，使得对每个
$A$-模 $M$，都有
$$
M =
\text{等化子}\left(
\xymatrix{
\prod\nolimits_{i \in I} M \otimes_A A_i \ar@<1ex>[r] \ar@<-1ex>[r] &
\prod\nolimits_{i, j \in I} M \otimes_A A_i \otimes_A A_j
}
\right)
$$
但不存在具有同一性质的有限子族。
\end{lemma}

\begin{proof}
见上述讨论。
\end{proof}

\noindent
继续使用素数集 $P$ 上的非主超滤子 $U$。设 $R$ 为环。对 $A \in U$，记
$R_A = S_A^{-1}R = R \otimes \mathbf{Z}_A$。
断言 II：给定闭子集 $T_A \subset \Spec(R_A)$（$A \in U$），并满足
$$
(\Spec(R_{A \cup B}) \to \Spec(R_A))^{-1}T_A =
(\Spec(R_{A \cup B}) \to \Spec(R_B))^{-1}T_B
$$
对所有 $A, B \in U$ 都成立，则存在闭子集 $T \subset \Spec(R)$，使得对所有
$A \in U$，都有 $T_A = (\Spec(R_A) \to \Spec(R))^{-1}(T)$。
对 $A \in U$，设 $I_A \subset R_A$ 为定义 $T_A$ 的根理想。则粘合条件蕴含
$S_{A \cup B}^{-1}I_A = S_{A \cup B}^{-1}I_B$
在 $R_{A \cup B}$ 中对所有 $A, B \in U$ 成立（因为局部化保持根理想这一性质）。
设 $I' \subset R$ 为映入 $I_P \subset R_P = R \otimes \mathbf{Q}$ 的元素之集。
于是对 $A \in U$ 可见：
\begin{enumerate}
\item $I_A \subset I'_A = S_A^{-1}I'$；并且
\item $M_A = I'_A/I_A$ 是挠模。
\end{enumerate}
当然，对 $A, B \in U$，得到典范等同
$S_{A \cup B}^{-1}M_A = S_{A \cup B}^{-1}M_B$。把挠模 $M_A$ 分解为它们的
$p$-准素分量后，读者容易证明：存在 $p$-幂挠 $R$-模 $M_p$，使得
$$
M_A = \bigoplus\nolimits_{p \not \in A} M_p
$$
并与上述典范等同相容。置 $M = \bigoplus_{p \in P} M_p$，便得到典范同构
$M_A = S_A^{-1}M$，且它们与上述典范等同相容。于是得到 $R$-模的典范映射
$$
I' \longrightarrow M
$$
% P12-ERR-0414
它对所有 $A \in U$ 都恢复映射 $I_A \to M_A$。这是上述所有相容性以及先前所证
断言的结果；该断言说 $M$ 是从 $\prod_{A \in U} M_A$ 到
$\prod_{A, B \in U} M_{A \cup B}$ 的两个映射的等化子。
令 $I = \Ker(I' \to M)$。则 $I$ 是理想，并且 $T = V(I)$ 是一个闭子集；
它对所有 $A \in U$ 都恢复闭子集 $T_A$。这证明了断言 II。

\begin{lemma}
\label{examples-lemma-Z-not-quasi-compact}
概形 $\Spec(\mathbf{Z})$ 在概形范畴的典范拓扑中不是拟紧的。
\end{lemma}

\begin{proof}
采用上述记号，考虑态射族
$$
\mathcal{W} = \{\Spec(\mathbf{Z}_A) \to \Spec(\mathbf{Z})\}_{A \in U}
$$
由《下降》引理 \ref{descent-lemma-universal-effective-epimorphism} 和上面证明的两个
断言，这是一个普遍有效满态射。在任意有纤维积的范畴中，普遍有效满态射给
$\mathcal{C}$ 赋予位点结构（模去一些容易解决的集合论问题），从而定义典范拓扑。
因此，$\mathcal{W}$ 是典范拓扑的一个覆盖。另一方面，我们已在上面看到，
任何有限子族
$$
\{\Spec(\mathbf{Z}_{A_i}) \to \Spec(\mathbf{Z})\}_{i = 1, \ldots, n},\quad
n \in \mathbf{N}, A_1, \ldots, A_n \in U
$$
都分解经过某个 $p$ 的 $\Spec(\mathbf{Z}[1/p])$。因此，这个有限族不可能是
普遍有效满态射；更一般地，任何普遍有效满态射
$\{g_j : T_j \to \Spec(\mathbf{Z})\}$ 都不能加细
$\{\Spec(\mathbf{Z}_{A_i}) \to \Spec(\mathbf{Z})\}_{i = 1, \ldots, n}$。
由《位点》定义 \ref{sites-definition-quasi-compact}，这说明
$\Spec(\mathbf{Z})$ 在典范拓扑中不是拟紧的。要看这里的拟紧概念与通常的
拓扑斯论定义一致，见《位点》引理 \ref{sites-lemma-quasi-compact}。
\end{proof}
